\begin{proof}[Proof of \cref{obj:thm:finite-oracle}]
Fix the parameters satisfying the stated assumptions. Throughout the proof, write
\[
c=\begin{pmatrix}-1\\1\end{pmatrix},
\qquad
v(t)=\begin{pmatrix}t\\t+1\end{pmatrix},
\qquad
F_\theta(\alpha,t)
=
\sum_{S\in\mathcal S}
\alpha_S\frac{(g_S^\top v(t))^2}{h_S(\theta)}
\quad
(\alpha\in\mathbb R^{\mathcal S},\ t\in\mathbb R).
\]
Thus \(c^\top v(t)=1\). The staircase quantities are those in
\cref{obj:lem:staircase-refinement}, and the optimized values are those in
\cref{obj:def:staircase-saddle}.

\begin{enumerate}
\item We first establish the existence facts needed for the minimax argument.
For \(S\in\mathcal S\) and \(x\in\mathcal X\), define
\[
b_S(x)=
\begin{cases}
e^\varepsilon,&x\in S,\\
1,&x\notin S.
\end{cases}
\]
Throughout this proof, relabel the four records indexed \(1,2,3,4\)
in \cref{obj:lem:staircase-refinement} by \(0,1,2,3\), respectively,
preserving their order and reindexing their probability coordinates
accordingly. The record probabilities are therefore
\[
(\pi_{\theta,0},\pi_{\theta,1},\pi_{\theta,2},\pi_{\theta,3})
=
\bigl(
(1-p)(1-\mu_0),\ (1-p)\mu_0,\
p(1-\mu_1),\ p\mu_1
\bigr).
\]
They are positive, and hence
\[
h_S(\theta)=\sum_{j=0}^3\pi_{\theta,j}b_S(x_j)>0
\qquad(S\in\mathcal S).
\]
Let the complementary-pattern weights be
\[
(\alpha_{\mathrm{pair}})_S=
\begin{cases}
(e^\varepsilon+1)^{-1},
&S=\{x_0\}\ \text{or}\ S=\mathcal X\setminus\{x_0\},\\
0,&\text{otherwise}.
\end{cases}
\]
Each record belongs to exactly one of these two patterns, so
\[
\sum_{S\in\mathcal S}
(\alpha_{\mathrm{pair}})_S b_S(x)
=\frac{e^\varepsilon+1}{e^\varepsilon+1}=1
\qquad(x\in\mathcal X).
\]
Thus \(\mathcal A_\varepsilon\) is nonempty. Its defining inequalities
and equalities are closed and convex. Moreover, since \(b_S(x)\ge1\),
\[
0\le\alpha_S\le\sum_{T\in\mathcal S}\alpha_T
\le\sum_{T\in\mathcal S}\alpha_Tb_T(x)=1
\qquad
(\alpha\in\mathcal A_\varepsilon,\ S\in\mathcal S).
\]
Consequently, \(\mathcal A_\varepsilon\) is compact.

For every feasible \(\alpha\), define the quadratic coefficients
\[
a_{\mathrm{quad}}(\alpha)
=
2\sum_{S\in\mathcal S}
\alpha_S
\frac{(g_S^\top v(0))(g_{S,0}+g_{S,1})}{h_S(\theta)},
\qquad
b_{\mathrm{quad}}(\alpha)
=
\sum_{S\in\mathcal S}
\alpha_S
\frac{(g_{S,0}+g_{S,1})^2}{h_S(\theta)}\ge0,
\]
where \(g_S=(g_{S,0},g_{S,1})^\top\). Expanding the squares gives
\[
F_\theta(\alpha,u)
=
F_\theta(\alpha,0)
+u\,a_{\mathrm{quad}}(\alpha)
+u^2b_{\mathrm{quad}}(\alpha)
\qquad(u\in\mathbb R).
\]
If \(b_{\mathrm{quad}}(\alpha)=0\), nonnegativity of
\(F_\theta(\alpha,u)\) for every \(u\) forces
\(a_{\mathrm{quad}}(\alpha)=0\): otherwise,
\[
F_\theta\!\left(
\alpha,-\frac{F_\theta(\alpha,0)+1}{a_{\mathrm{quad}}(\alpha)}
\right)=-1.
\]
If \(b_{\mathrm{quad}}(\alpha)>0\), completing the square gives a
minimizer. In both cases, a minimizing coordinate is
\[
t_\alpha=
\begin{cases}
0,&b_{\mathrm{quad}}(\alpha)=0,\\[2mm]
-\dfrac{a_{\mathrm{quad}}(\alpha)}
{2b_{\mathrm{quad}}(\alpha)},
&b_{\mathrm{quad}}(\alpha)>0,
\end{cases}
\]
and
\[
F_\theta(\alpha,u)-F_\theta(\alpha,t_\alpha)
=
b_{\mathrm{quad}}(\alpha)(u-t_\alpha)^2\ge0.
\]
In particular, the profile
\[
\mathscr P(\alpha)=\min_{t\in\mathbb R}F_\theta(\alpha,t)
\qquad(\alpha\in\mathcal A_\varepsilon)
\]
is finite and nonnegative. For each fixed \(t\), the objective is
continuous in \(\alpha\); its infimum over \(t\) is therefore upper
semicontinuous. Compactness supplies a feasible
\(\alpha_{\max}\) such that
\[
\mathscr P(\alpha)\le\mathscr P(\alpha_{\max})
\qquad(\alpha\in\mathcal A_\varepsilon),
\qquad
J^*(\theta,p,\varepsilon)=\mathscr P(\alpha_{\max}).
\]

Define the upper envelope by
\[
\mathscr E(t)=
\max_{\alpha\in\mathcal A_\varepsilon}F_\theta(\alpha,t)
\qquad(t\in\mathbb R).
\]
The maximum exists by compactness and continuity. Joint continuity
of \(F_\theta\), together with compactness of the weight set, also
makes \(\mathscr E\) continuous.

For the first singleton pattern,
\[
g_{\{x_0\}}
=
\begin{pmatrix}-(e^\varepsilon-1)(1-p)\\0\end{pmatrix}.
\]
Define the positive envelope coefficient
\[
\eta_{\mathrm{env}}
=
\frac{(e^\varepsilon-1)^2(1-p)^2}
{(e^\varepsilon+1)h_{\{x_0\}}(\theta)}>0.
\]
Nonnegativity of every summand yields
\[
\eta_{\mathrm{env}}t^2
\le F_\theta(\alpha_{\mathrm{pair}},t)
\le\mathscr E(t)
\qquad(t\in\mathbb R).
\]
Hence the sublevel set at \(\mathscr E(0)\) satisfies
\[
\{t\in\mathbb R:\mathscr E(t)\le\mathscr E(0)\}
\subseteq
\left[
-\sqrt{\frac{\mathscr E(0)}{\eta_{\mathrm{env}}}},
\sqrt{\frac{\mathscr E(0)}{\eta_{\mathrm{env}}}}
\right].
\]
It is nonempty and closed, and therefore compact. Minimizing on
this set gives a coordinate \(t_{\mathrm{env}}\in\mathbb R\) with
\[
\mathscr E(t_{\mathrm{env}})
\le\mathscr E(t)
\qquad(t\in\mathbb R).
\]
% lean: CausalSmith.Stat.LdpAteEfficiencySurface.informationObjective_exists_minimizer
% lean: CausalSmith.Stat.LdpAteEfficiencySurface.informationObjective_profile_exists_maximizer
% lean: CausalSmith.Stat.LdpAteEfficiencySurface.upperEnvelope_exists_minimizer

\item We now prove the minimax identity. Define
\[
\mathscr N(\alpha,t)=-F_\theta(\alpha,t)
\qquad
(\alpha\in\mathcal A_\varepsilon,\ t\in\mathbb R).
\]
For fixed \(t\), this function is continuous and affine in
\(\alpha\), hence quasiconvex. For fixed feasible \(\alpha\), it is
continuous and concave in \(t\), since
\(b_{\mathrm{quad}}(\alpha)\ge0\), hence quasiconcave. The weight
set is nonempty, compact, and convex, and \(\mathbb R\) is convex.
The extrema established above identify the quantities entering
Sion's minimax theorem \citep{Sion1958}:
\[
\sup_{t\in\mathbb R}\mathscr N(\alpha,t)=-\mathscr P(\alpha),
\qquad
\inf_{\alpha\in\mathcal A_\varepsilon}
\sup_{t\in\mathbb R}\mathscr N(\alpha,t)
=-J^*(\theta,p,\varepsilon),
\]
and
\[
\inf_{\alpha\in\mathcal A_\varepsilon}\mathscr N(\alpha,t)
=-\mathscr E(t),
\qquad
\sup_{t\in\mathbb R}
\inf_{\alpha\in\mathcal A_\varepsilon}\mathscr N(\alpha,t)
=-\inf_{t\in\mathbb R}\mathscr E(t).
\]
Sion's theorem equates the last quantities in these two displays.
Changing signs gives
\[
J^*(\theta,p,\varepsilon)
=
\inf_{t\in\mathbb R}\mathscr E(t)
=
\inf_{t\in\mathbb R}
\max_{\alpha\in\mathcal A_\varepsilon}F_\theta(\alpha,t).
\]
% lean: CausalSmith.Stat.LdpAteEfficiencySurface.finiteOracle_minimax

\item For fixed \(t\), define the vertex envelope
\[
\mathscr E_{\mathrm{vert}}(t)
=
\sup_{\alpha\in\operatorname{vert}(\mathcal A_\varepsilon)}
F_\theta(\alpha,t).
\]
The compact convex set \(\mathcal A_\varepsilon\) has extreme
points, and its continuous objective is bounded above there.
Since its vertices belong to the feasible set,
\[
\mathscr E_{\mathrm{vert}}(t)\le\mathscr E(t).
\]
Conversely, define the sublevel set
\[
\mathcal C_t=
\{\alpha\in\mathbb R^{\mathcal S}:
F_\theta(\alpha,t)\le\mathscr E_{\mathrm{vert}}(t)\}.
\]
It is closed by continuity and convex by affinity of the objective.
It contains all extreme points. The closed-convex-hull
characterization of a compact convex set therefore gives
\[
\mathcal A_\varepsilon
=
\overline{\operatorname{conv}
\bigl(\operatorname{vert}(\mathcal A_\varepsilon)\bigr)}
\subseteq\mathcal C_t.
\]
Evaluating at a maximizing feasible weight proves
\[
\mathscr E(t)\le\mathscr E_{\mathrm{vert}}(t),
\qquad
\mathscr E(t)=\mathscr E_{\mathrm{vert}}(t).
\]
The feasible set is a bounded polyhedron in
\(\mathbb R^{\mathcal S}\), so its vertex set is finite and the
vertex supremum is a maximum. Thus the two envelope value sets
coincide, and
\[
J^*(\theta,p,\varepsilon)
=
\inf_{t\in\mathbb R}
\max_{\alpha\in\operatorname{vert}(\mathcal A_\varepsilon)}
F_\theta(\alpha,t).
\]
% lean: CausalSmith.Stat.LdpAteEfficiencySurface.upperEnvelope_eq_vertexEnvelope

\item Let \(Q\) be any admissible stationary channel on an arbitrary
measurable output space. By \cref{obj:lem:staircase-refinement},
there are feasible weights \(\alpha_Q\) and a probability kernel
\(K_Q\) such that
\[
Q=K_Q\circ Q_{\alpha_Q},
\qquad
I_\theta(\alpha_Q)-I_\theta(Q)\succeq0.
\]
Choose a minimizing coordinate \(t_Q\in\mathbb R\) for the
quadratic objective of \(\alpha_Q\). The profile maximum established
above gives
\[
F_\theta(\alpha_Q,t_Q)
=\mathscr P(\alpha_Q)
\le\mathscr P(\alpha_{\max})
=J^*(\theta,p,\varepsilon).
\]
Expanding the information matrix shows, for every feasible
\(\alpha\) and every \(t\in\mathbb R\), that
\[
v(t)^\top I_\theta(\alpha)v(t)
=
\sum_{S\in\mathcal S}
\alpha_S\frac{(g_S^\top v(t))^2}{h_S(\theta)}
=F_\theta(\alpha,t).
\]
Applying the positive semidefinite difference to \(v(t_Q)\) gives
the intermediate bound
\[
v(t_Q)^\top I_\theta(Q)v(t_Q)
\le
v(t_Q)^\top I_\theta(\alpha_Q)v(t_Q)
=
F_\theta(\alpha_Q,t_Q)
\le J^*(\theta,p,\varepsilon).
\]

We next verify that the quantity on the right is positive. Define
the singleton weights
\[
(\alpha_{\mathrm{sing}})_S=
\begin{cases}
(e^\varepsilon+3)^{-1},&S=\{x_j\}\ \text{for some }j\in\{0,1,2,3\},\\
0,&\text{otherwise}.
\end{cases}
\]
For each record, one singleton contributes \(e^\varepsilon\)
and the other three contribute \(1\), so
\[
\sum_{S\in\mathcal S}
(\alpha_{\mathrm{sing}})_S b_S(x)
=\frac{e^\varepsilon+3}{e^\varepsilon+3}=1.
\]
These weights are feasible. Define the positive coefficients
\[
\eta_{\mathrm{control}}
=
\frac{(e^\varepsilon-1)^2(1-p)^2}
{(e^\varepsilon+3)h_{\{x_0\}}(\theta)},
\qquad
\eta_{\mathrm{treatment}}
=
\frac{(e^\varepsilon-1)^2p^2}
{(e^\varepsilon+3)h_{\{x_3\}}(\theta)}.
\]
The first and fourth singleton summands give, respectively,
\[
0\le\eta_{\mathrm{control}}t^2
\le F_\theta(\alpha_{\mathrm{sing}},t),
\qquad
0\le\eta_{\mathrm{treatment}}(t+1)^2
\le F_\theta(\alpha_{\mathrm{sing}},t).
\]
At a minimizing coordinate \(t_{\mathrm{sing}}\), a zero objective
would force both \(t_{\mathrm{sing}}=0\) and
\(t_{\mathrm{sing}}+1=0\). Therefore
\[
0<
F_\theta(\alpha_{\mathrm{sing}},t_{\mathrm{sing}})
=
\mathscr P(\alpha_{\mathrm{sing}})
\le J^*(\theta,p,\varepsilon).
\]

For completeness, the output information matrix is positive
semidefinite. Define its dominating measure and density quantities
by
\[
\nu_Q=\sum_{j=0}^3Q(\,\cdot\mid x_j),
\qquad
d_{Q,j}=\frac{dQ(\,\cdot\mid x_j)}{d\nu_Q},
\]
\[
f_{\theta,Q}(z)=\sum_{j=0}^3\pi_{\theta,j}d_{Q,j}(z),
\qquad
\dot f_{\theta,Q}(z)=
\sum_{j=0}^3(\nabla_\theta\pi_{\theta,j})d_{Q,j}(z).
\]
The information integrand is assigned zero when its denominator
vanishes. Symmetry follows from its outer-product form, and, for
every \(w\in\mathbb R^2\),
\[
w^\top I_\theta(Q)w
=
\int
\frac{(w^\top\dot f_{\theta,Q}(z))^2}
{f_{\theta,Q}(z)}\,d\nu_Q(z)\ge0.
\]
Here \(f_{\theta,Q}\ge0\); the interior parameters and the finite
record alphabet justify the finite-sum and integral operations.

We apply the contrast-variance bound directly, including its range
case split. If \(c\notin\operatorname{range}(I_\theta(Q))\), the
variance is \(+\infty\), which proves the required inequality.
Otherwise choose \(w_Q\in\mathbb R^2\) with
\(I_\theta(Q)w_Q=c\), and define
\[
\ell_Q=c^\top w_Q=w_Q^\top I_\theta(Q)w_Q.
\]
Positive semidefiniteness gives \(\ell_Q\ge0\). Equality would imply
\(I_\theta(Q)w_Q=0\), contradicting \(c\ne0\), so \(\ell_Q>0\).
Any other solution \(\widetilde w_Q\in\mathbb R^2\) satisfies
\[
c^\top\widetilde w_Q
=
w_Q^\top I_\theta(Q)\widetilde w_Q
=
w_Q^\top c
=\ell_Q.
\]
Thus the infimum in the contrast-variance definition equals
\(\ell_Q\).

Define the vector
\[
y_Q=\ell_Qv(t_Q)-w_Q\in\mathbb R^2.
\]
Using \(v(t_Q)^\top c=1\), expansion yields
\[
0\le y_Q^\top I_\theta(Q)y_Q
=
\ell_Q\left(
\ell_Qv(t_Q)^\top I_\theta(Q)v(t_Q)-1
\right).
\]
Since \(\ell_Q>0\), the intermediate bound above implies
\[
1
\le
\ell_Qv(t_Q)^\top I_\theta(Q)v(t_Q)
\le
\ell_QJ^*(\theta,p,\varepsilon).
\]
Taking the reciprocal of the positive optimized information gives
\[
V^*(\theta,p,\varepsilon)
=\frac{1}{J^*(\theta,p,\varepsilon)}
\le\ell_Q
=c^\top I_\theta(Q)^+c.
\]
% lean: CausalSmith.Stat.LdpAteEfficiencySurface.staircase_refinement_interior
% lean: CausalSmith.Stat.LdpAteEfficiencySurface.Jstar_pos_interior
% lean: CausalSmith.Stat.LdpAteEfficiencySurface.channelFisherInfo_posSemidef
% lean: CausalSmith.Stat.LdpAteEfficiencySurface.reciprocal_le_contrastVariance

\item Choose an optimal feasible weight
\(\alpha_{\mathrm{opt}}\in\mathcal A_\varepsilon\) and a minimizing
coordinate \(t_{\mathrm{opt}}\in\mathbb R\). Optimality of \(\alpha_{\mathrm{opt}}\) and minimality of
\(t_{\mathrm{opt}}\) give
\[
F_\theta(\alpha_{\mathrm{opt}},t_{\mathrm{opt}})
=
\mathscr P(\alpha_{\mathrm{opt}})
=
J^*(\theta,p,\varepsilon).
\]
Define
\[
I_{\mathrm{opt}}=I_\theta(\alpha_{\mathrm{opt}}),
\qquad
j_{\mathrm{opt}}
=
v(t_{\mathrm{opt}})^\top I_{\mathrm{opt}}v(t_{\mathrm{opt}})
=
J^*(\theta,p,\varepsilon)>0.
\]
The matrix is symmetric, and
\[
w^\top I_{\mathrm{opt}}w
=
\sum_{S\in\mathcal S}
(\alpha_{\mathrm{opt}})_S
\frac{(g_S^\top w)^2}{h_S(\theta)}
\ge0
\qquad(w\in\mathbb R^2).
\]
Moreover, the minimizing property becomes
\[
v(t_{\mathrm{opt}})^\top I_{\mathrm{opt}}v(t_{\mathrm{opt}})
\le
v(u)^\top I_{\mathrm{opt}}v(u)
\qquad(u\in\mathbb R).
\]

To compute its contrast variance, define
\[
\mathbf 1_2=\begin{pmatrix}1\\1\end{pmatrix},
\qquad
a_{\mathrm{shift}}=\mathbf 1_2^\top I_{\mathrm{opt}}\mathbf 1_2,
\qquad
b_{\mathrm{shift}}
=2\mathbf 1_2^\top I_{\mathrm{opt}}v(t_{\mathrm{opt}}).
\]
Since \(v(t_{\mathrm{opt}}+\zeta)
=v(t_{\mathrm{opt}})+\zeta\mathbf 1_2\), minimality implies
\[
0\le
\zeta b_{\mathrm{shift}}+\zeta^2a_{\mathrm{shift}}
\qquad(\zeta\in\mathbb R).
\]
This forces \(b_{\mathrm{shift}}=0\). Indeed, if it were nonzero,
define
\[
D_{\mathrm{shift}}=|a_{\mathrm{shift}}|+1>0,
\qquad
\zeta_{\mathrm{shift}}
=-\frac{b_{\mathrm{shift}}}{2D_{\mathrm{shift}}}.
\]
Then \(a_{\mathrm{shift}}\le D_{\mathrm{shift}}\) gives
\[
\zeta_{\mathrm{shift}}b_{\mathrm{shift}}
+\zeta_{\mathrm{shift}}^2a_{\mathrm{shift}}
=
-\frac{b_{\mathrm{shift}}^2}{2D_{\mathrm{shift}}}
+\frac{a_{\mathrm{shift}}b_{\mathrm{shift}}^2}
{4D_{\mathrm{shift}}^2}
\le
-\frac{b_{\mathrm{shift}}^2}{4D_{\mathrm{shift}}}<0,
\]
a contradiction.

Thus the two coordinates of
\(I_{\mathrm{opt}}v(t_{\mathrm{opt}})\) sum to zero. Combining this
with the definition of \(j_{\mathrm{opt}}\) gives
\[
I_{\mathrm{opt}}v(t_{\mathrm{opt}})
=
\begin{pmatrix}-j_{\mathrm{opt}}\\j_{\mathrm{opt}}\end{pmatrix}
=j_{\mathrm{opt}}c.
\]
Consequently, the vector
\[
w_{\mathrm{opt}}
=\frac{v(t_{\mathrm{opt}})}{j_{\mathrm{opt}}}
\]
satisfies
\[
I_{\mathrm{opt}}w_{\mathrm{opt}}=c,
\qquad
c^\top w_{\mathrm{opt}}=\frac1{j_{\mathrm{opt}}}.
\]
For every \(w\in\mathbb R^2\) satisfying \(I_{\mathrm{opt}}w=c\),
symmetry gives
\[
c^\top w
=w_{\mathrm{opt}}^\top I_{\mathrm{opt}}w
=w_{\mathrm{opt}}^\top c
=\frac1{j_{\mathrm{opt}}}.
\]
Thus the defining infimum gives
\[
c^\top I_\theta(\alpha_{\mathrm{opt}})^+c
=
\frac1{j_{\mathrm{opt}}}
=
V^*(\theta,p,\varepsilon).
\]

It remains to identify the corresponding channel information.
Feasibility makes each row of \(Q_{\alpha_{\mathrm{opt}}}\) a
probability distribution. For all \(S\in\mathcal S\) and
\(x,x'\in\mathcal X\),
\[
b_S(x)\le e^\varepsilon b_S(x'),
\qquad
(\alpha_{\mathrm{opt}})_Sb_S(x)
\le
e^\varepsilon(\alpha_{\mathrm{opt}})_Sb_S(x').
\]
Summing over any subset of the finite output alphabet gives the
required privacy inequality.

For \(S\in\mathcal S\), define
\[
r_S=\sum_{j=0}^3b_S(x_j)>0.
\]
The dominating measure and conditional densities of this channel
satisfy
\[
\nu_{Q_{\alpha_{\mathrm{opt}}}}(\{S\})
=(\alpha_{\mathrm{opt}})_Sr_S,
\qquad
d_{Q_{\alpha_{\mathrm{opt}}},j}(S)
=
\begin{cases}
b_S(x_j)/r_S,&(\alpha_{\mathrm{opt}})_S>0,\\
0,&(\alpha_{\mathrm{opt}})_S=0.
\end{cases}
\]
The zero-weight outputs have zero dominating mass. At every
positive-weight output, the output density and its derivative are
\[
f_{\theta,Q_{\alpha_{\mathrm{opt}}}}(S)
=\frac{h_S(\theta)}{r_S},
\qquad
\dot f_{\theta,Q_{\alpha_{\mathrm{opt}}}}(S)
=\frac{g_S}{r_S}.
\]
Hence integration over the finite alphabet gives
\[
\begin{aligned}
I_\theta(Q_{\alpha_{\mathrm{opt}}})
&=
\sum_{\substack{S\in\mathcal S\\(\alpha_{\mathrm{opt}})_S>0}}
(\alpha_{\mathrm{opt}})_Sr_S
\frac{g_Sg_S^\top}{h_S(\theta)r_S}\\
&=
\sum_{S\in\mathcal S}
(\alpha_{\mathrm{opt}})_S
\frac{g_Sg_S^\top}{h_S(\theta)}
=
I_\theta(\alpha_{\mathrm{opt}}).
\end{aligned}
\]
Together with the preceding variance calculation, this proves
\[
c^\top I_\theta(Q_{\alpha_{\mathrm{opt}}})^+c
=
c^\top I_\theta(\alpha_{\mathrm{opt}})^+c
=
V^*(\theta,p,\varepsilon).
\]
The universal lower bound and this admissible attaining channel
then yield
\[
\inf_Qc^\top I_\theta(Q)^+c=V^*(\theta,p,\varepsilon),
\]
with the infimum taken over the channel class in the statement.
% lean: CausalSmith.Stat.LdpAteEfficiencySurface.Jstar_exists_optimal_weight
% lean: CausalSmith.Stat.LdpAteEfficiencySurface.contrastVariance_eq_reciprocal_of_minimizer
% lean: CausalSmith.Stat.LdpAteEfficiencySurface.staircaseChannel_stationaryLDP
% lean: CausalSmith.Stat.LdpAteEfficiencySurface.channelFisherInfo_staircase_eq_informationMatrix

\item Finally, set
\[
t^*=t_{\mathrm{env}}\in\mathbb R.
\]
Since \(t_{\mathrm{env}}\) minimizes \(\mathscr E\), the minimax identity
gives
\[
\max_{\alpha\in\mathcal A_\varepsilon}
F_\theta(\alpha,t^*)
=
\mathscr E(t_{\mathrm{env}})
=
\inf_{t\in\mathbb R}\mathscr E(t)
=
J^*(\theta,p,\varepsilon),
\]
which proves the asserted attainment of the upper-envelope value.
% lean: CausalSmith.Stat.LdpAteEfficiencySurface.upperEnvelope_exists_minimizer
% lean: CausalSmith.Stat.LdpAteEfficiencySurface.finite_oracle
\end{enumerate}
\end{proof}

\begin{proof}[Proof of \cref{obj:thm:sequential-minimax}]
Throughout, nonnegative risks are interpreted in \([0,\infty]\). Set
\[
c=(-1,1)^\top,\qquad
J_0=J^*(\theta_0,p,\varepsilon),\qquad
V_0=V^*(\theta_0,p,\varepsilon),
\qquad
\iota(a)=\max\{a,0\}\in[0,\infty]\quad(a\in\mathbb R).
\]
For any procedure sequence \(\mathcal P\), write
\[
\operatorname{Risk}_{n,H}(\mathcal P)
=
\sup_{h\in\mathcal H_{n,H}(\theta_0)}
\mathbb E_{\theta_{n,h}}^{\mathcal P}
\!\left[(U_{n,h}^{\mathcal P})^2\right],
\qquad
\operatorname{Risk}_{\infty}(\mathcal P)
=
\sup_{H>0}\liminf_{n\to\infty}
\operatorname{Risk}_{n,H}(\mathcal P).
\]

\begin{enumerate}
\item
Fix an arbitrary sequentially private procedure sequence \(\mathcal P\), including its output spaces and estimator. Take the oracle coordinate \(t^*\) supplied by \cref{obj:thm:finite-oracle}, as used in \cref{obj:def:sequential-vt-handle}, and define
\[
v^*=v(t^*)=(t^*,t^*+1)^\top,
\qquad
G(\theta',t)
=
\max_{\alpha\in\mathcal A_\varepsilon}
F_{\theta'}(\alpha,t)
\quad
\bigl(\theta'\in(0,1)^2,\ t\in\mathbb R\bigr).
\]
Then
\[
c^\top v^*=1,
\qquad
G(\theta_0,t^*)=J_0.
\]
In particular, \(v^*\ne0\). For \(R>0\), put
\[
H(R)=R\|v^*\|_2>0,
\qquad
\vartheta_n(a)=\theta_0+\frac{a v^*}{\sqrt n}
\quad(n\ge1,\ a\in\mathbb R).
\]
For every fixed \(R\), the entire path \(\vartheta_n([-R,R])\) is interior for sufficiently large \(n\), because
\[
\sup_{|a|\le R}\|\vartheta_n(a)-\theta_0\|_2
\le \frac{R\|v^*\|_2}{\sqrt n}\longrightarrow0.
\]

Use the prior family in \cref{obj:def:sequential-vt-handle} at half the path radius:
\[
\overline w_R(a)=w_{R/2}(a)
=
\begin{cases}
\displaystyle
\frac{15}{8R}\left(1-\frac{4a^2}{R^2}\right)^2,
& |a|<R/2,\\[4pt]
0,& |a|\ge R/2.
\end{cases}
\]
Its derivative is
\[
\overline w_R'(a)
=
\begin{cases}
\displaystyle
-\frac{30a}{R^3}\left(1-\frac{4a^2}{R^2}\right),
& |a|<R/2,\\[4pt]
0,& |a|\ge R/2.
\end{cases}
\]
Thus \(\overline w_R\) is continuously differentiable, is nonnegative, and vanishes at both endpoints of \([-R,R]\). Direct integration gives
\[
\int_{-R}^{R}\overline w_R(a)\,da=1,
\qquad
\int_{-R}^{R}
\frac{\overline w_R'(a)^2}{\overline w_R(a)}\,da
=
\frac{480}{R^5}\int_{-R/2}^{R/2}a^2\,da
=
\frac{40}{R^2},
\]
where the quotient is zero wherever \(\overline w_R=0\).

For sufficiently large \(n\), define the prior-averaged information envelope
\[
\mathcal J_{n,R}
=
\int_{-R}^{R}
\overline w_R(a)G(\vartheta_n(a),t^*)\,da.
\]
The feasible weight set is compact: its row-normalization equations and nonnegative coordinates imply
\[
0\le\alpha_S\le1
\quad(S\in\mathcal S,\ \alpha\in\mathcal A_\varepsilon).
\]
Moreover, for interior \(\theta'\),
\[
1\le h_S(\theta')\le e^\varepsilon.
\]
Consequently, \(F_{\theta'}(\alpha,t^*)\) is jointly continuous in \((\theta',\alpha)\). Compactness gives continuity of its maximum, with
\[
|G(\theta',t^*)-G(\theta_0,t^*)|
\le
\max_{\alpha\in\mathcal A_\varepsilon}
|F_{\theta'}(\alpha,t^*)-F_{\theta_0}(\alpha,t^*)|.
\]
It follows that
\[
\sup_{|a|\le R}
|G(\vartheta_n(a),t^*)-J_0|
\longrightarrow0,
\qquad
\mathcal J_{n,R}\longrightarrow J_0.
\]
Equivalently, one may integrate the pointwise convergence using the eventual dominating function
\[
\bigl|\overline w_R(a)G(\vartheta_n(a),t^*)\bigr|
\le (|J_0|+1)\overline w_R(a).
\]
% lean: CausalSmith.Stat.LdpAteEfficiencySurface.sequentialVTPrior_facts
% lean: CausalSmith.Stat.LdpAteEfficiencySurface.tendsto_sequentialPriorInformationEnvelope

\item
We next obtain the transcript information bound. The product record law in \cref{obj:def:binary-trial-model}, with the category ordering in \cref{obj:synth_10}, has probabilities
\[
\begin{aligned}
\pi_{\theta',0}&=(1-p)(1-\theta'_0),&
\pi_{\theta',1}&=(1-p)\theta'_0,\\
\pi_{\theta',2}&=p(1-\theta'_1),&
\pi_{\theta',3}&=p\theta'_1
\end{aligned}
\qquad(\theta'\in(0,1)^2).
\]
Their directional derivatives along \(v^*\) are
\[
\dot\pi_j^*=(v^*)^\top\nabla_{\theta'}\pi_{\theta',j}
\quad(j=0,1,2,3),
\qquad
\sum_{j=0}^3\dot\pi_j^*=0.
\]

Fix \(n\ge1\). For \(0\le k\le n\), let
\[
\mathsf H_{n,k}=\prod_{i=1}^{k}\mathcal Z_{n,i}.
\]
For each input-category string
\(\mathbf j=(j_1,\ldots,j_k)\in\{0,1,2,3\}^k\), let
\(\Lambda_{n,k,\mathbf j}\) be the probability measure on \(\mathsf H_{n,k}\) obtained by successively applying the first \(k\) release kernels to that fixed input string. Define
\[
\nu_{n,k}
=
\sum_{\mathbf j\in\{0,1,2,3\}^k}\Lambda_{n,k,\mathbf j},
\qquad
\lambda_{n,k,\mathbf j}
=
\frac{d\Lambda_{n,k,\mathbf j}}{d\nu_{n,k}}.
\]
At \(k=0\), the history space is a singleton and its law is the unit point mass. Factorization of the transcript laws in \cref{obj:def:sequential-channels} yields the prefix density
\[
f_{n,k}^{\theta'}(\mathsf y)
=
\sum_{\mathbf j\in\{0,1,2,3\}^k}
\left(\prod_{i=1}^{k}\pi_{\theta',j_i}\right)
\lambda_{n,k,\mathbf j}(\mathsf y).
\]
Its directional derivative is the finite sum
\[
\dot f_{n,k}^{\theta'}(\mathsf y)
=
\sum_{\mathbf j\in\{0,1,2,3\}^k}
\left[
\sum_{i=1}^{k}
\dot\pi_{j_i}^*
\prod_{\substack{1\le i'\le k\\i'\ne i}}
\pi_{\theta',j_{i'}}
\right]
\lambda_{n,k,\mathbf j}(\mathsf y).
\]
Let \(\mathbb P_{n,k}^{\theta'}\) denote this prefix law, and define its score by
\[
s_{n,k}^{\theta'}(\mathsf y)
=
\begin{cases}
\dot f_{n,k}^{\theta'}(\mathsf y)/f_{n,k}^{\theta'}(\mathsf y),
&f_{n,k}^{\theta'}(\mathsf y)>0,\\
0,&f_{n,k}^{\theta'}(\mathsf y)=0.
\end{cases}
\]
All record probabilities are positive. Hence these scores are square integrable; indeed,
\[
|s_{n,k}^{\theta'}|
\le
k\max_{0\le j\le3}
\frac{|\dot\pi_j^*|}{\pi_{\theta',j}}
\quad
\mathbb P_{n,k}^{\theta'}\text{-almost everywhere}.
\]

For \(k<n\), freeze the prefix law at \(\theta'\) and form the joint channel
\[
\overline Q_{n,k}^{\theta'}(d\mathsf y,dz\mid x)
=
\mathbb P_{n,k}^{\theta'}(d\mathsf y)\,
Q_{k+1}^{(n)}(dz\mid x,\mathsf y).
\]
For a measurable joint event \(A\), write
\[
A_{\mathsf y}=\{z:(\mathsf y,z)\in A\}.
\]
The all-history privacy inequalities give
\[
\overline Q_{n,k}^{\theta'}(A\mid x)
=
\int Q_{k+1}^{(n)}(A_{\mathsf y}\mid x,\mathsf y)
\,\mathbb P_{n,k}^{\theta'}(d\mathsf y)
\le
e^\varepsilon\overline Q_{n,k}^{\theta'}(A\mid x').
\]
Thus this is a stationary private channel on the joint history--release space. Define its output law and directional score by
\[
\rho_{n,k}^{\theta'}
=
\sum_{j=0}^3
\pi_{\theta',j}\overline Q_{n,k}^{\theta'}(\,\cdot\mid x_j),
\qquad
\Delta_{n,k}^{\theta'}
=
\frac{d\left(\sum_{j=0}^3
\dot\pi_j^*\overline Q_{n,k}^{\theta'}(\,\cdot\mid x_j)\right)}
{d\rho_{n,k}^{\theta'}}.
\]
Choose a measurable version of the latter derivative, with value zero on a null set.

Differentiating the finite-mixture expression for the successor prefix law gives, almost everywhere under \(\rho_{n,k}^{\theta'}\),
\[
s_{n,k+1}^{\theta'}(\mathsf y,z)
=
s_{n,k}^{\theta'}(\mathsf y)
+
\Delta_{n,k}^{\theta'}(\mathsf y,z).
\]
Here the first summand differentiates the prefix law and the second differentiates the current-record mixture. For every measurable prefix event \(B\),
\[
\int_{B\times\mathcal Z_{n,k+1}}
\Delta_{n,k}^{\theta'}\,d\rho_{n,k}^{\theta'}
=
\sum_{j=0}^3\dot\pi_j^*\,
\mathbb P_{n,k}^{\theta'}(B)
=
0.
\]
Therefore the current score has conditional mean zero given the prefix, and
\[
\int
s_{n,k}^{\theta'}(\mathsf y)\Delta_{n,k}^{\theta'}(\mathsf y,z)
\,d\rho_{n,k}^{\theta'}(\mathsf y,z)
=
0.
\]

Define the prefix score energy
\[
\mathcal I_{n,k}(\theta')
=
\int (s_{n,k}^{\theta'})^2\,d\mathbb P_{n,k}^{\theta'}.
\]
Expanding the square in the score identity gives
\[
\mathcal I_{n,0}(\theta')=0,
\qquad
\mathcal I_{n,k+1}(\theta')
=
\mathcal I_{n,k}(\theta')
+
\int(\Delta_{n,k}^{\theta'})^2\,d\rho_{n,k}^{\theta'}.
\]
Apply \cref{obj:lem:staircase-refinement} to the joint channel. It supplies feasible weights
\(\beta_{n,k}^{\theta'}\in\mathcal A_\varepsilon\) with
\[
\begin{aligned}
\int(\Delta_{n,k}^{\theta'})^2\,d\rho_{n,k}^{\theta'}
&=(v^*)^\top
I_{\theta'}(\overline Q_{n,k}^{\theta'})v^*\\
&\le F_{\theta'}(\beta_{n,k}^{\theta'},t^*)
\le G(\theta',t^*).
\end{aligned}
\]
Induction on the prefix length now yields
\[
\mathcal I_{n,k}(\theta')\le kG(\theta',t^*)
\quad(0\le k\le n).
\]
For the local scalar path \(a\mapsto\vartheta_n(a)\), define
\[
s_n^{\mathrm{loc}}(a,z)
=
\frac{1}{\sqrt n}s_{n,n}^{\vartheta_n(a)}(z),
\qquad
I_n^{\mathrm{loc}}(a)
=
\int\bigl(s_n^{\mathrm{loc}}(a,z)\bigr)^2
\,\mathbb P_{n,n}^{\vartheta_n(a)}(dz).
\]
For sufficiently large \(n\), uniformly over \(|a|\le R\),
\[
0\le I_n^{\mathrm{loc}}(a)
=
\frac{\mathcal I_{n,n}(\vartheta_n(a))}{n}
\le G(\vartheta_n(a),t^*).
\]
Multiplying by the nonnegative prior density and integrating gives
\[
0\le
\int_{-R}^{R}\overline w_R(a)I_n^{\mathrm{loc}}(a)\,da
\le\mathcal J_{n,R}.
\]
The finite-mixture domination and the displayed score bounds also give the integrability needed for this calculation on arbitrary measurable output spaces.
% lean: CausalSmith.Stat.LdpAteEfficiencySurface.prefixScoreEnergy_succ
% lean: CausalSmith.Stat.LdpAteEfficiencySurface.prefixScoreEnergy_le_mul_upperEnvelope
% lean: CausalSmith.Stat.LdpAteEfficiencySurface.scaledSelected_fisher_integrable_and_control

\item
Fix \(R>0\) and take \(n\) sufficiently large that the preceding path is interior. Define the prior--transcript probability measure
\[
\mathbb B_{n,R}(da,dz)
=
\overline w_R(a)\,da\,
\mathbb P_{n,n}^{\vartheta_n(a)}(dz)
\quad(a\in[-R,R]),
\]
the estimation error and its integrated square by
\[
e_n(a,z)=\widehat\tau_n^{\mathcal P}(z)-\tau(\vartheta_n(a)),
\qquad
\mathcal E_{n,R}=\int e_n(a,z)^2\,\mathbb B_{n,R}(da,dz),
\]
and the joint score by
\[
\mathscr S_{n,R}(a,z)
=
s_n^{\mathrm{loc}}(a,z)
+
\begin{cases}
\overline w_R'(a)/\overline w_R(a),&\overline w_R(a)>0,\\
0,&\overline w_R(a)=0.
\end{cases}
\]

First suppose \(\mathcal E_{n,R}<\infty\). The likelihood score has mean zero at each path parameter. Hence its cross term with the prior score integrates to zero, and
\[
\int\mathscr S_{n,R}^2\,d\mathbb B_{n,R}
=
\int_{-R}^{R}\overline w_R(a)I_n^{\mathrm{loc}}(a)\,da
+
\frac{40}{R^2}.
\]
This quantity is positive and is at most \(\mathcal J_{n,R}+40/R^2\).

The target derivative along the path is
\[
\frac{d}{da}\tau(\vartheta_n(a))
=
\frac{c^\top v^*}{\sqrt n}
=
\frac1{\sqrt n}.
\]
Integration by parts therefore gives
\[
\int e_n\mathscr S_{n,R}\,d\mathbb B_{n,R}
=
\frac1{\sqrt n}.
\]
Indeed, the derivative of the product of the error and the joint density is the error times the joint-density derivative minus \(1/\sqrt n\) times that density. Its boundary integral vanishes because \(\overline w_R(\pm R)=0\). Finite integrated squared error and finite joint score energy justify the integrals by Cauchy--Schwarz. Applying the same inequality to the displayed identity yields
\[
\begin{aligned}
n\mathcal E_{n,R}
&\ge
\frac{1}{
\displaystyle
\int_{-R}^{R}\overline w_R(a)I_n^{\mathrm{loc}}(a)\,da
+40/R^2}\\
&\ge
\frac1{\mathcal J_{n,R}+40/R^2}.
\end{aligned}
\]

For every \(|a|\le R\),
\[
av^*\in\mathcal H_{n,H(R)}(\theta_0),
\qquad
\sqrt n\,e_n(a,z)=U_{n,av^*}^{\mathcal P}(z).
\]
Thus the normalized prior average is bounded by the local worst risk:
\[
\begin{aligned}
\iota(n\mathcal E_{n,R})
&=
\int_{-R}^{R}\overline w_R(a)
\mathbb E_{\vartheta_n(a)}^{\mathcal P}
\!\left[(U_{n,av^*}^{\mathcal P})^2\right]\,da\\
&\le \operatorname{Risk}_{n,H(R)}(\mathcal P).
\end{aligned}
\]
If instead \(\mathcal E_{n,R}=\infty\), the same nonnegative-integral identity makes this prior average infinite. Its domination by the local worst risk then gives
\(\operatorname{Risk}_{n,H(R)}(\mathcal P)=\infty\).
Both cases establish the eventual bound
\[
\iota\!\left(\frac1{\mathcal J_{n,R}+40/R^2}\right)
\le
\operatorname{Risk}_{n,H(R)}(\mathcal P).
\]
% lean: CausalSmith.Stat.LdpAteEfficiencySurface.ofReal_vanTreesEnvelope_le_localWorstRisk_of_fisher_control_allErrors

\item
We record positivity of the optimized information before passing to reciprocal limits. Define
\[
a_{\mathrm{ctl}}
=
\frac{(e^\varepsilon-1)^2(1-p)^2}
{e^\varepsilon(e^\varepsilon+3)},
\qquad
a_{\mathrm{trt}}
=
\frac{(e^\varepsilon-1)^2p^2}
{e^\varepsilon(e^\varepsilon+3)},
\qquad
J_{\mathrm{low}}
=
\frac{a_{\mathrm{ctl}}a_{\mathrm{trt}}}
{a_{\mathrm{ctl}}+a_{\mathrm{trt}}}.
\]
These constants are positive. The singleton staircase design
\[
\alpha^{\mathrm{RR}}_S
=
\begin{cases}
(e^\varepsilon+3)^{-1},&S\text{ is a singleton},\\
0,&\text{otherwise}
\end{cases}
\]
belongs to \(\mathcal A_\varepsilon\): each channel row contains one singleton entry with multiplier \(e^\varepsilon\) and three with multiplier one.

Using the singleton patterns \(\{x_0\}\) and \(\{x_3\}\), the upper bound \(h_S(\theta')\le e^\varepsilon\), and nonnegativity of the remaining terms gives, for every interior \(\theta'\) and every \(t\in\mathbb R\),
\[
\begin{aligned}
F_{\theta'}(\alpha^{\mathrm{RR}},t)
&\ge
a_{\mathrm{ctl}}t^2+a_{\mathrm{trt}}(t+1)^2\\
&=
(a_{\mathrm{ctl}}+a_{\mathrm{trt}})
\left(t+\frac{a_{\mathrm{trt}}}
{a_{\mathrm{ctl}}+a_{\mathrm{trt}}}\right)^2
+J_{\mathrm{low}}
\ge J_{\mathrm{low}}.
\end{aligned}
\]
The optimization in \cref{obj:def:staircase-saddle} consequently implies
\[
J^*(\theta',p,\varepsilon)\ge J_{\mathrm{low}}>0,
\qquad
J_0>0.
\]

For each fixed \(R>0\), the convergence of \(\mathcal J_{n,R}\) and the eventual lower bound therefore imply
\[
\iota\!\left(\frac1{J_0+40/R^2}\right)
\le
\liminf_{n\to\infty}
\operatorname{Risk}_{n,H(R)}(\mathcal P)
\le
\operatorname{Risk}_{\infty}(\mathcal P).
\]
To remove the prior-information penalty, take
\[
R_{\mathrm{rad},\ell}=\ell+1\quad(\ell\in\mathbb N).
\]
Then
\[
\frac{40}{R_{\mathrm{rad},\ell}^2}\longrightarrow0,
\qquad
\frac1{J_0+40/R_{\mathrm{rad},\ell}^2}\longrightarrow\frac1{J_0}=V_0.
\]
Passing to this limit proves
\[
\iota(V_0)\le\operatorname{Risk}_{\infty}(\mathcal P).
\]
Since \(\mathcal P\) and its measurable output spaces were arbitrary, this establishes the universal lower bound.
% lean: CausalSmith.Stat.LdpAteEfficiencySurface.inv_information_le_localAsymptoticRisk_of_priorRadius_family
% lean: CausalSmith.Stat.LdpAteEfficiencySurface.sequential_localAsymptoticRisk_lower_bound

\item
Choose the schedule
\[
m_n=\lfloor\sqrt n\rfloor\quad(n\in\mathbb N).
\]
For every integer \(b\ge0\),
\[
n\ge b^2\quad\Longrightarrow\quad m_n\ge b,
\]
so \(m_n\to\infty\). For \(n\ge1\),
\[
0\le\frac{m_n}{n}
\le\frac{\sqrt n}{n}
=\frac1{\sqrt n}\longrightarrow0.
\]
Also,
\[
1\le m_n<n\qquad(n\ge2).
\]
Thus this schedule satisfies
\cref{obj:ass:pilot-diverges,obj:ass:pilot-sublinear}.
% lean: CausalSmith.Stat.LdpAteEfficiencySurface.natSqrt_pilotRates

\item
Let
\[
\mathcal P^*
=
\text{the procedure in \cref{obj:def:pilot-estimator}
with this schedule and the fixed strong-saddle selector}.
\]
Apply \cref{obj:thm:pilot-attainment} at \(\theta_0\), with \(\gamma=1/2\). It supplies sequential privacy at every sample size, the common local limit
\[
L_{\mathcal P^*}=N(0,V_0),
\qquad
\int_{\mathbb R}u^2\,L_{\mathcal P^*}(du)=V_0<\infty,
\]
and, for every \(H>0\),
\[
\lim_{\substack{M\to\infty\\M\in\mathbb N}}
\limsup_{n\to\infty}
\sup_{h\in\mathcal H_{n,H}(\theta_0)}
\mathbb E_{\theta_{n,h}}^{\mathcal P^*}
\!\left[
(U_{n,h}^{\mathcal P^*})^2
\mathbf1\{|U_{n,h}^{\mathcal P^*}|>M\}
\right]
=0.
\]
Monotonicity in the threshold extends this limit to real \(M\to\infty\). These are precisely the regularity requirements, so
\[
\mathcal P^*\in\mathfrak P^{\mathrm{reg}}_\varepsilon(\theta_0).
\]
The conventions in \cref{obj:def:pilot-estimator} specify the procedure at every sample size. The following risk calculations use \(n\ge2\), when both sample portions are nonempty.
% lean: CausalSmith.Stat.LdpAteEfficiencySurface.pilot_attainment

\item
We establish exact risk attainment through the pilot-averaged conditional variance. First obtain uniform bounds for the selected correction. Define
\[
C_{\mathrm{zero}}
=
\sum_{S\in\mathcal S}(g_S^\top v(0))^2,
\qquad
T_{\mathrm{up}}
=
\sqrt{\frac{C_{\mathrm{zero}}}{a_{\mathrm{ctl}}}},
\]
\[
C_{\mathrm{grad}}
=
\sum_{S\in\mathcal S}(|g_{S,0}|+|g_{S,1}|),
\qquad
C_\phi
=
\frac{C_{\mathrm{grad}}(T_{\mathrm{up}}+1)}
{J_{\mathrm{low}}}.
\]
For an interior selector argument \(\eta\), set
\[
v_\eta=v(t_\eta),
\qquad
\mathcal I_\eta=I_\eta(\alpha_\eta),
\qquad
\phi_\eta(S)=
\frac{g_S^\top v_\eta}{J_\eta h_S(\eta)}
\quad(S\in\mathcal S).
\]
The strong-saddle conditions in \cref{obj:def:pilot-estimator} and the reference-design bound give
\[
J_{\mathrm{low}}
\le
a_{\mathrm{ctl}}t_\eta^2+a_{\mathrm{trt}}(t_\eta+1)^2
\le
F_\eta(\alpha^{\mathrm{RR}},t_\eta)
\le
J_\eta.
\]
On the other hand, minimizing in the profiling coordinate and using
\(\alpha_{\eta,S}\le1\) and \(h_S(\eta)\ge1\) gives
\[
J_\eta
=
F_\eta(\alpha_\eta,t_\eta)
\le F_\eta(\alpha_\eta,0)
\le C_{\mathrm{zero}}.
\]
It follows that
\[
|t_\eta|\le T_{\mathrm{up}},
\qquad
J_\eta\ge J_{\mathrm{low}}.
\]
Finally,
\[
|g_S^\top v_\eta|
\le
(|g_{S,0}|+|g_{S,1}|)(|t_\eta|+1)
\le
C_{\mathrm{grad}}(T_{\mathrm{up}}+1),
\]
and hence
\[
|\phi_\eta(S)|\le C_\phi
\quad(\eta\in(0,1)^2,\ S\in\mathcal S).
\]
% lean: CausalSmith.Stat.LdpAteEfficiencySurface.strongSelector_J_bounds
% lean: CausalSmith.Stat.LdpAteEfficiencySurface.strongSelector_t_bound
% lean: CausalSmith.Stat.LdpAteEfficiencySurface.strongSelector_phiTilde_bound

\item
For interior true and pilot parameters \(\theta'\) and \(\eta\), define
\[
\sigma^2(\theta',\eta)
=
\sum_{S\in\mathcal S}
\alpha_{\eta,S}h_S(\theta')\phi_\eta(S)^2
-
\bigl(\tau(\theta')-\tau(\eta)\bigr)^2.
\]
We derive the mean and second-moment identities entering this expression. Put
\[
\mathbf{1}_2=(1,1)^\top.
\]
Since \(v(t)=v(0)+t\mathbf{1}_2\), minimization of the selected quadratic objective gives
\[
0=
\left.\frac{d}{dt}F_\eta(\alpha_\eta,t)\right|_{t=t_\eta}
=
2\mathbf{1}_2^\top\mathcal I_\eta v_\eta.
\]
Thus \(\mathcal I_\eta v_\eta\) lies in the span of \(c\). Together with
\(c^\top v_\eta=1\) and \(v_\eta^\top\mathcal I_\eta v_\eta=J_\eta\), this yields
\[
\mathcal I_\eta v_\eta=J_\eta c.
\]

Channel normalization and the affine dependence of the pattern masses give
\[
\sum_S\alpha_{\eta,S}h_S(\theta')=1,
\qquad
\sum_S\alpha_{\eta,S}g_S=0,
\qquad
h_S(\theta')=h_S(\eta)+g_S^\top(\theta'-\eta).
\]
Therefore
\[
\begin{aligned}
\sum_S\alpha_{\eta,S}h_S(\theta')\phi_\eta(S)
&=
\frac1{J_\eta}
\left[
\sum_S\alpha_{\eta,S}g_S^\top v_\eta
+
(\theta'-\eta)^\top\mathcal I_\eta v_\eta
\right]\\
&=
c^\top(\theta'-\eta)
=
\tau(\theta')-\tau(\eta).
\end{aligned}
\]
At the diagonal,
\[
\begin{aligned}
\sum_S\alpha_{\eta,S}h_S(\eta)\phi_\eta(S)^2
&=
\frac1{J_\eta^2}
\sum_S
\alpha_{\eta,S}\frac{(g_S^\top v_\eta)^2}{h_S(\eta)}\\
&=\frac1{J_\eta}=V^*(\eta,p,\varepsilon).
\end{aligned}
\]
Consequently,
\[
\sigma^2(\theta',\eta)
=
\sum_S\alpha_{\eta,S}h_S(\theta')
\left[
\phi_\eta(S)-\bigl(\tau(\theta')-\tau(\eta)\bigr)
\right]^2,
\]
and
\[
0\le\sigma^2(\theta',\eta)\le C_\phi^2,
\qquad
\sigma^2(\eta,\eta)=V^*(\eta,p,\varepsilon).
\]
The upper bound follows by bounding the uncentered second moment by \(C_\phi^2\) and subtracting the squared mean.
% lean: CausalSmith.Stat.LdpAteEfficiencySurface.phiTilde_mean_identity
% lean: CausalSmith.Stat.LdpAteEfficiencySurface.phiTilde_secondMoment_identity
% lean: CausalSmith.Stat.LdpAteEfficiencySurface.adaptiveScoreVariance_nonneg
% lean: CausalSmith.Stat.LdpAteEfficiencySurface.adaptiveScoreVariance_diag

\item
For \(n\ge2\), define the main-sample size and pilot-prefix probabilities by
\[
N_n=n-m_n>0,
\]
\[
\omega_{n,\theta'}(\mathbf r)
=
\prod_{i=1}^{m_n}
\left[
\sum_{j=0}^3
\pi_{\theta',j}Q_{\mathrm{RR}}(r_i\mid x_j)
\right],
\qquad
\mathbf r\in\{0,1,2,3\}^{m_n},
\quad \theta'\in(0,1)^2.
\]
They satisfy
\[
\omega_{n,\theta'}(\mathbf r)\ge0,
\qquad
\sum_{\mathbf r\in\{0,1,2,3\}^{m_n}}
\omega_{n,\theta'}(\mathbf r)=1.
\]
Let \(\widetilde\theta(\mathbf r)\) be the clipped pilot estimate defined in \cref{obj:def:pilot-estimator}. For a main-output string
\(\mathbf s\in\mathcal S^{N_n}\), factorization gives
\[
\begin{aligned}
&\Pr_{\theta'}^{\mathcal P^*}
\!\left(H_{m_n}=\mathbf r,\,
(S_{m_n+1},\ldots,S_n)=\mathbf s\right)\\
&\qquad=
\omega_{n,\theta'}(\mathbf r)
\prod_{i=1}^{N_n}
\alpha_{\widetilde\theta(\mathbf r),s_i}
h_{s_i}(\theta').
\end{aligned}
\]
Transcripts outside these pilot and main-output formats have probability zero.

Given the pilot prefix, the main corrections are independent with mean
\(\tau(\theta')-\tau(\widetilde\theta(\mathbf r))\) and variance
\(\sigma^2(\theta',\widetilde\theta(\mathbf r))\). The estimator formula therefore gives
\[
\mathbb E_{\theta'}^{\mathcal P^*}
\!\left[
\widehat\tau^*-\tau(\theta')
\mid H_{m_n}=\mathbf r
\right]=0.
\]
Expanding the square of the centered sum, the off-diagonal terms vanish by independence and zero means, so
\[
\mathbb E_{\theta'}^{\mathcal P^*}
\!\left[
n(\widehat\tau^*-\tau(\theta'))^2
\mid H_{m_n}=\mathbf r
\right]
=
\frac n{N_n}\,
\sigma^2(\theta',\widetilde\theta(\mathbf r)).
\]
Define the prefix-averaged variance
\[
\overline\sigma_n^2(\theta')
=
\sum_{\mathbf r\in\{0,1,2,3\}^{m_n}}
\omega_{n,\theta'}(\mathbf r)
\sigma^2(\theta',\widetilde\theta(\mathbf r)).
\]
Summing the conditional identity over the finite prefix space gives, for every admissible \(h\),
\[
\mathbb E_{\theta_{n,h}}^{\mathcal P^*}
\!\left[(U_{n,h}^{\mathcal P^*})^2\right]
=
\frac n{N_n}\,
\overline\sigma_n^2(\theta_{n,h}),
\qquad
0\le\overline\sigma_n^2(\theta_{n,h})\le C_\phi^2.
\]
% lean: CausalSmith.Stat.LdpAteEfficiencySurface.pilotEstimator_localRisk_eq_scaled_prefixVarianceAverage
% lean: CausalSmith.Stat.LdpAteEfficiencySurface.pilotPrefixVarianceAverage_bounds

\item
We next show continuity of the variance at \((\theta_0,\theta_0)\). The uniform selected-coordinate bound supplies the compact interval
\[
K_{\mathrm{prof}}=[-T_{\mathrm{up}},T_{\mathrm{up}}].
\]
For every interior \(\eta\), the selected saddle lies in
\(\mathcal A_\varepsilon\times K_{\mathrm{prof}}\). Its saddle inequalities imply
\[
J^*(\eta,p,\varepsilon)
=
\max_{\alpha\in\mathcal A_\varepsilon}
\min_{t\in K_{\mathrm{prof}}}F_\eta(\alpha,t).
\]
Indeed, evaluating any feasible design at \(t_\eta\) bounds its restricted minimum above by \(J_\eta\), while the selected design has restricted minimum \(J_\eta\).

Joint continuity on the fixed compact weight--coordinate set gives
\[
\begin{aligned}
|J^*(\eta,p,\varepsilon)-J_0|
&\le
\sup_{\substack{\alpha\in\mathcal A_\varepsilon\\
t\in K_{\mathrm{prof}}}}
|F_\eta(\alpha,t)-F_{\theta_0}(\alpha,t)|
\longrightarrow0
\quad(\eta\to\theta_0).
\end{aligned}
\]
Since \(J_0>0\), reciprocation yields
\[
V^*(\eta,p,\varepsilon)\longrightarrow V_0.
\]

Subtracting the diagonal second-moment identity from the definition of \(\sigma^2\) gives the exact expression
\[
\begin{aligned}
\sigma^2(\theta',\eta)-V^*(\eta,p,\varepsilon)
&=
\sum_S
\alpha_{\eta,S}\bigl(h_S(\theta')-h_S(\eta)\bigr)
\phi_\eta(S)^2\\
&\qquad-
\bigl(\tau(\theta')-\tau(\eta)\bigr)^2.
\end{aligned}
\]
Using \(0\le\alpha_{\eta,S}\le1\) and \(|\phi_\eta(S)|\le C_\phi\), we obtain
\[
\begin{aligned}
|\sigma^2(\theta',\eta)-V^*(\eta,p,\varepsilon)|
&\le
C_\phi^2\sum_S|h_S(\theta')-h_S(\eta)|\\
&\qquad+
\bigl(\tau(\theta')-\tau(\eta)\bigr)^2.
\end{aligned}
\]
The right-hand side tends to zero as
\((\theta',\eta)\to(\theta_0,\theta_0)\). Hence
\[
\sigma^2(\theta',\eta)\longrightarrow V_0
\quad\text{as }(\theta',\eta)\to(\theta_0,\theta_0).
\]
% lean: CausalSmith.Stat.LdpAteEfficiencySurface.continuousAt_Jstar
% lean: CausalSmith.Stat.LdpAteEfficiencySurface.continuousAt_Vstar
% lean: CausalSmith.Stat.LdpAteEfficiencySurface.continuousAt_adaptiveScoreVariance_diag

\item
Fix \(H>0\) and an error tolerance \(\delta>0\). By the preceding continuity, choose \(r_{\mathrm{cont}}>0\) such that
\[
\|\theta'-\theta_0\|_\infty<r_{\mathrm{cont}},
\quad
\|\eta-\theta_0\|_\infty<r_{\mathrm{cont}}
\quad\Longrightarrow\quad
|\sigma^2(\theta',\eta)-V_0|<\delta/2
\]
for interior \(\theta'\) and \(\eta\). Define
\[
d_{\mathrm{int}}
=
\min\{\mu_0,1-\mu_0,\mu_1,1-\mu_1\}>0,
\qquad
r_{\mathrm{pil}}
=
\min\{r_{\mathrm{cont}}/3,d_{\mathrm{int}}/8\}>0,
\]
where \(\theta_0=(\mu_0,\mu_1)^\top\).

Uniformly over the bounded local set,
\[
\|\theta_{n,h}-\theta_0\|_\infty
\le\frac H{\sqrt n}.
\]
Thus, eventually, this distance is less than \(r_{\mathrm{pil}}\), and
\[
\delta_n=(m_n+2)^{-1}<d_{\mathrm{int}}/4.
\]
For these \(n\), every admissible local parameter satisfies
\[
\delta_n+r_{\mathrm{pil}}
\le
\min\{
(\theta_{n,h})_0,1-(\theta_{n,h})_0,
(\theta_{n,h})_1,1-(\theta_{n,h})_1
\}.
\]

For an interior \(\theta'\), define the bad-prefix mass
\[
\begin{aligned}
B_n(\theta';r)
&=
\sum_{\mathbf r\in\{0,1,2,3\}^{m_n}}
\omega_{n,\theta'}(\mathbf r)\,
\mathbf1\{
\|\widetilde\theta(\mathbf r)-\theta'\|_\infty\ge r
\}\\
&=
\Pr_{\theta'}^{\mathcal P^*}
\{\|\widetilde\theta-\theta'\|_\infty\ge r\},
\qquad r>0.
\end{aligned}
\]
The randomized-response probabilities in \cref{obj:synth_10} give
\[
\mathbb E_{\theta'}[\widehat u_1]
=
\frac{1+(e^\varepsilon-1)(1-p)\theta'_0}
{e^\varepsilon+3},
\qquad
\mathbb E_{\theta'}[\widehat u_3]
=
\frac{1+(e^\varepsilon-1)p\theta'_1}
{e^\varepsilon+3},
\]
and independence of the pilot releases gives
\[
\operatorname{Var}_{\theta'}(\widehat u_1)\le\frac1{m_n},
\qquad
\operatorname{Var}_{\theta'}(\widehat u_3)\le\frac1{m_n}.
\]
Define the positive frequency thresholds
\[
b_{\mathrm{ctl}}(r)
=
\frac{r(1-p)(e^\varepsilon-1)}{e^\varepsilon+3},
\qquad
b_{\mathrm{trt}}(r)
=
\frac{rp(e^\varepsilon-1)}{e^\varepsilon+3}.
\]
If both empirical-frequency deviations are smaller than their respective thresholds, the inverted arm estimates are within \(r\) of the true coordinates. Under the displayed interior-margin condition, clipping preserves these strict inequalities. Therefore Chebyshev's inequality and the union bound give
\[
B_n(\theta';r)
\le
\frac1{m_n b_{\mathrm{ctl}}(r)^2}
+
\frac1{m_n b_{\mathrm{trt}}(r)^2}.
\]
For fixed \(r>0\), this upper bound tends to zero as \(m_n\to\infty\).

On a good prefix at \(\theta'=\theta_{n,h}\), the pilot estimate satisfies
\[
\|\widetilde\theta(\mathbf r)-\theta_0\|_\infty
<
2r_{\mathrm{pil}}<r_{\mathrm{cont}},
\]
so its conditional variance differs from \(V_0\) by less than \(\delta/2\). On every prefix, the bound \(0\le\sigma^2\le C_\phi^2\) gives
\[
|\sigma^2(\theta_{n,h},\widetilde\theta(\mathbf r))-V_0|
\le C_\phi^2+|V_0|.
\]
Splitting the finite weighted average into good and bad prefixes consequently yields
\[
\begin{aligned}
|\overline\sigma_n^2(\theta_{n,h})-V_0|
&\le
\frac{\delta}{2}
+
(C_\phi^2+|V_0|)B_n(\theta_{n,h};r_{\mathrm{pil}})\\
&\le
\frac{\delta}{2}
+
(C_\phi^2+|V_0|)
\left[
\frac1{m_n b_{\mathrm{ctl}}(r_{\mathrm{pil}})^2}
+
\frac1{m_n b_{\mathrm{trt}}(r_{\mathrm{pil}})^2}
\right].
\end{aligned}
\]
The last term is eventually less than \(\delta/2\), uniformly over admissible \(h\). Hence
\[
\sup_{h\in\mathcal H_{n,H}(\theta_0)}
|\overline\sigma_n^2(\theta_{n,h})-V_0|
\longrightarrow0.
\]
% lean: CausalSmith.Stat.LdpAteEfficiencySurface.pilotPrefixBadMass_le
% lean: CausalSmith.Stat.LdpAteEfficiencySurface.pilotPrefixVarianceAverage_error_le
% lean: CausalSmith.Stat.LdpAteEfficiencySurface.eventually_uniform_pilotPrefixVarianceAverage

\item
The sublinear schedule gives
\[
\frac n{N_n}
=
\frac1{1-m_n/n}\longrightarrow1.
\]
Using the bound on the averaged variance,
\[
\begin{aligned}
\left|
\frac n{N_n}\overline\sigma_n^2(\theta_{n,h})-V_0
\right|
&\le
\left|\frac n{N_n}-1\right|C_\phi^2
+
|\overline\sigma_n^2(\theta_{n,h})-V_0|.
\end{aligned}
\]
Thus the exact risk identity implies
\[
\sup_{h\in\mathcal H_{n,H}(\theta_0)}
\left|
\mathbb E_{\theta_{n,h}}^{\mathcal P^*}
[(U_{n,h}^{\mathcal P^*})^2]-V_0
\right|
\longrightarrow0.
\]

In particular, for every \(\delta>0\), eventually
\[
\iota(V_0-\delta)
\le
\operatorname{Risk}_{n,H}(\mathcal P^*)
\le
\iota(V_0+\delta).
\]
For the lower inequality, use \(h=0\), which belongs to
\(\mathcal H_{n,H}(\theta_0)\); the upper inequality holds for every admissible perturbation and hence for their supremum. These bounds give convergence in the extended nonnegative reals:
\[
\operatorname{Risk}_{n,H}(\mathcal P^*)\longrightarrow\iota(V_0).
\]
The upper bounds also ensure that the risks are eventually finite. Consequently,
\[
\liminf_{n\to\infty}\operatorname{Risk}_{n,H}(\mathcal P^*)
=
\iota(V_0)
\qquad(H>0).
\]
Taking the supremum over positive radii, a nonempty family whose values all equal \(\iota(V_0)\), gives
\[
\operatorname{Risk}_{\infty}(\mathcal P^*)=\iota(V_0).
\]
This is the required attaining witness for the fixed selector.
% lean: CausalSmith.Stat.LdpAteEfficiencySurface.eventually_uniform_scaledPrefixVarianceAverage
% lean: CausalSmith.Stat.LdpAteEfficiencySurface.eventually_localWorstRisk_sandwich
% lean: CausalSmith.Stat.LdpAteEfficiencySurface.pilot_localWorstRisk_tendsto
% lean: CausalSmith.Stat.LdpAteEfficiencySurface.pilot_localAsymptoticRisk_eq

\item
Finally, the universal lower bound applies to every member of
\(\mathfrak P^{\mathrm{reg}}_\varepsilon(\theta_0)\), so
\[
\iota(V_0)
\le
\inf_{\mathcal P\in\mathfrak P^{\mathrm{reg}}_\varepsilon(\theta_0)}
\operatorname{Risk}_{\infty}(\mathcal P).
\]
The regular procedure \(\mathcal P^*\) constructed above belongs to this class and has risk \(\iota(V_0)\). Therefore
\[
\inf_{\mathcal P\in\mathfrak P^{\mathrm{reg}}_\varepsilon(\theta_0)}
\operatorname{Risk}_{\infty}(\mathcal P)
\le
\operatorname{Risk}_{\infty}(\mathcal P^*)
=
\iota(V_0).
\]
The two inequalities establish the claimed infimum equality in \([0,\infty]\).
% lean: CausalSmith.Stat.LdpAteEfficiencySurface.sequential_minimax
\end{enumerate}
\end{proof}

\begin{proof}[Proof of \cref{obj:thm:pilot-attainment}]
Fix the selector and pilot schedule in the statement, and let \(\mathcal P\) be the procedure in \cref{obj:def:pilot-estimator}. Use the staircase notation of \cref{obj:synth_4} and the contrast notation of \cref{obj:synth_6}; thus
\[
q=1-p,\qquad d=e^\varepsilon-1,\qquad c=(-1,1)^\top.
\]
Write
\[
N_n=\max\{n-m_n,0\}.
\]
Reciprocals of zero have the convention in \cref{obj:def:pilot-estimator,obj:synth_1}. For \(n\ge2\), the schedule assumptions give
\[
N_n=n-m_n>0,\qquad
N_n\longrightarrow\infty,\qquad
\frac{n}{N_n}\longrightarrow1.
\]
We use the local alternatives and scaled errors of \cref{obj:synth_2,obj:synth_3}, with baseline \(\theta\).

\begin{enumerate}
\item \emph{Bounds and centered conditional corrections.}
Define the reference weights, supported on the four singleton patterns, by
\[
\alpha^{\mathrm{RR}}_S=
\begin{cases}
(e^\varepsilon+3)^{-1},& |S|=1,\\
0,& |S|\ne1,
\end{cases}
\qquad S\in\mathcal S.
\]
Each record belongs to exactly one singleton pattern, so
\[
\sum_{S\in\mathcal S}\alpha^{\mathrm{RR}}_S b_S(j)
=\frac{e^\varepsilon+3}{e^\varepsilon+3}=1,
\qquad j\in\{0,1,2,3\}.
\]
Hence \(\alpha^{\mathrm{RR}}\in\mathcal A_\varepsilon\). For every interior parameter \(\vartheta\),
\[
1\le h_S(\vartheta)\le e^\varepsilon,
\qquad
0\le\alpha_S\le1
\quad(\alpha\in\mathcal A_\varepsilon).
\]
The first inequality follows because \(h_S(\vartheta)\) is an average of entries of \(b_S\), all of which lie in \([1,e^\varepsilon]\). The second follows from any channel-row normalization and \(b_S(j)\ge1\).

Define the constants
\[
\underline a_0=\frac{d^2q^2}{e^\varepsilon(e^\varepsilon+3)},
\qquad
\underline a_1=\frac{d^2p^2}{e^\varepsilon(e^\varepsilon+3)},
\qquad
\underline J=\frac{\underline a_0\underline a_1}
{\underline a_0+\underline a_1},
\]
and
\[
\overline F_0=\sum_{S\in\mathcal S}(g_S^\top v(0))^2,
\qquad
\overline t=\sqrt{\frac{\overline F_0}{\underline a_0}},
\qquad
\overline g=\sum_{S\in\mathcal S}(|g_{S,0}|+|g_{S,1}|),
\qquad
\overline\phi=\frac{\overline g(\overline t+1)}{\underline J}.
\]
Here \(\underline a_0,\underline a_1,\underline J>0\). The contributions of the singleton patterns \(\{0\}\) and \(\{3\}\) give, for every \(u\in\mathbb R\),
\[
F_\vartheta(\alpha^{\mathrm{RR}},u)
\ge \underline a_0u^2+\underline a_1(u+1)^2
=(\underline a_0+\underline a_1)
\left(u+\frac{\underline a_1}{\underline a_0+\underline a_1}\right)^2
+\underline J.
\]
Also, termwise use of \(\alpha_S\le1\) and \(h_S(\vartheta)\ge1\) gives
\[
F_\vartheta(\alpha,0)\le\overline F_0.
\]
The two saddle inequalities therefore imply
\[
\begin{aligned}
\underline J
&\le F_\vartheta(\alpha^{\mathrm{RR}},t_\vartheta)
\le F_\vartheta(\alpha_\vartheta,t_\vartheta)
=J_\vartheta\\
&\le F_\vartheta(\alpha_\vartheta,0)
\le\overline F_0,
\end{aligned}
\qquad
\underline a_0t_\vartheta^2\le J_\vartheta\le\overline F_0.
\]
Consequently,
\[
J_\vartheta=J^*(\vartheta,p,\varepsilon)\ge\underline J>0,
\qquad |t_\vartheta|\le\overline t.
\]

For \(\vartheta\in(0,1)^2\) and \(S\in\mathcal S\), define the correction at the selection parameter by
\[
\phi_\vartheta(S)
=\frac{g_S^\top v(t_\vartheta)}
{J_\vartheta h_S(\vartheta)}.
\]
At \(\vartheta=\widetilde\theta\), this is the correction in \cref{obj:synth_9}. Since
\[
|g_S^\top v(u)|
\le (|g_{S,0}|+|g_{S,1}|)(|u|+1)
\le\overline g(|u|+1),
\]
the preceding bounds yield
\[
|\phi_\vartheta(S)|\le\overline\phi
\qquad
(\vartheta\in(0,1)^2,\ S\in\mathcal S).
\]

We next establish the exact centering identities. Define
\[
\mathbf1_2=(1,1)^\top.
\]
Minimization in the profiling coordinate gives
\[
0=\left.\frac{d}{du}F_\vartheta(\alpha_\vartheta,u)
\right|_{u=t_\vartheta}
=2\mathbf1_2^\top I_\vartheta(\alpha_\vartheta)v(t_\vartheta).
\]
Thus \(I_\vartheta(\alpha_\vartheta)v(t_\vartheta)\) is a multiple of \(c\). Because
\[
c^\top v(t_\vartheta)=1,
\qquad
v(t_\vartheta)^\top I_\vartheta(\alpha_\vartheta)v(t_\vartheta)
=J_\vartheta,
\]
that multiple is \(J_\vartheta\), and hence
\[
I_\vartheta(\alpha_\vartheta)v(t_\vartheta)=J_\vartheta c.
\]
Channel normalization and the affine formula for \(h_S\) also give
\[
\sum_S\alpha_{\vartheta,S}g_S=0,
\qquad
h_S(\xi)=h_S(\vartheta)+g_S^\top(\xi-\vartheta).
\]

For \(\xi,\vartheta\in(0,1)^2\), define the main-release probabilities and centered correction by
\[
\omega_{\xi,\vartheta}(S)
=\alpha_{\vartheta,S}h_S(\xi),
\qquad
\zeta_{\xi,\vartheta}(S)
=\phi_\vartheta(S)-c^\top(\xi-\vartheta).
\]
These probabilities are nonnegative and sum to one. Expanding the affine denominator gives
\[
\begin{aligned}
\sum_S\omega_{\xi,\vartheta}(S)\phi_\vartheta(S)
&=\frac1{J_\vartheta}
\left\{
\sum_S\alpha_{\vartheta,S}g_S^\top v(t_\vartheta)
+(\xi-\vartheta)^\top
 I_\vartheta(\alpha_\vartheta)v(t_\vartheta)
\right\}\\
&=c^\top(\xi-\vartheta).
\end{aligned}
\]
At the diagonal,
\[
\sum_S\omega_{\vartheta,\vartheta}(S)\phi_\vartheta(S)^2
=\frac{F_\vartheta(\alpha_\vartheta,t_\vartheta)}{J_\vartheta^2}
=\frac1{J_\vartheta}
=V^*(\vartheta,p,\varepsilon).
\]
In particular,
\[
\sum_S\omega_{\xi,\vartheta}(S)\zeta_{\xi,\vartheta}(S)=0,
\qquad
|\zeta_{\xi,\vartheta}(S)|\le2\overline\phi.
\]
For the last inequality, the exact mean identity and the probability normalization imply
\[
|c^\top(\xi-\vartheta)|
=\left|\sum_S\omega_{\xi,\vartheta}(S)\phi_\vartheta(S)\right|
\le\overline\phi.
\]
% lean: CausalSmith.Stat.LdpAteEfficiencySurface.pilotSubsequenceRowModel

\item \emph{Continuity of the conditional variance.}
For \(\xi,\vartheta\in(0,1)^2\), define
\[
\sigma^2(\xi,\vartheta)
=\sum_S\omega_{\xi,\vartheta}(S)\zeta_{\xi,\vartheta}(S)^2
=\sum_S\alpha_{\vartheta,S}h_S(\xi)\phi_\vartheta(S)^2
-\{c^\top(\xi-\vartheta)\}^2.
\]
The diagonal second-moment identity gives
\[
\sigma^2(\vartheta,\vartheta)=V^*(\vartheta,p,\varepsilon).
\]
Subtracting this identity from the preceding formula yields
\[
\begin{aligned}
\sigma^2(\xi,\vartheta)-V^*(\vartheta,p,\varepsilon)
={}&\sum_S\alpha_{\vartheta,S}
\{h_S(\xi)-h_S(\vartheta)\}\phi_\vartheta(S)^2\\
&-\{c^\top(\xi-\vartheta)\}^2.
\end{aligned}
\]
Therefore,
\[
\left|\sigma^2(\xi,\vartheta)-V^*(\vartheta,p,\varepsilon)\right|
\le
\overline\phi^2\sum_S|h_S(\xi)-h_S(\vartheta)|
+\{c^\top(\xi-\vartheta)\}^2.
\]

To establish continuity of the optimized value, define
\[
K=[-\overline t,\overline t],
\qquad
\mathcal E_\xi(u)=\max_{\alpha\in\mathcal A_\varepsilon}
F_\xi(\alpha,u),
\quad \xi\in(0,1)^2,\ u\in\mathbb R.
\]
The saddle inequalities give
\[
\mathcal E_\xi(t_\xi)=J_\xi,
\qquad
\mathcal E_\xi(u)\ge
F_\xi(\alpha_\xi,u)\ge J_\xi.
\]
Since \(t_\xi\in K\),
\[
J^*(\xi,p,\varepsilon)=\min_{u\in K}\mathcal E_\xi(u).
\]
The feasible set is nonempty, closed, and contained in \([0,1]^{14}\), hence compact. The objective is jointly continuous on interior parameters, feasible weights, and \(K\). The elementary bounds for maxima and minima give
\[
\left|J^*(\xi,p,\varepsilon)-J^*(\theta,p,\varepsilon)\right|
\le
\sup_{\substack{\alpha\in\mathcal A_\varepsilon\\u\in K}}
|F_\xi(\alpha,u)-F_\theta(\alpha,u)|
\longrightarrow0
\quad(\xi\to\theta).
\]
The convergence follows from joint continuity on a compact parameter neighborhood times \(\mathcal A_\varepsilon\times K\). Positivity of \(J^*\) then gives
\[
V^*(\xi,p,\varepsilon)\longrightarrow V^*(\theta,p,\varepsilon).
\]
Combining this with the displayed variance-error bound proves
\[
\sigma^2(\xi,\vartheta)\longrightarrow V^*(\theta,p,\varepsilon)
\qquad
\text{as }(\xi,\vartheta)\to(\theta,\theta)
\text{ through interior pairs}.
\]
% lean: CausalSmith.Stat.LdpAteEfficiencySurface.continuousAt_adaptiveScoreVariance_diag

\item \emph{Gaussian convergence of deterministic conditional rows.}
Fix \(h\in\mathbb R^2\). To define valid row probabilities at every index, set
\[
\xi_{n,h}=
\begin{cases}
\theta_{n,h},&\theta_{n,h}\in(0,1)^2,\\
\theta,&\theta_{n,h}\notin(0,1)^2.
\end{cases}
\]
By \cref{obj:synth_2}, \(\theta_{n,h}\to\theta\); thus it is eventually interior, and
\[
\xi_{n,h}\to\theta,
\qquad
\xi_{n,h}=\theta_{n,h}\quad\text{eventually}.
\]

For \(n\ge2\) and \(\vartheta\in(0,1)^2\), let
\(\mathbb P^{\mathrm{row}}_{n,h,\vartheta}\) be the product law of \(N_n\) independent pattern releases with
\[
\mathbb P^{\mathrm{row}}_{n,h,\vartheta}(S_j=S)
=\omega_{\xi_{n,h},\vartheta}(S),
\qquad 1\le j\le N_n,\quad S\in\mathcal S.
\]
Define its total-sample-scaled distribution function by
\[
\mathcal C_{n,h}(x;\vartheta)
=
\mathbb P^{\mathrm{row}}_{n,h,\vartheta}
\left\{
\frac{\sqrt n}{N_n}
\sum_{j=1}^{N_n}\zeta_{\xi_{n,h},\vartheta}(S_j)\le x
\right\},
\qquad x\in\mathbb R.
\]
Also define the target distribution function
\[
\mathcal G_\theta(x)
=\Phi\left(\frac{x}{\sqrt{V^*(\theta,p,\varepsilon)}}\right).
\]

Take any strictly increasing sequence of integers \(n_k\ge2\) and any interior sequence \(\vartheta_k\to\theta\). Under the corresponding row law, put
\[
Z_{k,j}=\zeta_{\xi_{n_k,h},\vartheta_k}(S_j),
\qquad 1\le j\le N_{n_k},
\]
and
\[
\mathcal D_k(y)
=
\mathbb P^{\mathrm{row}}_{n_k,h,\vartheta_k}
\left\{
N_{n_k}^{-1/2}\sum_{j=1}^{N_{n_k}}Z_{k,j}\le y
\right\},
\qquad y\in\mathbb R.
\]
The rows are independent, centered, and bounded by \(2\overline\phi\), and their single-release variances satisfy
\[
\mathbb E^{\mathrm{row}}[Z_{k,1}^2]
=\sigma^2(\xi_{n_k,h},\vartheta_k)
\longrightarrow V^*(\theta,p,\varepsilon)>0.
\]
The increments normalized by
\(\sqrt{N_{n_k}V^*(\theta,p,\varepsilon)}\) therefore satisfy
\[
\sum_{j=1}^{N_{n_k}}
\mathbb E^{\mathrm{row}}
\left[
\left(\frac{Z_{k,j}}
{\sqrt{N_{n_k}V^*(\theta,p,\varepsilon)}}\right)^2
\right]
=
\frac{\sigma^2(\xi_{n_k,h},\vartheta_k)}
{V^*(\theta,p,\varepsilon)}
\longrightarrow1,
\]
and
\[
\sum_{j=1}^{N_{n_k}}
\mathbb E^{\mathrm{row}}
\left[
\left|\frac{Z_{k,j}}
{\sqrt{N_{n_k}V^*(\theta,p,\varepsilon)}}\right|^4
\right]
\le
\frac{(2\overline\phi)^4}
{N_{n_k}\{V^*(\theta,p,\varepsilon)\}^2}
\longrightarrow0.
\]
The triangular-array Lyapunov central limit theorem consequently gives
\[
\mathcal D_k(y)\longrightarrow\mathcal G_\theta(y)
\qquad(y\in\mathbb R).
\]

Define
\[
\lambda_k=\sqrt{\frac{n_k}{N_{n_k}}},
\qquad
\chi_k=\frac{x}{\lambda_k}.
\]
Then \(\lambda_k\to1\), \(\chi_k\to x\), and
\[
\mathcal C_{n_k,h}(x;\vartheta_k)=\mathcal D_k(\chi_k).
\]
For every \(a>0\), eventually
\[
\mathcal D_k(x-a)\le\mathcal D_k(\chi_k)\le\mathcal D_k(x+a).
\]
Pointwise convergence and continuity of \(\mathcal G_\theta\), followed by \(a\downarrow0\), prove
\[
\mathcal C_{n_k,h}(x;\vartheta_k)\longrightarrow\mathcal G_\theta(x).
\]
% lean: CausalSmith.Stat.LdpAteEfficiencySurface.adaptivePilotRowScaledCDF_subsequence_tendsto_gaussian

\item \emph{Averaging over the private pilot.}
For \(n\ge2\) and a pilot release vector
\(\mathbf r=(r_1,\ldots,r_{m_n})\in\{0,1,2,3\}^{m_n}\), define
\[
\vartheta_n(\mathbf r)
=\widetilde\theta
\quad\text{computed from }R_i=r_i
\text{ by }\cref{obj:def:pilot-estimator}.
\]
Clipping ensures \(\vartheta_n(\mathbf r)\in(0,1)^2\) for every \(\mathbf r\). Define the pilot weights
\[
w_{n,h}(\mathbf r)
=
\prod_{i=1}^{m_n}
\frac{1+d\,\pi_{\xi_{n,h},r_i}}{e^\varepsilon+3}.
\]
By \cref{obj:synth_10}, these weights are nonnegative and sum to one. Factorization of the successive release kernels gives, whenever \(\theta_{n,h}\) is interior,
\[
\Pr_{\theta_{n,h}}\{U_{n,h}^{\mathcal P}\le x\}
=
\sum_{\mathbf r}w_{n,h}(\mathbf r)
\mathcal C_{n,h}(x;\vartheta_n(\mathbf r)).
\]
Indeed, the exact mean identity gives, for each fixed pilot,
\[
U_{n,h}^{\mathcal P}
=
\frac{\sqrt n}{N_n}
\sum_{j=1}^{N_n}
\zeta_{\theta_{n,h},\vartheta_n(\mathbf r)}(S_j).
\]

We record the pilot concentration bound used in this averaging. Define
\[
\ell_0=\frac{qd}{e^\varepsilon+3},
\qquad
\ell_1=\frac{pd}{e^\varepsilon+3},
\qquad
\mathcal B_n(\varrho)
=
\frac1{m_n(\varrho\ell_0)^2}
+\frac1{m_n(\varrho\ell_1)^2},
\quad \varrho>0,\ n\ge2.
\]
For an interior parameter \(\xi\), the means of the pilot indicators for categories \(1\) and \(3\) are
\[
\frac{1+dq\xi_0}{e^\varepsilon+3},
\qquad
\frac{1+dp\xi_1}{e^\varepsilon+3}.
\]
Each indicator has variance at most one, so independence and Chebyshev's inequality give, for \(k\in\{1,3\}\) and \(b>0\),
\[
\Pr_\xi\{|\widehat u_k-\mathbb E_\xi\widehat u_k|\ge b\}
\le\frac1{m_nb^2}.
\]
Frequency inversion scales the two coordinate errors by
\(\ell_0^{-1}\) and \(\ell_1^{-1}\). If
\[
\delta_n+\varrho\le
\min\{\xi_0,1-\xi_0,\xi_1,1-\xi_1\},
\]
then an inverted coordinate within \(\varrho\) of \(\xi\) lies inside the clipping interval. Taking the contrapositive coordinatewise and applying the union bound yields
\[
\Pr_\xi\{\|\widetilde\theta-\xi\|_\infty\ge\varrho\}
\le\mathcal B_n(\varrho),
\qquad
\|\widetilde\theta-\xi\|_\infty
=\max_{k\in\{0,1\}}|\widetilde\theta_k-\xi_k|.
\]

Define, for \(n\ge2\),
\[
\beta_n=m_n^{-1/4},
\qquad
\varrho_{n,h}
=
2\max\left\{
\beta_n,\sqrt{\|\theta_{n,h}-\theta\|_\infty}
\right\},
\]
and
\[
\Delta_\theta
=\min\{\theta_0,1-\theta_0,\theta_1,1-\theta_1\}>0.
\]
Pilot divergence and local convergence imply
\[
\beta_n\to0,\qquad \varrho_{n,h}\to0.
\]
Eventually,
\[
\beta_n<\Delta_\theta/4,\qquad
\delta_n<\Delta_\theta/4,\qquad
\|\theta_{n,h}-\theta\|_\infty<\Delta_\theta/4,
\]
so the concentration bound applies at \(\xi=\theta_{n,h}\) with radius \(\beta_n\). Its right-hand side becomes
\[
\mathcal B_n(\beta_n)
=
\frac1{\sqrt{m_n}}
\left(\frac1{\ell_0^2}+\frac1{\ell_1^2}\right)
\longrightarrow0.
\]
Moreover, eventually
\[
\|\theta_{n,h}-\theta\|_\infty
\le\sqrt{\|\theta_{n,h}-\theta\|_\infty}.
\]
Hence
\[
\|\vartheta_n(\mathbf r)-\theta_{n,h}\|_\infty<\beta_n
\quad\Longrightarrow\quad
\|\vartheta_n(\mathbf r)-\theta\|_\infty\le\varrho_{n,h}.
\]

The deterministic-row convergence is uniform over these shrinking pilot neighborhoods. Explicitly, for every \(a>0\), eventually
\[
\|\vartheta_n(\mathbf r)-\theta\|_\infty\le\varrho_{n,h}
\quad\Longrightarrow\quad
|\mathcal C_{n,h}(x;\vartheta_n(\mathbf r))
-\mathcal G_\theta(x)|<a.
\]
If this implication failed, one could choose strictly increasing \(n_k\) and pilot vectors \(\mathbf r_k\) such that
\[
\|\vartheta_{n_k}(\mathbf r_k)-\theta\|_\infty
\le\varrho_{n_k,h}\longrightarrow0,
\qquad
|\mathcal C_{n_k,h}(x;\vartheta_{n_k}(\mathbf r_k))
-\mathcal G_\theta(x)|\ge a.
\]
This contradicts the deterministic-row limit.

Both distribution functions lie in \([0,1]\). Splitting the pilot average into the displayed neighborhood and its complement therefore gives, eventually,
\[
\left|
\sum_{\mathbf r}w_{n,h}(\mathbf r)
\mathcal C_{n,h}(x;\vartheta_n(\mathbf r))
-\mathcal G_\theta(x)
\right|
\le a+\mathcal B_n(\beta_n).
\]
Letting \(n\to\infty\), then \(a\downarrow0\), proves
\[
\Pr_{\theta_{n,h}}\{U_{n,h}^{\mathcal P}\le x\}
\longrightarrow\mathcal G_\theta(x)
\qquad(h\in\mathbb R^2,\ x\in\mathbb R).
\]
% lean: CausalSmith.Stat.LdpAteEfficiencySurface.pilot_scaledError_cdf_tendsto_of_all_conditionalRows

\item \emph{The weak local limit.}
For an interior parameter \(\xi\), the transcript law is a probability law: it is a mixture of conditional transcript probability laws, with nonnegative input-path weights satisfying
\[
\sum_{(j_1,\ldots,j_n)\in\{0,1,2,3\}^n}
\prod_{i=1}^n\pi_{\xi,j_i}
=\left(\sum_{j=0}^3\pi_{\xi,j}\right)^n=1.
\]
This also holds at \(n=0\), with the empty-product convention.

For each fixed \(h\), define a Borel probability law on \(\mathbb R\) by
\[
\nu_{n,h}=
\begin{cases}
\mathcal L_{\theta_{n,h}}(U_{n,h}^{\mathcal P}),
&\theta_{n,h}\in(0,1)^2,\\
N(0,V^*(\theta,p,\varepsilon)),
&\theta_{n,h}\notin(0,1)^2.
\end{cases}
\]
The local parameter is eventually interior, so the preceding distribution-function limit applies to \(\nu_{n,h}\). Convergence at every point of the continuous Gaussian distribution function implies weak convergence. Equivalently, for every bounded continuous \(f:\mathbb R\to\mathbb R\),
\[
\int f(u)\,d\nu_{n,h}(u)
\longrightarrow
\int f(u)\,dN(0,V^*(\theta,p,\varepsilon))(u).
\]
Eventually the left-hand side is
\(\mathbb E_{\theta_{n,h}}[f(U_{n,h}^{\mathcal P})]\), proving the asserted common weak local limit. Taking \(h=0\) gives the stated convergence at \(\theta\).
% lean: CausalSmith.Stat.LdpAteEfficiencySurface.weakLocalLimit_of_scaledError_cdf_tendsto

\item \emph{Consistency of the empirical variance.}
The empirical variance is the quantity in \cref{obj:synth_1}. Fix \(\delta>0\). Continuity of \(\sigma^2\) at \((\theta,\theta)\) supplies \(R_{\mathrm{var}}>0\) such that
\[
\|\vartheta-\theta\|_\infty<R_{\mathrm{var}},
\quad \vartheta\in(0,1)^2
\quad\Longrightarrow\quad
|\sigma^2(\theta,\vartheta)-V^*(\theta,p,\varepsilon)|<\delta/2.
\]
Set
\[
\varrho_{\mathrm{var}}
=\min\{R_{\mathrm{var}}/2,\Delta_\theta/4\}>0.
\]
Eventually \(\delta_n<\Delta_\theta/2\), so
\[
\delta_n+\varrho_{\mathrm{var}}\le\Delta_\theta,
\qquad
\Pr_\theta\{\|\widetilde\theta-\theta\|_\infty
\ge\varrho_{\mathrm{var}}\}
\le\mathcal B_n(\varrho_{\mathrm{var}})\longrightarrow0.
\]

Condition on a pilot vector, and write its estimate as \(\vartheta\). Define the centered main-sample moments
\[
\overline\zeta_n
=\frac1{N_n}\sum_{j=1}^{N_n}\zeta_{\theta,\vartheta}(S_j),
\qquad
\overline{\zeta^2}_n
=\frac1{N_n}\sum_{j=1}^{N_n}\zeta_{\theta,\vartheta}(S_j)^2.
\]
Translation invariance of empirical variance gives
\[
\widehat V^*=\overline{\zeta^2}_n-\overline\zeta_n^{\,2}.
\]
Conditional independence, centering, and boundedness give
\[
\mathbb E_\theta[\overline\zeta_n^{\,2}\mid H_{m_n}]
=\frac{\sigma^2(\theta,\vartheta)}{N_n}
\le\frac{(2\overline\phi)^2}{N_n},
\]
and
\[
\operatorname{Var}_\theta(\overline{\zeta^2}_n\mid H_{m_n})
\le\frac{(2\overline\phi)^4}{N_n}.
\]
On a pilot with
\(|\sigma^2(\theta,\vartheta)-V^*(\theta,p,\varepsilon)|<\delta/2\),
the event \(|\widehat V^*-V^*(\theta,p,\varepsilon)|>\delta\) implies
\[
|\overline{\zeta^2}_n-\sigma^2(\theta,\vartheta)|\ge\delta/4
\quad\text{or}\quad
\overline\zeta_n^{\,2}\ge\delta/4.
\]
Chebyshev's and Markov's inequalities consequently bound its conditional probability by
\[
\frac{(2\overline\phi)^4}{N_n(\delta/4)^2}
+\frac{(2\overline\phi)^2}{N_n(\delta/4)}.
\]
On the complementary pilots the conditional probability is at most one. Averaging gives, eventually,
\[
\begin{aligned}
\Pr_\theta\{|\widehat V^*-V^*(\theta,p,\varepsilon)|>\delta\}
\le{}&
\mathcal B_n(\varrho_{\mathrm{var}})
+\frac{(2\overline\phi)^4}{N_n(\delta/4)^2}\\
&+\frac{(2\overline\phi)^2}{N_n(\delta/4)}
\longrightarrow0.
\end{aligned}
\]
% lean: CausalSmith.Stat.LdpAteEfficiencySurface.pilot_VhatStar_consistent

\item \emph{Uniform squared tails and regularity.}
Fix \(H>0\), and use the admissible perturbation set of \cref{obj:synth_6}. For sufficiently large \(n\),
\[
\frac n{N_n}\le2.
\]
For every admissible \(h\), condition on the pilot and denote the \(N_n\) centered corrections by
\[
Z_j=\zeta_{\theta_{n,h},\widetilde\theta}(S_j),
\qquad 1\le j\le N_n.
\]
They are conditionally independent and identically distributed, with mean zero and
\[
|Z_j|\le2\overline\phi,
\qquad
U_{n,h}^{\mathcal P}
=\frac{\sqrt n}{N_n}\sum_{j=1}^{N_n}Z_j.
\]
Define their conditional moments by
\[
\mathfrak m_2=\mathbb E_{\theta_{n,h}}[Z_1^2\mid H_{m_n}],
\qquad
\mathfrak m_4=\mathbb E_{\theta_{n,h}}[Z_1^4\mid H_{m_n}].
\]
Then
\[
0\le\mathfrak m_2\le(2\overline\phi)^2,
\qquad
0\le\mathfrak m_4\le(2\overline\phi)^4.
\]
In the fourth-power expansion, terms having an index that occurs exactly once have zero conditional expectation. The remaining terms give
\[
\mathbb E_{\theta_{n,h}}
\left[\left(\sum_{j=1}^{N_n}Z_j\right)^4\middle|H_{m_n}\right]
=
N_n\mathfrak m_4+3N_n(N_n-1)\mathfrak m_2^2.
\]
Since \(N_n\ge1\),
\[
\mathbb E_{\theta_{n,h}}
\left[\left(N_n^{-1/2}\sum_{j=1}^{N_n}Z_j\right)^4
\middle|H_{m_n}\right]
\le4(2\overline\phi)^4.
\]
Thus, uniformly over the pilot and admissible \(h\),
\[
\mathbb E_{\theta_{n,h}}
[(U_{n,h}^{\mathcal P})^4\mid H_{m_n}]
\le
4\left(\frac n{N_n}\right)^2(2\overline\phi)^4
\le16(2\overline\phi)^4.
\]
Averaging preserves this bound. For every integer \(M>0\), the pointwise inequality
\[
u^2\mathbf1\{|u|>M\}\le\frac{u^4}{M^2},
\qquad u\in\mathbb R,
\]
therefore yields
\[
0\le
\limsup_{n\to\infty}
\sup_{h\in\mathcal H_{n,H}(\theta)}
\mathbb E_{\theta_{n,h}}
\left[(U_{n,h}^{\mathcal P})^2
\mathbf1\{|U_{n,h}^{\mathcal P}|>M\}\right]
\le\frac{16(2\overline\phi)^4}{M^2}.
\]
Letting \(M\to\infty\) proves the required uniform squared-tail convergence.

The Gaussian law
\[
L_{\mathcal P}=N(0,V^*(\theta,p,\varepsilon))
\]
is a probability law with
\[
\int u^2\,dL_{\mathcal P}(u)<\infty.
\]
Together with the common weak local limit and the squared-tail bound, this verifies the regularity conditions in \cref{obj:synth_6}.
% lean: CausalSmith.Stat.LdpAteEfficiencySurface.regular_of_weakLocalLimit_and_uniformTails

\item \emph{Standardizing the limiting distribution function.}
The reference-design argument gives
\[
J^*(\theta,p,\varepsilon)>0,
\qquad
V^*(\theta,p,\varepsilon)>0,
\qquad
\sqrt{V^*(\theta,p,\varepsilon)}>0.
\]
Define, for each sample size,
\[
\mathsf T_n
=\frac{U_{n,0}^{\mathcal P}}
{\sqrt{V^*(\theta,p,\varepsilon)}}.
\]
Since \(\theta_{n,0}=\theta\), positivity of the denominator gives, for every \(x\in\mathbb R\),
\[
\{\mathsf T_n\le x\}
=
\left\{U_{n,0}^{\mathcal P}
\le x\sqrt{V^*(\theta,p,\varepsilon)}\right\}.
\]
The Gaussian scaling identity is
\[
N(0,V^*(\theta,p,\varepsilon))
\left((-\infty,x\sqrt{V^*(\theta,p,\varepsilon)}]\right)
=\Phi(x).
\]
Applying the local distribution-function limit at this threshold proves
\[
\Pr_\theta\{\mathsf T_n\le x\}\longrightarrow\Phi(x)
\qquad(x\in\mathbb R).
\]
% lean: CausalSmith.Stat.LdpAteEfficiencySurface.pilot_attainment

\item \emph{Transfer to the Wald pivot.}
First, consistency also controls non-strict deviation events. For every \(b>0\),
\[
\{|\widehat V^*-V^*(\theta,p,\varepsilon)|\ge b\}
\subseteq
\{|\widehat V^*-V^*(\theta,p,\varepsilon)|>b/2\},
\]
and hence
\[
\Pr_\theta\{|\widehat V^*-V^*(\theta,p,\varepsilon)|\ge b\}
\longrightarrow0.
\]
The normal distribution function is continuous and strictly increasing, with limits zero and one at the two ends of the real line. The critical-value definition thus gives
\[
\Phi(z_{1-\gamma/2})=1-\gamma/2.
\]
Normal symmetry gives \(\Phi(0)=1/2\). Since \(1-\gamma/2>1/2\),
\[
z_{1-\gamma/2}>0.
\]

Define the estimation error and fallback pivot, for \(n\ge1\), by
\[
\Delta_{\tau,n}=\widehat\tau^*-\tau,
\qquad
\mathsf W_n=
\begin{cases}
0,&\widehat V^*=0,\ \Delta_{\tau,n}=0,\\
z_{1-\gamma/2}+1,&\widehat V^*=0,\ \Delta_{\tau,n}\ne0,\\
\sqrt n\,\Delta_{\tau,n}/\sqrt{\widehat V^*},
&\widehat V^*>0.
\end{cases}
\]
The empirical variance is nonnegative on every transcript, since it is a nonnegative multiple of a sum of squares. When it is positive, ordinary square-root algebra gives
\[
|\mathsf W_n|\le z_{1-\gamma/2}
\quad\Longleftrightarrow\quad
|\Delta_{\tau,n}|
\le z_{1-\gamma/2}\sqrt{\widehat V^*/n}.
\]
When it is zero, both events hold precisely when \(\Delta_{\tau,n}=0\), by the two fallback branches. Thus the equivalence holds on every transcript.

We show that \(\mathsf W_n-\mathsf T_n\) converges to zero in probability. The standardized distribution-function limit first gives tightness. Explicitly, for any \(\eta>0\), define
\[
\eta_{\mathrm{tail}}=\min\{\eta/8,1/8\}.
\]
Choose \(R_{\mathrm{tight}}>0\) so that the standard normal probabilities of
\((-\infty,-R_{\mathrm{tight}}]\) and
\([R_{\mathrm{tight}},\infty)\) are each at most
\(\eta_{\mathrm{tail}}\). Distribution-function convergence gives weak convergence of \(\mathsf T_n\), and the two endpoints have zero normal mass. Consequently,
\[
\Pr_\theta\{\mathsf T_n\notin[-R_{\mathrm{tight}},R_{\mathrm{tight}}]\}
\longrightarrow
1-\Phi(R_{\mathrm{tight}})+\Phi(-R_{\mathrm{tight}})
\le2\eta_{\mathrm{tail}}<\eta.
\]
Taking
\[
M_{\mathrm{tight}}=R_{\mathrm{tight}}+1
\]
therefore gives, eventually,
\[
\Pr_\theta\{M_{\mathrm{tight}}\le|\mathsf T_n|\}<\eta.
\]

Now fix \(\delta>0\) and \(\eta>0\), and choose such a positive
\(M_{\mathrm{tight}}\) with tail probability eventually below \(\eta/2\).
Define
\[
\mathfrak g_{\mathrm{scale}}(u)
=\frac{\sqrt{V^*(\theta,p,\varepsilon)}}{\sqrt u},
\qquad u>0.
\]
This function is continuous at \(V^*(\theta,p,\varepsilon)\), where it equals one. Choose \(\varrho_{\mathrm{scale},0}>0\) so that
\[
u>0,\quad |u-V^*(\theta,p,\varepsilon)|<\varrho_{\mathrm{scale},0}
\quad\Longrightarrow\quad
|\mathfrak g_{\mathrm{scale}}(u)-1|
<\delta/M_{\mathrm{tight}},
\]
and set
\[
\varrho_{\mathrm{scale}}
=
\min\{\varrho_{\mathrm{scale},0}/2,
V^*(\theta,p,\varepsilon)/2\}>0.
\]
On the event
\[
|\widehat V^*-V^*(\theta,p,\varepsilon)|
<\varrho_{\mathrm{scale}},
\qquad
|\mathsf T_n|<M_{\mathrm{tight}},
\]
the estimated variance is positive and
\[
\mathsf W_n
=\mathsf T_n\mathfrak g_{\mathrm{scale}}(\widehat V^*),
\qquad
|\mathsf W_n-\mathsf T_n|
\le M_{\mathrm{tight}}
|\mathfrak g_{\mathrm{scale}}(\widehat V^*)-1|
<\delta.
\]
It follows that
\[
\begin{aligned}
\Pr_\theta\{|\mathsf W_n-\mathsf T_n|\ge\delta\}
\le{}&
\Pr_\theta\{|\widehat V^*-V^*(\theta,p,\varepsilon)|
\ge\varrho_{\mathrm{scale}}\}\\
&+\Pr_\theta\{|\mathsf T_n|\ge M_{\mathrm{tight}}\}
<\eta
\end{aligned}
\]
eventually. Thus this deviation probability tends to zero.

For every \(x\in\mathbb R\) and \(\delta>0\),
\[
\begin{aligned}
\Pr_\theta\{\mathsf T_n\le x-\delta\}
-\Pr_\theta\{|\mathsf W_n-\mathsf T_n|\ge\delta\}
&\le\Pr_\theta\{\mathsf W_n\le x\}\\
&\le\Pr_\theta\{\mathsf T_n\le x+\delta\}
+\Pr_\theta\{|\mathsf W_n-\mathsf T_n|\ge\delta\}.
\end{aligned}
\]
Taking limits and then letting \(\delta\downarrow0\) proves convergence of the pivot distribution functions to \(\Phi\). The closed interval
\([-z_{1-\gamma/2},z_{1-\gamma/2}]\) has zero normal boundary mass, so
\[
\Pr_\theta\{|\mathsf W_n|\le z_{1-\gamma/2}\}
\longrightarrow
N(0,1)([-z_{1-\gamma/2},z_{1-\gamma/2}]).
\]
The event equivalence transfers this limit to the stated Wald event.
% lean: CausalSmith.Stat.LdpAteEfficiencySurface.pilot_wald_coverage_of_standardized_cdf_and_VhatStar_consistent

\item \emph{Privacy at every sample size and history.}
The randomized-response probabilities in \cref{obj:synth_10} are normalized, and for every pair of private records and every reported category,
\[
Q_{\mathrm{RR}}(k\mid x_j)
\le e^\varepsilon Q_{\mathrm{RR}}(k\mid x_{j'}).
\]
Summing over an output event proves the pilot privacy inequality.

At any main stage, including at an unattainable history, the history-based parameter estimate clips both coordinates to
\[
[\delta_n,1-\delta_n]\subset(0,1),
\qquad
0<\delta_n=(m_n+2)^{-1}\le1/2.
\]
The selected weights are therefore feasible. For every pattern \(S\),
\[
Q_{\alpha_\vartheta}(S\mid x_j)
=\alpha_{\vartheta,S}b_S(j),
\qquad
\sum_SQ_{\alpha_\vartheta}(S\mid x_j)=1,
\]
and
\[
\alpha_{\vartheta,S}b_S(j)
\le e^\varepsilon\alpha_{\vartheta,S}b_S(j').
\]
This inequality also holds for zero weights. Summing over patterns proves privacy for every measurable main-release event. At each fixed history, the same selected channel is used on the two compared private records. The two stage cases therefore verify the Markov property and the all-history inequality in \cref{obj:ass:sequential-ldp}. At \(n=0\), the stage conditions are vacuous.
% lean: CausalSmith.Stat.LdpAteEfficiencySurface.pilotProtocol_private

\item \emph{Exact pilot-adapted unbiasedness.}
Fix \(n\ge2\). Factorization gives, for a pilot vector \(\mathbf r\) and a main pattern vector
\(\mathbf s=(s_1,\ldots,s_{N_n})\in\mathcal S^{N_n}\),
\[
\Pr_\theta\{H_{m_n}=\mathbf r,\ (S_{m_n+1},\ldots,S_n)=\mathbf s\}
=
w_{n,0}(\mathbf r)
\prod_{j=1}^{N_n}
\omega_{\theta,\vartheta_n(\mathbf r)}(s_j).
\]
Transcripts with a release type incompatible with its stage have zero mass.

For each fixed pilot, product normalization and the exact correction mean give
\[
\begin{aligned}
&\sum_{\mathbf s\in\mathcal S^{N_n}}
\left(\prod_{j=1}^{N_n}
\omega_{\theta,\vartheta_n(\mathbf r)}(s_j)\right)
\left\{
c^\top\vartheta_n(\mathbf r)
+\frac1{N_n}\sum_{j=1}^{N_n}
\phi_{\vartheta_n(\mathbf r)}(s_j)
-c^\top\theta
\right\}\\
&\hspace{2cm}
=c^\top\vartheta_n(\mathbf r)
+\frac{N_n}{N_n}
c^\top(\theta-\vartheta_n(\mathbf r))
-c^\top\theta
=0.
\end{aligned}
\]
Multiplication by \(w_{n,0}(\mathbf r)\) shows that the integral of
\(\widehat\tau^*-\tau\) over every pilot-prefix cell is zero. A prefix incompatible with the pilot release type has zero mass, so its cell integral is also zero.

A function \(f\) satisfying the statement is constant on each such cell. Summing the zero cell integrals, each multiplied by its corresponding value of \(f\), gives
\[
\mathbb E_\theta[(\widehat\tau^*-\tau)f]=0.
\]
The transcript space is finite and the estimator and \(f\) are real-valued, so their product is integrable under the transcript probability law.
% lean: CausalSmith.Stat.LdpAteEfficiencySurface.pilot_conditional_unbiasedness

\item \emph{Evaluation of the limiting coverage.}
Normal symmetry and the critical-value identity give
\[
\Phi(-z_{1-\gamma/2})
=1-\Phi(z_{1-\gamma/2})
=\gamma/2.
\]
The standard normal law has zero mass at the interval endpoints. Hence
\[
N(0,1)([-z_{1-\gamma/2},z_{1-\gamma/2}])
=
\Phi(z_{1-\gamma/2})-\Phi(-z_{1-\gamma/2})
=
1-\gamma.
\]
Substitution into the Wald-event limit proves
\[
\Pr_\theta\left\{
|\widehat\tau^*-\tau|
\le z_{1-\gamma/2}\sqrt{\widehat V^*/n}
\right\}\longrightarrow1-\gamma,
\]
completing all the asserted conclusions.
% lean: CausalSmith.Stat.LdpAteEfficiencySurface.gaussianMeasure_waldCriticalValue_Icc
\end{enumerate}
\end{proof}

\begin{proof}[Proof of \cref{obj:thm:five-output-upper-bound}]
\begin{enumerate}
\item
Use the record and pattern indexing of \cref{obj:synth_8}. For every \(S\in\mathcal S\), define
\[
\begin{aligned}
b_{S,j}
&:=1+(e^\varepsilon-1)\mathbf 1_{\{j\in S\}},
&&j\in\{0,1,2,3\},\\
\mathbf b_S&:=(b_{S,0},b_{S,1},b_{S,2},b_{S,3})^\top,\\
g_S
&:=(e^\varepsilon-1)
\begin{pmatrix}
(1-p)(\mathbf 1_{\{1\in S\}}-\mathbf 1_{\{0\in S\}})\\
p(\mathbf 1_{\{3\in S\}}-\mathbf 1_{\{2\in S\}})
\end{pmatrix},\\
h_S(\theta)&:=\sum_{j=0}^{3}\pi_{\theta,j}b_{S,j},
\qquad
v(t):=(t,t+1)^\top,\quad t\in\mathbb R.
\end{aligned}
\]
The record probabilities are
\[
(\pi_{\theta,0},\pi_{\theta,1},\pi_{\theta,2},\pi_{\theta,3})
=
\bigl((1-p)(1-\mu_0),(1-p)\mu_0,
p(1-\mu_1),p\mu_1\bigr).
\]
The interior hypotheses make each probability positive, and
\(b_{S,j}\ge1\). Consequently,
\[
h_S(\theta)>0\qquad(S\in\mathcal S).
\]
Write
\[
\begin{aligned}
\mathcal A_\varepsilon
&=\left\{\beta\in\mathbb R^{\mathcal S}:
\beta_S\ge0,\ 
\sum_{S\in\mathcal S}\beta_S\mathbf b_S=\mathbf1
\right\},
\qquad \mathbf1:=(1,1,1,1)^\top,\\
f_S(t)&:=\frac{(g_S^\top v(t))^2}{h_S(\theta)},
\qquad
F_\theta(\beta,t)=\sum_{S\in\mathcal S}\beta_Sf_S(t),\\
E(t)&:=\sup_{\beta\in\mathcal A_\varepsilon}F_\theta(\beta,t),
\qquad t\in\mathbb R.
\end{aligned}
\]
By \cref{obj:thm:finite-oracle}, there is \(t^*\in\mathbb R\) such that
\[
E(t^*)=\inf_{t\in\mathbb R}E(t)
=J^*(\theta,p,\varepsilon).
\]
% lean: CausalSmith.Stat.LdpAteEfficiencySurface.exists_sparse_staircase_saddle

\item
We first obtain a weight vector attaining the profiled optimum.
The complementary patterns \(S_1=\{0\}\) and \(S_{14}=\{1,2,3\}\)
give the feasible vector
\[
\alpha^{\mathrm{base}}_S:=
\begin{cases}
(e^\varepsilon+1)^{-1},&S\in\{S_1,S_{14}\},\\
0,&\text{otherwise},
\end{cases}
\qquad
\sum_{S\in\mathcal S}\alpha^{\mathrm{base}}_S\mathbf b_S
=
\frac{\mathbf b_{S_1}+\mathbf b_{S_{14}}}{e^\varepsilon+1}
=\mathbf1.
\]
Thus \(\mathcal A_\varepsilon\) is nonempty. It is closed, since its
constraints are finitely many continuous equalities and inequalities.
Moreover, for every \(\beta\in\mathcal A_\varepsilon\) and \(S\in\mathcal S\),
\[
0\le\beta_S
\le\sum_{T\in\mathcal S}\beta_T
\le\sum_{T\in\mathcal S}\beta_Tb_{T,0}
=1.
\]
Hence \(\mathcal A_\varepsilon\) is compact.

Define the profiled objective on this set by
\[
\Phi(\beta):=\inf_{t\in\mathbb R}F_\theta(\beta,t),
\qquad \beta\in\mathcal A_\varepsilon.
\]
Nonnegativity of the weights and positivity of the denominators give
\[
0\le\Phi(\beta)\le F_\theta(\beta,0)<\infty.
\]
For each fixed \(t\), \(F_\theta(\,\cdot\,,t)\) is continuous.
Its infimum over \(t\) is therefore upper semicontinuous.
The maximum of \(\Phi\) on the nonempty compact feasible set is attained.
Using the optimized value in \cref{obj:def:staircase-saddle}, choose
\(\alpha^{\mathrm{opt}}\in\mathcal A_\varepsilon\) with
\[
\Phi(\alpha^{\mathrm{opt}})
=J^*(\theta,p,\varepsilon).
\]
In particular,
\[
J^*(\theta,p,\varepsilon)
\le F_\theta(\alpha^{\mathrm{opt}},t)
\qquad(t\in\mathbb R).
\]
% lean: CausalSmith.Stat.LdpAteEfficiencySurface.Jstar_exists_optimal_weight

\item
Continuity and compactness also supply a maximizer at the fixed
coordinate \(t^*\): choose
\[
\alpha^{\mathrm{env}}
\in\operatorname*{arg\,max}_{\beta\in\mathcal A_\varepsilon}
F_\theta(\beta,t^*).
\]
Then
\[
E(t^*)=F_\theta(\alpha^{\mathrm{env}},t^*),
\]
and the preceding lower bound yields
\[
J^*(\theta,p,\varepsilon)
\le F_\theta(\alpha^{\mathrm{opt}},t^*)
\le F_\theta(\alpha^{\mathrm{env}},t^*)
=E(t^*)
=J^*(\theta,p,\varepsilon).
\]
Therefore
\[
F_\theta(\alpha^{\mathrm{opt}},t^*)
=J^*(\theta,p,\varepsilon),
\qquad
F_\theta(\beta,t^*)
\le F_\theta(\alpha^{\mathrm{opt}},t^*)
\quad(\beta\in\mathcal A_\varepsilon).
\]
% lean: CausalSmith.Stat.LdpAteEfficiencySurface.informationObjective_exists_maximizer

\item
Define, for \(S\in\mathcal S\), \(t\in\mathbb R\), and
\(\beta\in\mathbb R^{\mathcal S}\),
\[
\begin{aligned}
\dot f_S(t)
&:=\frac{2(g_S^\top v(t))(g_{S,0}+g_{S,1})}{h_S(\theta)},\\
D_\beta(t)&:=\sum_{S\in\mathcal S}\beta_S\dot f_S(t),\\
C_\beta&:=\sum_{S\in\mathcal S}
\beta_S\frac{(g_{S,0}+g_{S,1})^2}{h_S(\theta)}.
\end{aligned}
\]
Expanding each square gives the identity
\[
F_\theta(\beta,t+\lambda)-F_\theta(\beta,t)
=
\lambda D_\beta(t)+\lambda^2C_\beta
\qquad(t,\lambda\in\mathbb R).
\]
The preceding step shows that \(t^*\) minimizes
\(F_\theta(\alpha^{\mathrm{opt}},\,\cdot\,)\).
If \(D_{\alpha^{\mathrm{opt}}}(t^*)\ne0\), take
\[
\lambda_{\mathrm{stat}}
:=
-\frac{D_{\alpha^{\mathrm{opt}}}(t^*)}
{2(|C_{\alpha^{\mathrm{opt}}}|+1)}.
\]
The quadratic identity then gives
\[
\begin{aligned}
0
&\le
F_\theta(\alpha^{\mathrm{opt}},t^*+\lambda_{\mathrm{stat}})
-F_\theta(\alpha^{\mathrm{opt}},t^*)\\
&=
\frac{D_{\alpha^{\mathrm{opt}}}(t^*)^2}
{4(|C_{\alpha^{\mathrm{opt}}}|+1)^2}
\left[
C_{\alpha^{\mathrm{opt}}}
-2(|C_{\alpha^{\mathrm{opt}}}|+1)
\right]
<0,
\end{aligned}
\]
where the strict inequality uses
\(C_{\alpha^{\mathrm{opt}}}\le|C_{\alpha^{\mathrm{opt}}}|\).
This contradiction proves
\[
D_{\alpha^{\mathrm{opt}}}(t^*)=0.
\]
% lean: CausalSmith.Stat.LdpAteEfficiencySurface.informationObjective_stationary_of_min

\item
For every real weight vector, define its indexed support by
\[
\operatorname{supp}(\beta)
:=\{S\in\mathcal S:\beta_S\ne0\},
\qquad \beta\in\mathbb R^{\mathcal S}.
\]
Consider
\[
\mathcal W:=
\left\{\beta\in\mathcal A_\varepsilon:
D_\beta(t^*)=D_{\alpha^{\mathrm{opt}}}(t^*),\
F_\theta(\beta,t^*)=F_\theta(\alpha^{\mathrm{opt}},t^*)
\right\}.
\]
This set contains \(\alpha^{\mathrm{opt}}\).
Choose a member of minimum support cardinality:
\[
\alpha^{\mathrm{sp}}
\in\operatorname*{arg\,min}_{\beta\in\mathcal W}
|\operatorname{supp}(\beta)|,
\qquad
\mathcal K_{\mathrm{sp}}:=\operatorname{supp}(\alpha^{\mathrm{sp}}).
\]
Such a choice exists because support cardinalities are nonnegative
integers. By construction,
\[
D_{\alpha^{\mathrm{sp}}}(t^*)=0,
\qquad
F_\theta(\alpha^{\mathrm{sp}},t^*)=J^*(\theta,p,\varepsilon),
\]
and
\[
F_\theta(\beta,t^*)\le F_\theta(\alpha^{\mathrm{sp}},t^*)
\qquad(\beta\in\mathcal A_\varepsilon).
\]

Suppose that \(|\mathcal K_{\mathrm{sp}}|>5\).
The more than five columns
\[
\begin{pmatrix}\mathbf b_S\\ \dot f_S(t^*)\end{pmatrix}
\in\mathbb R^5,
\qquad S\in\mathcal K_{\mathrm{sp}},
\]
are linearly dependent. Extending their dependence coefficients by zero
outside \(\mathcal K_{\mathrm{sp}}\) gives
\[
\begin{gathered}
\Delta=(\Delta_S)_{S\in\mathcal S}
\in\mathbb R^{\mathcal S}\setminus\{0\},
\qquad
\Delta_S=0\quad(S\notin\mathcal K_{\mathrm{sp}}),\\
\sum_{S\in\mathcal S}\Delta_S\mathbf b_S=0,
\qquad
\sum_{S\in\mathcal S}\Delta_S\dot f_S(t^*)=0.
\end{gathered}
\]
% lean: Causalean.Mathlib.Optimization.exists_kernel_direction

\item
The objective moment of this direction also vanishes.
To see this, define
\[
\delta_S:=
\begin{cases}
1,&\alpha^{\mathrm{sp}}_S=0,\\
\displaystyle\frac{\alpha^{\mathrm{sp}}_S}{|\Delta_S|+1},
&\alpha^{\mathrm{sp}}_S>0,
\end{cases}
\qquad
\delta:=\min_{S\in\mathcal S}\delta_S.
\]
Each \(\delta_S\) is positive. Since \(\mathcal S\) is finite and
nonempty, and \(\Delta_S=0\) wherever \(\alpha^{\mathrm{sp}}_S=0\),
\[
\delta>0,
\qquad
\delta|\Delta_S|\le\alpha^{\mathrm{sp}}_S
\quad(S\in\mathcal S).
\]
Thus both perturbations are feasible:
\[
\alpha^{\mathrm{sp}}\pm\delta\Delta\ge0,
\qquad
\sum_{S\in\mathcal S}
(\alpha^{\mathrm{sp}}_S\pm\delta\Delta_S)\mathbf b_S
=\mathbf1.
\]
Their objective values satisfy
\[
F_\theta(\alpha^{\mathrm{sp}}\pm\delta\Delta,t^*)
=
F_\theta(\alpha^{\mathrm{sp}},t^*)
\pm\delta\sum_{S\in\mathcal S}\Delta_Sf_S(t^*)
\le F_\theta(\alpha^{\mathrm{sp}},t^*).
\]
Using both signs and \(\delta>0\) proves
\[
\sum_{S\in\mathcal S}\Delta_Sf_S(t^*)=0.
\]
% lean: Causalean.Mathlib.Optimization.objective_orthogonal_of_maximal

\item
We now move to a boundary point with smaller support.
Choose \(S_{\mathrm{pivot}}\in\mathcal S\) with
\(\Delta_{S_{\mathrm{pivot}}}\ne0\), and define
\[
\sigma:=
\begin{cases}
1,&\Delta_{S_{\mathrm{pivot}}}>0,\\
-1,&\Delta_{S_{\mathrm{pivot}}}<0,
\end{cases}
\qquad
\Delta^{\mathrm{mov}}_S:=\sigma\Delta_S,
\qquad
\mathcal K_+:=\{S\in\mathcal S:\Delta^{\mathrm{mov}}_S>0\}.
\]
Then \(\mathcal K_+\) is nonempty. Define
\[
\lambda_{\mathrm{bdry}}
:=
\min_{S\in\mathcal K_+}
\frac{\alpha^{\mathrm{sp}}_S}{\Delta^{\mathrm{mov}}_S}
\ge0,
\qquad
\alpha^{\mathrm{bdry}}_S
:=
\alpha^{\mathrm{sp}}_S
-\lambda_{\mathrm{bdry}}\Delta^{\mathrm{mov}}_S.
\]
For \(S\in\mathcal K_+\), the defining minimum gives
\[
\lambda_{\mathrm{bdry}}\Delta^{\mathrm{mov}}_S
\le\alpha^{\mathrm{sp}}_S.
\]
For \(S\notin\mathcal K_+\),
\[
\alpha^{\mathrm{bdry}}_S
\ge\alpha^{\mathrm{sp}}_S\ge0.
\]
Hence \(\alpha^{\mathrm{bdry}}\ge0\).
Choose an index attaining the minimum, so that
\[
S_{\mathrm{hit}}\in\mathcal K_+,
\qquad
\lambda_{\mathrm{bdry}}
=
\frac{\alpha^{\mathrm{sp}}_{S_{\mathrm{hit}}}}
{\Delta^{\mathrm{mov}}_{S_{\mathrm{hit}}}}.
\]
The support condition on \(\Delta\) implies
\(\alpha^{\mathrm{sp}}_{S_{\mathrm{hit}}}>0\), whereas
\[
\alpha^{\mathrm{bdry}}_{S_{\mathrm{hit}}}=0.
\]
Coordinates outside \(\mathcal K_{\mathrm{sp}}\) remain zero.
Consequently,
\[
\operatorname{supp}(\alpha^{\mathrm{bdry}})
\subsetneq\operatorname{supp}(\alpha^{\mathrm{sp}}).
\]

The three vanishing moments established above give
\[
\begin{aligned}
\sum_{S\in\mathcal S}\alpha^{\mathrm{bdry}}_S\mathbf b_S
&=\mathbf1,\\
D_{\alpha^{\mathrm{bdry}}}(t^*)
&=D_{\alpha^{\mathrm{sp}}}(t^*)=0,\\
F_\theta(\alpha^{\mathrm{bdry}},t^*)
&=F_\theta(\alpha^{\mathrm{sp}},t^*)
=J^*(\theta,p,\varepsilon).
\end{aligned}
\]
Thus \(\alpha^{\mathrm{bdry}}\in\mathcal W\), contradicting the
minimum-support choice. We conclude that
\[
|\operatorname{supp}(\alpha^{\mathrm{sp}})|\le5.
\]
% lean: Causalean.Mathlib.Optimization.exists_sparse_maximizer

\item
Nonnegative weights and positive denominators imply
\[
C_{\alpha^{\mathrm{sp}}}
=
\sum_{S\in\mathcal S}
\alpha^{\mathrm{sp}}_S
\frac{(g_{S,0}+g_{S,1})^2}{h_S(\theta)}
\ge0.
\]
Since \(D_{\alpha^{\mathrm{sp}}}(t^*)=0\), the quadratic identity yields
\[
F_\theta(\alpha^{\mathrm{sp}},t)
-F_\theta(\alpha^{\mathrm{sp}},t^*)
=
(t-t^*)^2C_{\alpha^{\mathrm{sp}}}\ge0
\qquad(t\in\mathbb R).
\]
Therefore \(t^*\) is a global minimizer, including when the quadratic
coefficient is zero, and
\[
\inf_{t\in\mathbb R}F_\theta(\alpha^{\mathrm{sp}},t)
=
F_\theta(\alpha^{\mathrm{sp}},t^*)
=
J^*(\theta,p,\varepsilon).
\]
% lean: CausalSmith.Stat.LdpAteEfficiencySurface.informationObjective_min_of_stationary

\item
Finally, set
\[
\alpha:=\alpha^{\mathrm{sp}},
\qquad
\mathcal S_{\mathrm{out}}
:=
\{S\in\mathcal S:(P_\theta Q_\alpha)(\{S\})\ne0\}.
\]
For these feasible weights, the staircase construction gives
\[
Q_\alpha(\{S\}\mid x_j)=\alpha_Sb_{S,j},
\qquad
(P_\theta Q_\alpha)(\{S\})
=
\sum_{j=0}^{3}\pi_{\theta,j}\alpha_Sb_{S,j}.
\]
In particular,
\[
\alpha_S=0
\quad\Longrightarrow\quad
Q_\alpha(\{S\}\mid x_j)=0\ \text{for every }j
\quad\Longrightarrow\quad
(P_\theta Q_\alpha)(\{S\})=0.
\]
Thus
\[
\mathcal S_{\mathrm{out}}\subseteq\operatorname{supp}(\alpha),
\qquad
\kappa(Q_\alpha)
=
|\mathcal S_{\mathrm{out}}|
\le|\operatorname{supp}(\alpha)|
\le5.
\]
Together with feasibility and the profiled equality from the preceding
step, this proves the claim.
% lean: CausalSmith.Stat.LdpAteEfficiencySurface.outputCardinality_staircase_le_support_card
\end{enumerate}
\end{proof}

\begin{proof}[Proof of \cref{obj:thm:five-output-region}]
\begin{enumerate}
\item \textit{Certificate coordinates and profiling equations.}
Let \(\mathcal K_5\) be the active-pattern set and let \(\mathcal U\) be the interior parameter domain:
\[
\mathcal K_5=\{S_3,S_4,S_6,S_8,S_9\},
\qquad
\mathcal U=(0,1)^3\times(0,\infty).
\]
For \(\xi=(p,\mu_0,\mu_1,\varepsilon)\in\mathcal U\), set
\[
\theta(\xi)=(\mu_0,\mu_1)^\top,\qquad
q=1-p,\qquad
r_\varepsilon=e^\varepsilon,\qquad d_\varepsilon=r_\varepsilon-1.
\]
Use the pattern masses, gradients, profiling directions, and projected scores from
\cref{obj:lem:staircase-refinement}. Write their dependence on \(\xi\) explicitly as
\[
h_S(\xi)=h_S(\theta(\xi)),\qquad
g_S(\xi)=\nabla_\theta h_S(\theta(\xi)),\qquad
f_S(\xi,t)=\frac{(g_S(\xi)^\top v(t))^2}{h_S(\xi)},
\qquad
\rho_S(\xi,t)=\frac{g_S(\xi)^\top v(t)}{h_S(\xi)}.
\]
Here \(p,\varepsilon\) are fixed in the gradient. In particular,
\[
v(t)=\begin{pmatrix}t\\t+1\end{pmatrix},
\qquad
e_{\mathrm{prof}}=\begin{pmatrix}1\\1\end{pmatrix},
\qquad
F_\theta(\alpha,t)=\sum_{S\in\mathcal S}\alpha_S f_S(\xi,t).
\]

Index the four records by \(x_0,x_1,x_2,x_3\). Their probabilities, in this order, are
\[
(\pi_{\theta,0},\pi_{\theta,1},\pi_{\theta,2},\pi_{\theta,3})
=
\bigl(q(1-\mu_0),q\mu_0,p(1-\mu_1),p\mu_1\bigr).
\]
The pattern enumeration is specified by
\[
x_j\in S_m
\quad\Longleftrightarrow\quad
\left\lfloor\frac{m}{2^j}\right\rfloor\equiv1\pmod 2,
\qquad
0\le j\le3,\quad 1\le m\le14.
\]
Define the ray vectors and ray matrix by
\[
(b_S(\varepsilon))_j=
\begin{cases}
r_\varepsilon,&x_j\in S,\\
1,&x_j\notin S,
\end{cases}
\qquad
B_\varepsilon=(b_{S_1}(\varepsilon),\ldots,b_{S_{14}}(\varepsilon)).
\]
All record probabilities and ray coordinates are positive, so
\[
h_S(\xi)=\sum_{j=0}^3\pi_{\theta,j}(b_S(\varepsilon))_j>0.
\]

The stationarity, dual, and score conditions described in
\cref{obj:def:support-regions} enter the following full certificate definition.
A point \(\xi\in\mathbb R^4\) belongs to \(\mathcal R_5\) if and only if
\(\xi\in\mathcal U\) and there exist
\(\alpha\in\mathbb R^{14}\), \(t\in\mathbb R\), and \(\eta\in\mathbb R^4\) satisfying
\[
\begin{gathered}
\alpha\ge0,\qquad B_\varepsilon\alpha=\mathbf1,\qquad
\alpha_S>0\ (S\in\mathcal K_5),\qquad
\alpha_S=0\ (S\notin\mathcal K_5),\\
\partial_tF_\theta(\alpha,t)=0,\\
b_S(\varepsilon)^\top\eta=f_S(\xi,t)\quad(S\in\mathcal K_5),\\
f_S(\xi,t)<b_S(\varepsilon)^\top\eta\quad(S\notin\mathcal K_5),\\
\rho_S(\xi,t)\ne\rho_T(\xi,t)
\quad(S,T\in\mathcal K_5,\ S\ne T).
\end{gathered}
\]
The derivative appearing here is
\[
\partial_tF_\theta(\alpha,t)
=
2\sum_S\alpha_S
\frac{(g_S(\xi)^\top v(t))(g_S(\xi)^\top e_{\mathrm{prof}})}
     {h_S(\xi)}.
\]

The five active columns of the ray matrix are
\[
B_\varepsilon\big|_{\mathcal K_5}
=
\begin{pmatrix}
r_\varepsilon&1&1&1&r_\varepsilon\\
r_\varepsilon&1&r_\varepsilon&1&1\\
1&r_\varepsilon&r_\varepsilon&1&1\\
1&1&1&r_\varepsilon&r_\varepsilon
\end{pmatrix},
\]
with columns ordered as \(S_3,S_4,S_6,S_8,S_9\).
Subtracting consecutive normalization equations, and using
\(r_\varepsilon-1>0\), gives
\[
\alpha_{S_3}=\alpha_{S_4}=\alpha_{S_8},\qquad
\alpha_{S_6}=\alpha_{S_9},\qquad
(r_\varepsilon+2)\alpha_{S_3}
+(r_\varepsilon+1)\alpha_{S_6}=1.
\]
Thus every feasible vector supported within \(\mathcal K_5\) has the form
\[
\alpha^u_S(\varepsilon)=
\begin{cases}
u/(r_\varepsilon+1),&S\in\{S_6,S_9\},\\
(1-u)/(r_\varepsilon+2),&S\in\{S_3,S_4,S_8\},\\
0,&S\notin\mathcal K_5,
\end{cases}
\qquad 0\le u\le1.
\]
Conversely, the displayed matrix verifies feasibility of this vector.
Its five active weights are positive precisely when \(0<u<1\).

Define the two normalized constituent objectives and their coefficients by
\[
\begin{aligned}
\Phi_{\mathrm{bin}}(\xi,t)
&=\frac{f_{S_6}(\xi,t)+f_{S_9}(\xi,t)}{r_\varepsilon+1},\\
\Phi_{\mathrm{ter}}(\xi,t)
&=\frac{f_{S_3}(\xi,t)+f_{S_4}(\xi,t)+f_{S_8}(\xi,t)}
        {r_\varepsilon+2},\\
a_{\mathrm{bin}}(\xi)
&=\frac{d_\varepsilon^2}{r_\varepsilon+1}
  \left(\frac1{h_{S_6}(\xi)}+\frac1{h_{S_9}(\xi)}\right),\\
a_{\mathrm{ter}}(\xi)
&=\frac{d_\varepsilon^2p^2}{r_\varepsilon+2}
  \left(\frac1{h_{S_4}(\xi)}+\frac1{h_{S_8}(\xi)}\right).
\end{aligned}
\]
Both coefficients are positive throughout \(\mathcal U\).
The active gradients are
\[
g_{S_3}=0,\quad
g_{S_4}=\begin{pmatrix}0\\-d_\varepsilon p\end{pmatrix},\quad
g_{S_6}=\begin{pmatrix}d_\varepsilon q\\-d_\varepsilon p\end{pmatrix},\quad
g_{S_8}=\begin{pmatrix}0\\d_\varepsilon p\end{pmatrix},\quad
g_{S_9}=-g_{S_6}.
\]
Consequently,
\[
\begin{aligned}
\Phi_{\mathrm{bin}}(\xi,t)
&=a_{\mathrm{bin}}(\xi)\{p-(1-2p)t\}^2,\\
\Phi_{\mathrm{ter}}(\xi,t)
&=a_{\mathrm{ter}}(\xi)(t+1)^2,\\
F_\theta(\alpha^u,t)
&=u\Phi_{\mathrm{bin}}(\xi,t)+(1-u)\Phi_{\mathrm{ter}}(\xi,t).
\end{aligned}
\]
The binary rays sum to \((r_\varepsilon+1)\mathbf1\), and the ternary rays
sum to \((r_\varepsilon+2)\mathbf1\). Averaging the active dual equalities
over these two families therefore yields
\[
\Phi_{\mathrm{bin}}(\xi,t)
=
\sum_{j=0}^3\eta_j
=
\Phi_{\mathrm{ter}}(\xi,t).
\]
Finally, the certificate stationarity equation becomes
\[
u\,\partial_t\Phi_{\mathrm{bin}}(\xi,t)
+(1-u)\,\partial_t\Phi_{\mathrm{ter}}(\xi,t)=0.
\]
% lean: CausalSmith.Stat.LdpAteEfficiencySurface.r5_stationarity_eq_mixture

\item \textit{Continuous continuation and openness.}
Fix
\[
\xi_\circ=(p_\circ,\mu_{0,\circ},\mu_{1,\circ},\varepsilon_\circ)
\in\mathcal R_5
\]
and a certificate \((\alpha_\circ,t_\circ,\eta_\circ)\).
Define, for \(\xi\in\mathcal U\),
\[
s_{\mathrm{loc}}(\xi)
=
\sqrt{\frac{a_{\mathrm{ter}}(\xi)}{a_{\mathrm{bin}}(\xi)}}.
\]
The constituent equality implies
\[
\{p_\circ-(1-2p_\circ)t_\circ\}^2
=
\{s_{\mathrm{loc}}(\xi_\circ)(t_\circ+1)\}^2.
\]
Splitting equality of squares gives a sign
\(\sigma\in\{-1,1\}\) such that
\[
p_\circ-(1-2p_\circ)t_\circ
=
\sigma s_{\mathrm{loc}}(\xi_\circ)(t_\circ+1).
\]
Its rearrangement is
\[
\{1-2p_\circ+\sigma s_{\mathrm{loc}}(\xi_\circ)\}t_\circ
=
p_\circ-\sigma s_{\mathrm{loc}}(\xi_\circ).
\]
If the coefficient on the left were zero, the right side would also be
zero, giving \(p_\circ=1\). Thus, on a neighborhood where its denominator
is nonzero, define
\[
t_\sigma(\xi)
=
\frac{p-\sigma s_{\mathrm{loc}}(\xi)}
     {1-2p+\sigma s_{\mathrm{loc}}(\xi)}.
\]
This function is continuous at \(\xi_\circ\), equals \(t_\circ\) there,
and satisfies
\[
\Phi_{\mathrm{bin}}(\xi,t_\sigma(\xi))
=
\Phi_{\mathrm{ter}}(\xi,t_\sigma(\xi)).
\]

At the original certificate,
\[
\partial_t\Phi_{\mathrm{ter}}(\xi_\circ,t_\circ)
=
2a_{\mathrm{ter}}(\xi_\circ)(t_\circ+1)\ne0.
\]
Indeed, vanishing would give \(t_\circ=-1\); the constituent equality
would then give
\[
a_{\mathrm{bin}}(\xi_\circ)(1-p_\circ)^2=0,
\]
contrary to \(p_\circ<1\).
If the two constituent derivatives were equal, stationarity would make
this common derivative zero. Their difference is therefore nonzero, and
stationarity uniquely determines the mixture coordinate:
\[
u_\circ=(r_{\varepsilon_\circ}+1)\alpha_{\circ,S_6}
=
-\frac{\partial_t\Phi_{\mathrm{ter}}(\xi_\circ,t_\circ)}
       {\partial_t\Phi_{\mathrm{bin}}(\xi_\circ,t_\circ)
        -\partial_t\Phi_{\mathrm{ter}}(\xi_\circ,t_\circ)}.
\]
On a neighborhood where the derivative difference remains nonzero, set
\[
u_\sigma(\xi)
=
-\frac{\partial_t\Phi_{\mathrm{ter}}(\xi,t_\sigma(\xi))}
       {\partial_t\Phi_{\mathrm{bin}}(\xi,t_\sigma(\xi))
        -\partial_t\Phi_{\mathrm{ter}}(\xi,t_\sigma(\xi))}.
\]
It is continuous at \(\xi_\circ\), equals \(u_\circ\) there, and remains in
\((0,1)\) after shrinking the neighborhood. These arguments also apply
when \(p_\circ=1/2\).

We next continue the dual certificate while preserving its value at the
original point. Define
\[
\mathsf D_\varepsilon=
\begin{pmatrix}
r_\varepsilon&r_\varepsilon&1&1\\
1&1&r_\varepsilon&1\\
1&r_\varepsilon&r_\varepsilon&1\\
1&1&1&r_\varepsilon
\end{pmatrix},
\qquad
\boldsymbol f_{\mathrm{act}}(\xi)=
\begin{pmatrix}
f_{S_3}(\xi,t_\sigma(\xi))\\
f_{S_4}(\xi,t_\sigma(\xi))\\
f_{S_6}(\xi,t_\sigma(\xi))\\
f_{S_8}(\xi,t_\sigma(\xi))
\end{pmatrix}.
\]
For \(\varepsilon>0\), an explicit right inverse is
\[
\mathsf N_\varepsilon=
\begin{pmatrix}
\dfrac{r_\varepsilon+1}{d_\varepsilon(r_\varepsilon+2)}
&\dfrac{r_\varepsilon+1}{d_\varepsilon(r_\varepsilon+2)}
&-\dfrac1{d_\varepsilon}
&-\dfrac1{d_\varepsilon(r_\varepsilon+2)}\\
0&-\dfrac1{d_\varepsilon}&\dfrac1{d_\varepsilon}&0\\
-\dfrac1{d_\varepsilon(r_\varepsilon+2)}
&\dfrac{r_\varepsilon+1}{d_\varepsilon(r_\varepsilon+2)}
&0&-\dfrac1{d_\varepsilon(r_\varepsilon+2)}\\
-\dfrac1{d_\varepsilon(r_\varepsilon+2)}
&-\dfrac1{d_\varepsilon(r_\varepsilon+2)}
&0&\dfrac{r_\varepsilon+1}{d_\varepsilon(r_\varepsilon+2)}
\end{pmatrix}.
\]
Direct multiplication gives
\[
\mathsf D_\varepsilon\mathsf N_\varepsilon=\operatorname{Id}_4.
\]
Set
\[
\mathsf N_\circ=\mathsf N_{\varepsilon_\circ}.
\]
At \(\xi_\circ\), the matrix
\(\mathsf D_{\varepsilon_\circ}\mathsf N_\circ\) is the identity.
Its determinant remains nonzero nearby. On that neighborhood define
\[
\eta(\xi)
=
\eta_\circ+
\mathsf N_\circ
(\mathsf D_\varepsilon\mathsf N_\circ)^{-1}
\bigl(\boldsymbol f_{\mathrm{act}}(\xi)
      -\mathsf D_\varepsilon\eta_\circ\bigr).
\]
Matrix inversion is continuous there. The original active equalities
make the parenthesized residual zero at \(\xi_\circ\), and multiplication
by \(\mathsf D_\varepsilon\) gives
\[
\eta(\xi_\circ)=\eta_\circ,
\qquad
\mathsf D_\varepsilon\eta(\xi)=\boldsymbol f_{\mathrm{act}}(\xi).
\]
Thus the four active equalities indexed by \(S_3,S_4,S_6,S_8\) continue.

The fifth equality follows from the constituent equality. Indeed,
\[
\frac{b_{S_6}^\top\eta+b_{S_9}^\top\eta}{r_\varepsilon+1}
=
\frac{b_{S_3}^\top\eta+b_{S_4}^\top\eta+b_{S_8}^\top\eta}
     {r_\varepsilon+2},
\]
because both sides equal \(\mathbf1^\top\eta\).
Subtracting this identity from
\(\Phi_{\mathrm{bin}}(\xi,t_\sigma)=
\Phi_{\mathrm{ter}}(\xi,t_\sigma)\), using the four continued
equalities, gives
\[
b_{S_9}(\varepsilon)^\top\eta(\xi)
=
f_{S_9}(\xi,t_\sigma(\xi)).
\]

All pattern masses stay positive nearby. Hence the finitely many
information functions, projected scores, and dual-ray values are
continuous at \(\xi_\circ\). The strict inactive slacks and nonzero score
differences therefore persist simultaneously:
\[
\begin{gathered}
f_S(\xi,t_\sigma(\xi))
<
b_S(\varepsilon)^\top\eta(\xi)
\quad(S\notin\mathcal K_5),\\
\rho_S(\xi,t_\sigma(\xi))
\ne\rho_T(\xi,t_\sigma(\xi))
\quad(S,T\in\mathcal K_5,\ S\ne T).
\end{gathered}
\]
The strict interior inequalities, the branch denominator, and the
derivative difference likewise persist. For every point in the resulting
neighborhood, choose
\[
\alpha(\xi)=\alpha^{u_\sigma(\xi)}(\varepsilon).
\]
The normalization calculation gives feasibility and five positive
weights, while the defining formula for \(u_\sigma\) gives
\[
\partial_tF_{\theta(\xi)}
\bigl(\alpha(\xi),t_\sigma(\xi)\bigr)=0.
\]
Together with the continued dual equalities, strict slacks, and score
separation, these are all the certificate conditions. Every point of
\(\mathcal R_5\) consequently has a neighborhood in \(\mathbb R^4\)
contained in \(\mathcal R_5\), proving openness.
% lean: CausalSmith.Stat.LdpAteEfficiencySurface.isOpen_R5_continuousCertificate

\item \textit{The centered certificate and nonemptiness.}
Fix \(0<p<1/2\), and define
\[
\theta_{\mathrm{cen}}=\begin{pmatrix}1/2\\1/2\end{pmatrix},
\qquad
\xi_{\mathrm{cen}}(p)=(p,1/2,1/2,\log3),
\]
\[
\begin{aligned}
s_{\mathrm{cen}}&=\sqrt{\frac{8p^2}{5(1+p)}},
&
D_{\mathrm{cen}}&=1-2p+s_{\mathrm{cen}},\\
t_{\mathrm{cen}}&=\frac{p-s_{\mathrm{cen}}}{D_{\mathrm{cen}}},
&
\lambda_{\mathrm{cen}}&=\frac{s_{\mathrm{cen}}}{D_{\mathrm{cen}}},\\
J_{\mathrm{cen}}&=\frac{s_{\mathrm{cen}}^2(1-p)^2}{D_{\mathrm{cen}}^2},
&
\eta_{\mathrm{cen}}&=\frac{J_{\mathrm{cen}}}{4}(-1,-1,3,3)^\top.
\end{aligned}
\]
The identity
\[
5(1+p)s_{\mathrm{cen}}^2=8p^2
\]
and \(0<p<1/2\) give
\[
p<s_{\mathrm{cen}}<2p,\qquad
D_{\mathrm{cen}}>0,\qquad
0<\lambda_{\mathrm{cen}}<1,\qquad
t_{\mathrm{cen}}<0.
\]
Set
\[
\alpha_{\mathrm{cen}}=\alpha^{\lambda_{\mathrm{cen}}}(\log3).
\]
Thus
\[
\alpha_{\mathrm{cen},S_6}
=\alpha_{\mathrm{cen},S_9}=\frac{\lambda_{\mathrm{cen}}}{4},
\qquad
\alpha_{\mathrm{cen},S_3}
=\alpha_{\mathrm{cen},S_4}
=\alpha_{\mathrm{cen},S_8}
=\frac{1-\lambda_{\mathrm{cen}}}{5},
\]
and all other weights vanish. Its row sums are
\[
5\,\frac{1-\lambda_{\mathrm{cen}}}{5}
+
4\,\frac{\lambda_{\mathrm{cen}}}{4}=1,
\]
so it is feasible and has precisely the five stated positive weights.

At these parameters the constituent objectives reduce to
\[
\Phi_{\mathrm{bin}}(\xi_{\mathrm{cen}}(p),t)
=\{p-(1-2p)t\}^2,
\qquad
\Phi_{\mathrm{ter}}(\xi_{\mathrm{cen}}(p),t)
=s_{\mathrm{cen}}^2(t+1)^2.
\]
Define the positive score magnitudes
\[
\rho_{\mathrm{bin}}
=p-(1-2p)t_{\mathrm{cen}}
=\frac{s_{\mathrm{cen}}q}{D_{\mathrm{cen}}},
\qquad
\rho_{\mathrm{ter}}
=\frac{2p(t_{\mathrm{cen}}+1)}{1+p}.
\]
Direct rearrangement gives
\[
t_{\mathrm{cen}}+1=\frac{q}{D_{\mathrm{cen}}},
\qquad
s_{\mathrm{cen}}(t_{\mathrm{cen}}+1)=\rho_{\mathrm{bin}},
\qquad
J_{\mathrm{cen}}=\rho_{\mathrm{bin}}^2>0.
\]
Consequently,
\[
\begin{aligned}
\partial_tF_{\theta_{\mathrm{cen}}}
(\alpha_{\mathrm{cen}},t_{\mathrm{cen}})
&=-2\lambda_{\mathrm{cen}}(1-2p)\rho_{\mathrm{bin}}
  +2(1-\lambda_{\mathrm{cen}})s_{\mathrm{cen}}\rho_{\mathrm{bin}}\\
&=0.
\end{aligned}
\]

For completeness, the information and dual calculations over all
fourteen patterns are
\[
\begin{array}{c|c|c}
\text{patterns}
&b_S(\log3)^\top\eta_{\mathrm{cen}}/J_{\mathrm{cen}}
&f_S(\xi_{\mathrm{cen}}(p),t_{\mathrm{cen}})/J_{\mathrm{cen}}\\ \hline
S_1,S_2&1/2&
\dfrac{4q^2t_{\mathrm{cen}}^2}{(2-p)J_{\mathrm{cen}}}\\
S_3&0&0\\
S_4,S_8&5/2&5/2\\
S_5,S_{10}&2&
\dfrac{2(p+t_{\mathrm{cen}})^2}{J_{\mathrm{cen}}}\\
S_6,S_9&2&2\\
S_7,S_{11}&3/2&
\dfrac{5(1+p)}{2(3-p)}\\
S_{12}&4&0\\
S_{13},S_{14}&7/2&
\dfrac{4q^2t_{\mathrm{cen}}^2}{(2+p)J_{\mathrm{cen}}}
\end{array}
\]
This table follows by inserting the displayed record probabilities and
binary-mask enumeration into \(h_S\), \(g_S\), and \(b_S^\top\eta_{\mathrm{cen}}\).
For example, for \(S_4=\{x_2\}\),
\[
h_{S_4}(\xi_{\mathrm{cen}}(p))=1+p,\qquad
g_{S_4}(\xi_{\mathrm{cen}}(p))=(0,-2p)^\top,
\]
so
\[
f_{S_4}(\xi_{\mathrm{cen}}(p),t_{\mathrm{cen}})
=
\frac{4p^2(t_{\mathrm{cen}}+1)^2}{1+p}
=\frac52J_{\mathrm{cen}},
\]
and its ray gives the same dual value.

We verify strictness in the five inactive families. First,
\[
5s_{\mathrm{cen}}^2<8p^2
\quad\Longrightarrow\quad
s_{\mathrm{cen}}<\frac{4p}{3}.
\]
Together with \(p<s_{\mathrm{cen}}\), this yields
\[
8(p-s_{\mathrm{cen}})^2
<
\frac89p^2
<
\frac32p^2
<
(2-p)s_{\mathrm{cen}}^2.
\]
Multiplying by \(q^2/D_{\mathrm{cen}}^2>0\) and dividing by \(2(2-p)\)
proves
\[
\frac{4q^2t_{\mathrm{cen}}^2}{2-p}<\frac12J_{\mathrm{cen}},
\]
the strict slack for \(S_1,S_2\).

Next,
\[
p+t_{\mathrm{cen}}
=\frac{q(2p-s_{\mathrm{cen}})}{D_{\mathrm{cen}}},
\qquad
0<p+t_{\mathrm{cen}}<\frac{qs_{\mathrm{cen}}}{D_{\mathrm{cen}}}
=\rho_{\mathrm{bin}}.
\]
Squaring proves
\[
2(p+t_{\mathrm{cen}})^2<2J_{\mathrm{cen}},
\]
which handles \(S_5,S_{10}\).

For \(S_7,S_{11}\), the same square identity gives
\[
3s_{\mathrm{cen}}^2(3-p)-8p^2
=
\frac{8p^2(4-8p)}{5(1+p)}>0.
\]
Using \(t_{\mathrm{cen}}+1=q/D_{\mathrm{cen}}\), this is exactly
\[
\frac{4p^2(t_{\mathrm{cen}}+1)^2}{3-p}
<
\frac32J_{\mathrm{cen}}.
\]
For \(S_{12}\), strictness is \(0<4J_{\mathrm{cen}}\).
Finally,
\[
8(p-s_{\mathrm{cen}})^2
<
(2-p)s_{\mathrm{cen}}^2
<
7(2+p)s_{\mathrm{cen}}^2
\]
gives
\[
\frac{4q^2t_{\mathrm{cen}}^2}{2+p}
<
\frac72J_{\mathrm{cen}},
\]
as required for \(S_{13},S_{14}\).

The active projected scores, in the order \(S_3,S_4,S_6,S_8,S_9\), are
\[
\bigl(0,-\rho_{\mathrm{ter}},-\rho_{\mathrm{bin}},
          \rho_{\mathrm{ter}},\rho_{\mathrm{bin}}\bigr).
\]
Both magnitudes are positive, and
\[
\left(\frac{\rho_{\mathrm{ter}}}{\rho_{\mathrm{bin}}}\right)^2
=
\frac{5}{2(1+p)}>1.
\]
Hence all five scores are distinct: the comparisons between positive
magnitudes are strict, and the remaining comparisons follow from their
signs. Reversing the order of a pair gives the reverse comparison.

We have verified feasibility, positive active weights, stationarity,
all active dual equalities, all strict inactive slacks, and score
separation. Therefore
\[
\xi_{\mathrm{cen}}(p)\in\mathcal R_5
\qquad(0<p<1/2).
\]
Taking \(p=1/4\) proves that \(\mathcal R_5\) is nonempty.
% lean: CausalSmith.Stat.LdpAteEfficiencySurface.R5_nonempty_centered

\item \textit{Centered optimality and uniqueness.}
Stationarity and the nonnegative quadratic coefficients give
\[
\begin{aligned}
&F_{\theta_{\mathrm{cen}}}(\alpha_{\mathrm{cen}},t)
-F_{\theta_{\mathrm{cen}}}(\alpha_{\mathrm{cen}},t_{\mathrm{cen}})\\
&\qquad=
\bigl\{\lambda_{\mathrm{cen}}(1-2p)^2
 +(1-\lambda_{\mathrm{cen}})s_{\mathrm{cen}}^2\bigr\}
(t-t_{\mathrm{cen}})^2\ge0.
\end{aligned}
\]
Its value at \(t_{\mathrm{cen}}\) is \(J_{\mathrm{cen}}\).
For every feasible \(\beta\), nonnegativity of its weights and the dual
bounds give
\[
\begin{aligned}
\inf_{t\in\mathbb R}F_{\theta_{\mathrm{cen}}}(\beta,t)
&\le F_{\theta_{\mathrm{cen}}}(\beta,t_{\mathrm{cen}})\\
&\le\sum_S\beta_S b_S(\log3)^\top\eta_{\mathrm{cen}}\\
&=(B_{\log3}\beta)^\top\eta_{\mathrm{cen}}
=\mathbf1^\top\eta_{\mathrm{cen}}
=J_{\mathrm{cen}}.
\end{aligned}
\]
The infima here are bounded below by zero and are taken over nonempty
sets. The feasible vector \(\alpha_{\mathrm{cen}}\) reaches the upper
bound, so the optimization in \cref{obj:def:staircase-saddle} yields
\[
J^*(\theta_{\mathrm{cen}},p,\log3)
=
J_{\mathrm{cen}}
=
\inf_tF_{\theta_{\mathrm{cen}}}(\alpha_{\mathrm{cen}},t).
\]

Suppose a feasible \(\beta\) has this optimized profile. Every inequality
in the preceding bound is then an equality. In particular,
\[
\sum_S\beta_S
\bigl\{b_S(\log3)^\top\eta_{\mathrm{cen}}
       -f_S(\xi_{\mathrm{cen}}(p),t_{\mathrm{cen}})\bigr\}=0.
\]
Each summand is nonnegative, and every inactive slack is strictly
positive. Thus
\[
\beta_S=0\qquad(S\notin\mathcal K_5).
\]
Also,
\[
F_{\theta_{\mathrm{cen}}}(\beta,t_{\mathrm{cen}})
=
J_{\mathrm{cen}}
\le F_{\theta_{\mathrm{cen}}}(\beta,t)
\quad(t\in\mathbb R),
\]
so differentiation of this quadratic at its minimum gives
\[
\partial_tF_{\theta_{\mathrm{cen}}}(\beta,t_{\mathrm{cen}})=0.
\]

The row constraints already show that
\[
\beta_{S_3}=\beta_{S_4}=\beta_{S_8}
=\frac{1-4\beta_{S_6}}5,
\qquad
\beta_{S_9}=\beta_{S_6}.
\]
To solve the remaining stationarity equation, define its reduced
coefficient
\[
C_{\mathrm{cen}}
=
10s_{\mathrm{cen}}p^2+5s_{\mathrm{cen}}p
-5s_{\mathrm{cen}}-8p^2.
\]
Substitute these row coordinates into the displayed derivative formula,
use \(t_{\mathrm{cen}}=(p-s_{\mathrm{cen}})/D_{\mathrm{cen}}\), and clear
the positive denominators. The resulting equation is
\[
(p-1)\bigl(C_{\mathrm{cen}}\beta_{S_6}+2p^2\bigr)=0.
\]
Since \(p<1/2\), it follows that
\[
C_{\mathrm{cen}}\beta_{S_6}=-2p^2.
\]
The square identity for \(s_{\mathrm{cen}}\) further gives
\[
\begin{aligned}
C_{\mathrm{cen}}s_{\mathrm{cen}}
&=5(2p-1)(1+p)s_{\mathrm{cen}}^2-8p^2s_{\mathrm{cen}}\\
&=-8p^2D_{\mathrm{cen}}<0.
\end{aligned}
\]
Thus \(C_{\mathrm{cen}}\ne0\), and
\[
C_{\mathrm{cen}}\bigl(4D_{\mathrm{cen}}\beta_{S_6}\bigr)
=
-8p^2D_{\mathrm{cen}}
=
C_{\mathrm{cen}}s_{\mathrm{cen}}.
\]
Cancellation proves
\[
\beta_{S_6}=\frac{s_{\mathrm{cen}}}{4D_{\mathrm{cen}}}
=\alpha_{\mathrm{cen},S_6}.
\]
The row-coordinate identities determine the other four active weights,
and both vectors vanish elsewhere. Hence
\[
\beta=\alpha_{\mathrm{cen}},
\]
proving the centered uniqueness assertion.
% lean: CausalSmith.Stat.LdpAteEfficiencySurface.centeredR5_profile_optimal_unique

\item \textit{Attainment, refinement, and the centered cardinality bound.}
We first give the directional inequality used for attaining channels.
For every \(\xi\in\mathcal U\), define the singleton design
\[
\alpha^{\mathrm{sgl}}_S=
\begin{cases}
1/(r_\varepsilon+3),&S\in\{S_1,S_2,S_4,S_8\},\\
0,&\text{otherwise}.
\end{cases}
\]
These four patterns are the four singleton subsets. Their rays sum to
\((r_\varepsilon+3)\mathbf1\), proving feasibility. Its objective is a
quadratic with nonnegative terms, and in particular
\[
F_\theta(\alpha^{\mathrm{sgl}},t)
\ge
\frac{d_\varepsilon^2q^2}
     {(r_\varepsilon+3)h_{S_1}(\xi)}t^2,
\qquad
F_\theta(\alpha^{\mathrm{sgl}},t)
\ge
\frac{d_\varepsilon^2p^2}
     {(r_\varepsilon+3)h_{S_8}(\xi)}(t+1)^2.
\]
Both coefficients are positive. This quadratic has a minimizer, and a
zero minimum would force both \(t=0\) and \(t+1=0\). Therefore
\[
0<\inf_tF_\theta(\alpha^{\mathrm{sgl}},t)
\le J^*(\theta,p,\varepsilon),
\]
where the final comparison is supplied by the optimization in
\cref{obj:def:staircase-saddle}.

For a stationary private channel \(Q\), write
\[
\mathsf P_Q=P_\theta Q,\qquad
\nu_Q=\sum_{j=0}^3Q(x_j,\cdot),
\]
and define its output score by
\[
s_Q(z)=
\begin{cases}
\displaystyle
\frac{\nabla_\theta\{d\mathsf P_Q/d\nu_Q\}(z)}
     {\{d\mathsf P_Q/d\nu_Q\}(z)},
&\{d\mathsf P_Q/d\nu_Q\}(z)>0,\\[6pt]
0,&\{d\mathsf P_Q/d\nu_Q\}(z)=0.
\end{cases}
\]
The score representation of the information matrix in
\cref{obj:lem:staircase-refinement} gives, with
\(I_Q=I_\theta(Q)\),
\[
I_Q=\int s_Q(z)s_Q(z)^\top\,d\mathsf P_Q(z),
\qquad
v_{\mathrm{sol}}^\top I_Qv_{\mathrm{sol}}
=
\int(v_{\mathrm{sol}}^\top s_Q(z))^2\,d\mathsf P_Q(z)\ge0
\quad(v_{\mathrm{sol}}\in\mathbb R^2).
\]
Thus \(I_Q\) is positive semidefinite.

Let \(\mathcal V(I)\) denote the extended contrast variance:
\[
\mathcal V(I)=
\begin{cases}
\displaystyle
\max\!\left\{\inf_{\{v_{\mathrm{sol}}\in\mathbb R^2:
                         Iv_{\mathrm{sol}}=c\}}
                  c^\top v_{\mathrm{sol}},\,0\right\},
&c\in\operatorname{range}(I),\\
+\infty,&c\notin\operatorname{range}(I),
\end{cases}
\qquad
c=\begin{pmatrix}-1\\1\end{pmatrix}.
\]
If \(Q\) attains the variance in the statement, its variance is finite,
so there is a vector \(v_{\mathrm{sol}}\in\mathbb R^2\) satisfying
\[
I_Qv_{\mathrm{sol}}=c.
\]
Define
\[
V_{\mathrm{sol}}=c^\top v_{\mathrm{sol}}
=v_{\mathrm{sol}}^\top I_Qv_{\mathrm{sol}}>0.
\]
To justify strict positivity, zero quadratic value for a positive
semidefinite matrix implies \(I_Qv_{\mathrm{sol}}=0\), contradicting
\(c\ne0\). Any other solution has the same contrast pairing, since
\[
I_Qv_{\mathrm{sol}}'=c
\quad\Longrightarrow\quad
c^\top v_{\mathrm{sol}}'
=
v_{\mathrm{sol}}^\top I_Qv_{\mathrm{sol}}'
=
v_{\mathrm{sol}}^\top c
=
V_{\mathrm{sol}}.
\]
Consequently attainment and positivity of \(J^*\) imply
\[
V_{\mathrm{sol}}=\mathcal V(I_Q)
=V^*(\theta,p,\varepsilon)
=\frac1{J^*(\theta,p,\varepsilon)}.
\]
Here both values are positive, so the nonnegative-part convention
preserves their real values.

For arbitrary \(t\in\mathbb R\), define the residual vector
\[
v_{\mathrm{res}}=V_{\mathrm{sol}}v(t)-v_{\mathrm{sol}}.
\]
Since \(c^\top v(t)=1\), positive semidefiniteness gives
\[
\begin{aligned}
0
&\le v_{\mathrm{res}}^\top I_Qv_{\mathrm{res}}\\
&=V_{\mathrm{sol}}^2v(t)^\top I_Qv(t)-V_{\mathrm{sol}}.
\end{aligned}
\]
Dividing by \(V_{\mathrm{sol}}>0\) and substituting its value proves
\[
J^*(\theta,p,\varepsilon)\le v(t)^\top I_Qv(t)
\qquad(t\in\mathbb R).
\]
% lean: CausalSmith.Stat.LdpAteEfficiencySurface.attainment_forces_direction_lower

Return to the centered parameters and an attaining \(Q\). Use the
constant privacy sequence \(\varepsilon_n=\log3\); it satisfies
\cref{obj:ass:fixed-privacy}. The centered parameters also satisfy
\cref{obj:ass:interior-assignment,obj:ass:interior-means}.
By \cref{obj:lem:staircase-refinement}, there are feasible weights
\(\beta\) and a probability kernel \(K\) such that
\[
Q=K\circ Q_\beta,
\qquad
I_{\theta_{\mathrm{cen}}}(\beta)-I_{\theta_{\mathrm{cen}}}(Q)\succeq0.
\]
For every \(t\), the directional lower bound and this matrix order imply
\[
J_{\mathrm{cen}}
\le
v(t)^\top I_{\theta_{\mathrm{cen}}}(Q)v(t)
\le
v(t)^\top I_{\theta_{\mathrm{cen}}}(\beta)v(t)
=
F_{\theta_{\mathrm{cen}}}(\beta,t).
\]
The centered dual bound supplies the reverse inequality at
\(t_{\mathrm{cen}}\). Hence
\[
F_{\theta_{\mathrm{cen}}}(\beta,t_{\mathrm{cen}})
=J_{\mathrm{cen}}.
\]
Equality in the strict dual bound forces \(\beta\) to vanish outside
\(\mathcal K_5\). The preceding inequality for every \(t\) makes
\(t_{\mathrm{cen}}\) a minimizer for \(\beta\), so its derivative there is
zero. The supported stationary calculation above therefore gives
\[
\beta=\alpha_{\mathrm{cen}},
\qquad
\beta_S>0\quad(S\in\mathcal K_5).
\]
Moreover, the directional inequalities at \(t_{\mathrm{cen}}\) give
\[
v(t_{\mathrm{cen}})^\top I_{\theta_{\mathrm{cen}}}(\beta)
v(t_{\mathrm{cen}})
=
v(t_{\mathrm{cen}})^\top I_{\theta_{\mathrm{cen}}}(Q)
v(t_{\mathrm{cen}})
=
J_{\mathrm{cen}}.
\]
The score-equality conclusion of
\cref{obj:lem:staircase-refinement} now applies. If two distinct active
conditional output measures were not mutually singular, their projected
scores would agree. The centered score calculation proves that they
differ. Thus, writing \(K_S=K(S,\cdot)\),
\[
K_S\perp K_T
\qquad(S,T\in\mathcal K_5,\ S\ne T).
\]

We spell out how these five singular probability measures give the
cardinality bound. Evaluations of output-space measures on arbitrary
subsets use their outer-measure extensions. For any probability measure
\(\mathsf P_{\mathrm{out}}\) on \(Z\), define
\[
\mathsf P_{\mathrm{out}}^{\mathrm{out}}(A)
=
\inf\bigl\{\mathsf P_{\mathrm{out}}(E):
A\subseteq E\subseteq Z,\ E\text{ measurable}\bigr\},
\qquad A\subseteq Z.
\]
We also write \(\mathsf P_{\mathrm{out}}(A)\) for this extension; it agrees
with the original measure on measurable sets. In particular, singleton
evaluations of \(K_S\) and \(\mathsf P_Q\) below use this convention, as
does the output cardinality:
\[
\kappa(Q)=
\begin{cases}
\#\{z\in Z:\mathsf P_Q^{\mathrm{out}}(\{z\})\ne0\},&Z\text{ finite},\\
+\infty,&Z\text{ infinite}.
\end{cases}
\]
These extensions are monotone and finitely subadditive. Absolute
continuity transfers to their evaluations on arbitrary sets: a set of
outer measure zero is contained in a measurable null set, obtained by
intersecting measurable supersets whose masses tend to zero; the
dominated measure also vanishes on that measurable null set and hence
on the original set.

For every measurable output event \(A\), factorization
and the positive record probability \(\pi_{\theta,0}>0\) give
\[
\beta_S(b_S(\varepsilon))_0K_S(A)\le Q(x_0,A),
\qquad
\pi_{\theta,0}Q(x_0,A)\le\mathsf P_Q(A).
\]
The coefficients on the left are positive for active \(S\). Therefore
\[
K_S\ll Q(x_0,\cdot)\ll\mathsf P_Q
\qquad(S\in\mathcal K_5).
\]
Suppose first that the output space \(Z\) is finite. For each
\(S\in\mathcal K_5\), choose a point \(z_S\in Z\) such that
\[
K_S(\{z_S\})>0.
\]
Such a point exists because otherwise finite subadditivity would give
\[
1=K_S(Z)
=K_S\!\left(\bigcup_{z\in Z}\{z\}\right)
\le\sum_{z\in Z}K_S(\{z\})=0.
\]
This contradiction also covers an empty output space, with the sum
interpreted as zero. Absolute continuity then gives
\[
\mathsf P_Q(\{z_S\})>0\qquad(S\in\mathcal K_5),
\]
since a zero output-law mass would force \(K_S(\{z_S\})=0\).

To prove that the chosen points are distinct, fix distinct
\(S,T\in\mathcal K_5\). Mutual singularity supplies a measurable set
\(E_{S,T}\subseteq Z\) with
\[
K_S(E_{S,T})=0,\qquad K_T(Z\setminus E_{S,T})=0.
\]
Suppose \(z_S=z_T\). If \(z_S\in E_{S,T}\), measure monotonicity gives
\[
0<K_S(\{z_S\})\le K_S(E_{S,T})=0.
\]
If \(z_S\notin E_{S,T}\), then \(z_T\in Z\setminus E_{S,T}\), and
\[
0<K_T(\{z_T\})\le K_T(Z\setminus E_{S,T})=0.
\]
Both cases are contradictory. Thus the five selected points are distinct
and have positive output-law mass, proving
\[
5=|\mathcal K_5|
=\#\{z_S:S\in\mathcal K_5\}
\le
\#\{z\in Z:\mathsf P_Q(\{z\})\ne0\}
=\kappa(Q).
\]
For an infinite output space, the stated convention gives
\(\kappa(Q)=+\infty\), and the same bound follows.
% lean: CausalSmith.Stat.LdpAteEfficiencySurface.centeredR5_outputCardinality_ge_five

\item \textit{Exchange of the treatment labels.}
Suppose \(1/2<p<1\) and put \(q=1-p\). Then
\[
0<q<1/2,
\qquad
(q,1/2,1/2,\log3)\in\mathcal R_5.
\]
Define the input exchange, its action on patterns and weights, and the
coordinate-exchange matrix by
\[
\begin{gathered}
\mathfrak s(x_0)=x_2,\quad
\mathfrak s(x_1)=x_3,\quad
\mathfrak s(x_2)=x_0,\quad
\mathfrak s(x_3)=x_1,\\
\mathfrak s S=\{\mathfrak s x:x\in S\},\qquad
\alpha^{\mathrm{sw}}_S=\alpha_{\mathfrak s S},\qquad
\Pi=\begin{pmatrix}0&1\\1&0\end{pmatrix}.
\end{gathered}
\]
Both exchanges are involutions. The ray identity
\[
(b_{\mathfrak s S}(\varepsilon))_j
=b_S(\varepsilon)_{\mathfrak s(j)}
\]
shows, by reindexing the normalization sums, that
\[
\alpha\in\mathcal A_\varepsilon
\quad\Longleftrightarrow\quad
\alpha^{\mathrm{sw}}\in\mathcal A_\varepsilon.
\]
On the centered slice, exchanging the record probabilities gives
\[
h_{\mathfrak s S}(\xi_{\mathrm{cen}}(p))
=h_S(\xi_{\mathrm{cen}}(q)),
\qquad
g_{\mathfrak s S}(\xi_{\mathrm{cen}}(p))
=\Pi g_S(\xi_{\mathrm{cen}}(q)).
\]
Since
\[
\Pi v(t)=-v(-1-t),
\]
the information contributions satisfy
\[
f_{\mathfrak s S}(\xi_{\mathrm{cen}}(p),t)
=
f_S(\xi_{\mathrm{cen}}(q),-1-t).
\]
Reindexing the sum and using the bijection \(t\mapsto-1-t\) consequently
gives
\[
\inf_t\sum_S\alpha^{\mathrm{sw}}_S
 f_S(\xi_{\mathrm{cen}}(p),t)
=
\inf_t\sum_S\alpha_S f_S(\xi_{\mathrm{cen}}(q),t).
\]
The exchange is also a bijection of feasible weights. Taking the
optimized profiles on both sides proves
\[
J^*(\theta_{\mathrm{cen}},p,\log3)
=
J^*(\theta_{\mathrm{cen}},q,\log3),
\qquad
V^*(\theta_{\mathrm{cen}},p,\log3)
=
V^*(\theta_{\mathrm{cen}},q,\log3).
\]

Choose the centered optimizer at \(q\) and exchange its weights. This
produces a feasible optimizer at \(p\) with support
\[
\mathfrak s\mathcal K_5=\{\mathfrak s S:S\in\mathcal K_5\}.
\]
If another feasible vector \(\beta\) is optimal at \(p\), then
\(\beta^{\mathrm{sw}}\) is optimal at \(q\). Centered uniqueness gives
\[
\beta^{\mathrm{sw}}=\alpha_{\mathrm{cen}}(q).
\]
Applying the involution again proves
\[
\beta=\alpha_{\mathrm{cen}}(q)^{\mathrm{sw}},
\]
the claimed uniqueness at the original assignment probability.

For the channel assertion, define
\[
Q^{\mathrm{sw}}(x,A)=Q(\mathfrak s x,A).
\]
This is a probability kernel and satisfies the same privacy inequalities
as \(Q\), since those inequalities hold for every input pair.
In the next identities, quantities for \(Q^{\mathrm{sw}}\) are evaluated
at assignment probability \(q\), and those for \(Q\) at \(p\).
Reindexing the four record probabilities gives equality of the output
laws; reindexing their derivatives exchanges the two parameter
coordinates. Therefore
\[
P_{\theta_{\mathrm{cen}}}Q^{\mathrm{sw}}
=
P_{\theta_{\mathrm{cen}}}Q,
\qquad
I_{\theta_{\mathrm{cen}}}(Q^{\mathrm{sw}})
=
\Pi I_{\theta_{\mathrm{cen}}}(Q)\Pi.
\]
The contrast variance is preserved by this matrix exchange. Indeed,
\(\Pi c=-c\), and for every \(v_{\mathrm{sol}}\in\mathbb R^2\),
\[
Iv_{\mathrm{sol}}=c
\quad\Longleftrightarrow\quad
(\Pi I\Pi)(-\Pi v_{\mathrm{sol}})=c,
\qquad
c^\top(-\Pi v_{\mathrm{sol}})=c^\top v_{\mathrm{sol}}.
\]
Thus the solution-value sets coincide; if they are empty, both contrast
variances equal \(+\infty\). In either case,
\[
\mathcal V(\Pi I\Pi)=\mathcal V(I).
\]
An attaining \(Q\) at \(p\) consequently gives an attaining
\(Q^{\mathrm{sw}}\) at \(q\). The centered bound yields
\(\kappa(Q^{\mathrm{sw}})\ge5\). Equality of their output laws, on the same
output space, gives
\[
\kappa(Q^{\mathrm{sw}})=\kappa(Q),
\]
and hence the desired bound at \(p\).
% lean: CausalSmith.Stat.LdpAteEfficiencySurface.reflected_five_output_region

\item \textit{Uniqueness for a general five-pattern certificate.}
Fix \(\xi=(p,\mu_0,\mu_1,\varepsilon)\in\mathcal R_5\), put
\(\theta=\theta(\xi)\), and choose a certificate \((\alpha_0,t_0,\eta)\).
For any feasible weight vector, its objective is nonnegative.
For the certificate vector, expansion of the quadratic gives
\[
\begin{aligned}
F_\theta(\alpha_0,t)-F_\theta(\alpha_0,t_0)
&=(t-t_0)\partial_tF_\theta(\alpha_0,t_0)\\
&\quad+(t-t_0)^2
 \sum_S\alpha_{0,S}
 \frac{(g_S(\xi)^\top e_{\mathrm{prof}})^2}{h_S(\xi)}\\
&\ge0,
\end{aligned}
\]
where stationarity makes the linear term zero.
The active equalities, the vanishing inactive weights, and normalization
give
\[
F_\theta(\alpha_0,t_0)
=
\sum_S\alpha_{0,S}b_S(\varepsilon)^\top\eta
=
\sum_{j=0}^3\eta_j.
\]
For every feasible \(\beta\),
\[
\inf_tF_\theta(\beta,t)
\le F_\theta(\beta,t_0)
\le\sum_S\beta_Sb_S(\varepsilon)^\top\eta
=\sum_{j=0}^3\eta_j.
\]
The certificate vector attains this upper bound. Hence
\[
J^*(\theta,p,\varepsilon)
=
\inf_tF_\theta(\alpha_0,t)
=
F_\theta(\alpha_0,t_0)
=
\sum_{j=0}^3\eta_j.
\]
Moreover, equality of another feasible profile with \(J^*\) forces
equality in the dual bound and therefore
\[
\beta_S=0\qquad(S\notin\mathcal K_5),
\]
by the same nonnegative slack-sum argument.

The constituent coefficients are positive, and the active dual
equalities give
\[
a_{\mathrm{bin}}(\xi)\{p-(1-2p)t_0\}^2
=
a_{\mathrm{ter}}(\xi)(t_0+1)^2.
\]
As already derived, \(t_0=-1\) would force \(p=1\), so
\[
\partial_t\Phi_{\mathrm{ter}}(\xi,t_0)\ne0.
\]
Define the certificate mixture coordinate by
\[
u_0=(r_\varepsilon+1)\alpha_{0,S_6}.
\]
The normalization and stationarity calculations give
\[
u_0
=
-\frac{\partial_t\Phi_{\mathrm{ter}}(\xi,t_0)}
       {\partial_t\Phi_{\mathrm{bin}}(\xi,t_0)
        -\partial_t\Phi_{\mathrm{ter}}(\xi,t_0)},
\qquad
\partial_t\Phi_{\mathrm{bin}}(\xi,t_0)
-\partial_t\Phi_{\mathrm{ter}}(\xi,t_0)\ne0.
\]

Now let \(\beta\) be any feasible optimized vector. Its nonnegative
objective has profile \(J^*\), and the dual upper bound at \(t_0\) gives
\[
J^*\le F_\theta(\beta,t_0)\le J^*.
\]
It follows that
\[
F_\theta(\beta,t_0)=J^*
\le F_\theta(\beta,t)\quad(t\in\mathbb R),
\qquad
\partial_tF_\theta(\beta,t_0)=0.
\]
Since \(\beta\) vanishes outside \(\mathcal K_5\), define
\[
u_\beta=(r_\varepsilon+1)\beta_{S_6}.
\]
Its stationarity equation yields
\[
u_\beta
=
-\frac{\partial_t\Phi_{\mathrm{ter}}(\xi,t_0)}
       {\partial_t\Phi_{\mathrm{bin}}(\xi,t_0)
        -\partial_t\Phi_{\mathrm{ter}}(\xi,t_0)}
=u_0.
\]
The supported normalization formula determines every weight from this
coordinate. Hence \(\beta=\alpha_0\). In particular, the unique optimized
vector has positive weight on all five active patterns.
% lean: CausalSmith.Stat.LdpAteEfficiencySurface.r5Certificate_unique_optimum

\item \textit{The cardinality bound throughout the region.}
Keep the general certificate from the preceding argument and let \(Q\)
be any stationary \(\varepsilon\)-LDP channel attaining the specified
extended contrast variance. The certificate includes the interior
conditions and \(\varepsilon>0\). With the constant sequence
\(\varepsilon_n=\varepsilon\), all hypotheses of
\cref{obj:lem:staircase-refinement} hold. It supplies feasible weights
\(\beta\) and a probability kernel \(K\) with
\[
Q=K\circ Q_\beta,
\qquad
I_\theta(\beta)-I_\theta(Q)\succeq0.
\]
The directional attainment bound already derived applies to these
parameters. Consequently, for every \(t\in\mathbb R\),
\[
J^*(\theta,p,\varepsilon)
\le v(t)^\top I_\theta(Q)v(t)
\le v(t)^\top I_\theta(\beta)v(t)
=F_\theta(\beta,t).
\]
At the certificate direction, the dual bound is
\[
F_\theta(\beta,t_0)\le
\sum_{j=0}^3\eta_j
=
J^*(\theta,p,\varepsilon).
\]
Thus
\[
F_\theta(\beta,t_0)=J^*(\theta,p,\varepsilon)
\le F_\theta(\beta,t)\quad(t\in\mathbb R),
\]
and
\[
\inf_tF_\theta(\beta,t)=J^*(\theta,p,\varepsilon).
\]
The general certificate uniqueness just proved gives
\[
\beta=\alpha_0,\qquad \beta_S>0\quad(S\in\mathcal K_5).
\]
The same upper and lower bounds at \(t_0\) also give directional equality:
\[
v(t_0)^\top I_\theta(\beta)v(t_0)
=
v(t_0)^\top I_\theta(Q)v(t_0)
=
J^*(\theta,p,\varepsilon).
\]
By the score-equality conclusion of
\cref{obj:lem:staircase-refinement}, overlapping active conditional
measures would have equal projected scores. The region's certificate
separates every pair, so
\[
K(S,\cdot)\perp K(T,\cdot)
\qquad(S,T\in\mathcal K_5,\ S\ne T).
\]
Factorization, positive active weights, positive ray coordinates, and
\(\pi_{\theta,0}=q(1-\mu_0)>0\) give
\[
K(S,\cdot)\ll Q(x_0,\cdot)\ll P_\theta Q
\qquad(S\in\mathcal K_5).
\]
Suppose first that the output space \(Z\) is finite. With the outer-measure
convention above, the conditional-atom argument chooses, for every
\(S\in\mathcal K_5\), a point
\(z_S\in Z\) satisfying
\[
K(S,\{z_S\})>0.
\]
Indeed, if all singleton masses of this probability measure were zero,
finite subadditivity would give
\[
1=K(S,Z)\le\sum_{z\in Z}K(S,\{z\})=0,
\]
including the empty-space case. Absolute continuity transfers positivity:
\[
(P_\theta Q)(\{z_S\})>0\qquad(S\in\mathcal K_5).
\]
For distinct \(S,T\in\mathcal K_5\), choose a measurable separating set
\(E_{S,T}\subseteq Z\) such that
\[
K(S,E_{S,T})=0,\qquad K(T,Z\setminus E_{S,T})=0.
\]
If \(z_S=z_T\), membership in this set gives the two cases
\[
\begin{aligned}
z_S\in E_{S,T}
&\quad\Longrightarrow\quad
0<K(S,\{z_S\})\le K(S,E_{S,T})=0,\\
z_S\notin E_{S,T}
&\quad\Longrightarrow\quad
0<K(T,\{z_T\})\le K(T,Z\setminus E_{S,T})=0.
\end{aligned}
\]
Hence the selected points are distinct, and
\[
5=|\mathcal K_5|
=\#\{z_S:S\in\mathcal K_5\}
\le\#\{z\in Z:(P_\theta Q)(\{z\})\ne0\}
=\kappa(Q).
\]
For an infinite output space the cardinality convention gives
\(\kappa(Q)=+\infty\). Thus in both cases
\[
\kappa(Q)\ge|\mathcal K_5|=5.
\]
This proves the universal assertion throughout \(\mathcal R_5\), and
completes all the claims.
% lean: CausalSmith.Stat.LdpAteEfficiencySurface.r5Certificate_outputCardinality_ge_five
\end{enumerate}
% lean: CausalSmith.Stat.LdpAteEfficiencySurface.five_output_region
\end{proof}

\begin{proof}[Proof of \cref{obj:thm:three-output-region}]
\begin{enumerate}
\item \textit{Certificate notation.}
Write the parameter domain and the active pattern set as
\[
\mathcal D:=(0,1)^3\times(0,\infty),
\qquad
\mathcal K:=\{S_1,S_6,S_8\}.
\]
For a parameter vector, use the coordinate notation
\[
\xi=(p(\xi),\mu_0(\xi),\mu_1(\xi),\varepsilon(\xi)),
\qquad
\theta(\xi):=(\mu_0(\xi),\mu_1(\xi))^\top.
\]
The fixed pattern enumeration is
\[
\mathcal S:=\{S_k:1\le k\le14\},
\qquad
S_k:=\{j\in\{0,1,2,3\}:\text{the \(j\)-th binary digit of \(k\) is \(1\)}\}.
\]
In particular,
\[
S_1=\{0\},\qquad S_6=\{1,2\},\qquad S_8=\{3\}.
\]
For \(S\in\mathcal S\), define its staircase ray by
\[
b_{S,j}(\varepsilon):=
\begin{cases}
e^\varepsilon,&j\in S,\\
1,&j\notin S,
\end{cases}
\qquad j\in\{0,1,2,3\},
\qquad
b_S(\varepsilon):=
(b_{S,0}(\varepsilon),\ldots,b_{S,3}(\varepsilon))^\top.
\]
The record probabilities, pattern gradients, and pattern masses are
\[
\pi_\theta(p):=
\begin{pmatrix}
(1-p)(1-\mu_0)\\
(1-p)\mu_0\\
p(1-\mu_1)\\
p\mu_1
\end{pmatrix},
\qquad
g_S(p,\varepsilon):=
\begin{pmatrix}
(1-p)(b_{S,1}(\varepsilon)-b_{S,0}(\varepsilon))\\
p(b_{S,3}(\varepsilon)-b_{S,2}(\varepsilon))
\end{pmatrix},
\]
\[
h_S(\theta;p,\varepsilon):=
\sum_{j=0}^3\pi_{\theta,j}(p)b_{S,j}(\varepsilon),
\qquad
v(t):=(t,t+1)^\top,\quad t\in\mathbb R.
\]
The additional arguments on \(h_S\) display its dependence on assignment and privacy. For \(\xi\in\mathcal D\), define the pattern information and its profiling slope by
\[
f_S(\xi,t):=
\frac{\bigl(g_S(p(\xi),\varepsilon(\xi))^\top v(t)\bigr)^2}
{h_S(\theta(\xi);p(\xi),\varepsilon(\xi))},
\]
\[
\sigma_S(\xi,t):=
\frac{
2\bigl(g_S(p(\xi),\varepsilon(\xi))^\top v(t)\bigr)
\bigl(g_{S,0}(p(\xi),\varepsilon(\xi))
      +g_{S,1}(p(\xi),\varepsilon(\xi))\bigr)}
{h_S(\theta(\xi);p(\xi),\varepsilon(\xi))}.
\]
The masses are positive on this domain, as shown below, so these quotients are well defined. The feasible set and objective are
\[
\mathcal A_\varepsilon:=
\left\{\beta\in[0,\infty)^{14}:
\sum_{S\in\mathcal S}\beta_S b_{S,j}(\varepsilon)=1
\text{ for every }j\in\{0,1,2,3\}\right\},
\]
\[
F_{\theta(\xi)}(\beta,t)
=\sum_{S\in\mathcal S}\beta_S f_S(\xi,t),
\]
where the objective in this display is evaluated at \(p(\xi)\) and \(\varepsilon(\xi)\). These are the quantities entering \cref{obj:def:staircase-saddle}. The three-pattern certificate of \cref{obj:def:support-regions} has support \(\mathcal K\), positive active weights, zero weighted profiling slope, active dual equalities, and strict inactive dual inequalities.
% lean: CausalSmith.Stat.LdpAteEfficiencySurface.certifiedRegion

\item \textit{Canonical weights for an arbitrary certificate.}
Fix \(\xi_0\in\mathcal R_3\), and choose witnessing quantities
\[
\beta^{(0)}\in\mathcal A_{\varepsilon(\xi_0)},
\qquad t_{\mathrm{wit}}\in\mathbb R,
\qquad \eta^{(0)}\in\mathbb R^4.
\]
Their support is \(\mathcal K\), so
\[
\beta^{(0)}_S=0\quad(S\notin\mathcal K).
\]
Put
\[
r_0:=e^{\varepsilon(\xi_0)}>1.
\]
Normalization at input coordinates \(0,1,3\) gives
\[
\begin{aligned}
r_0\beta^{(0)}_{S_1}+\beta^{(0)}_{S_6}
 +\beta^{(0)}_{S_8}&=1,\\
\beta^{(0)}_{S_1}+r_0\beta^{(0)}_{S_6}
 +\beta^{(0)}_{S_8}&=1,\\
\beta^{(0)}_{S_1}+\beta^{(0)}_{S_6}
 +r_0\beta^{(0)}_{S_8}&=1.
\end{aligned}
\]
Subtracting consecutive equations yields
\[
(r_0-1)(\beta^{(0)}_{S_1}-\beta^{(0)}_{S_6})=0,
\qquad
(r_0-1)(\beta^{(0)}_{S_6}-\beta^{(0)}_{S_8})=0.
\]
Consequently,
\[
\beta^{(0)}=\alpha^{(3)}(\varepsilon(\xi_0)),
\qquad
\alpha^{(3)}_S(\varepsilon):=
\begin{cases}
(e^\varepsilon+2)^{-1},&S\in\mathcal K,\\
0,&S\notin\mathcal K,
\end{cases}
\quad \varepsilon>0.
\]
For every \(\varepsilon>0\), the active weights are positive. Since the three active patterns partition the inputs,
\[
\sum_{S\in\mathcal K}b_{S,j}(\varepsilon)=e^\varepsilon+2
\quad(j=0,1,2,3).
\]
Thus
\[
\alpha^{(3)}(\varepsilon)\in\mathcal A_\varepsilon,
\qquad
\{S:\alpha^{(3)}_S(\varepsilon)\ne0\}=\mathcal K.
\]
% lean: CausalSmith.Stat.LdpAteEfficiencySurface.threeMaskWeights_unique
% lean: CausalSmith.Stat.LdpAteEfficiencySurface.threeMaskWeights_feasible
% lean: CausalSmith.Stat.LdpAteEfficiencySurface.threeMaskWeights_support

\item \textit{A continuous stationary direction.}
For every \(\xi\in\mathcal D\), all record probabilities are nonnegative and all staircase-ray entries are positive. In particular,
\[
h_S(\theta(\xi);p(\xi),\varepsilon(\xi))
\ge
(1-p(\xi))(1-\mu_0(\xi))b_{S,0}(\varepsilon(\xi))
>0
\quad(S\in\mathcal S).
\]
Define the canonical stationarity function and its linear coefficient by
\[
L_3(\xi,t):=
\sum_{S\in\mathcal S}
\alpha^{(3)}_S(\varepsilon(\xi))\sigma_S(\xi,t),
\qquad
a_3(\xi):=L_3(\xi,1)-L_3(\xi,0).
\]
The identity
\[
g_S^\top v(t)=g_{S,1}+t(g_{S,0}+g_{S,1})
\]
gives, with all quantities evaluated at \(\xi\),
\[
\sigma_S(\xi,t)
=
\sigma_S(\xi,0)
+
t\,\frac{2(g_{S,0}+g_{S,1})^2}{h_S(\theta(\xi))}.
\]
Hence
\[
L_3(\xi,t)=L_3(\xi,0)+t\,a_3(\xi),
\]
\[
a_3(\xi)=
\sum_{S\in\mathcal S}
\alpha^{(3)}_S(\varepsilon(\xi))
\frac{2(g_{S,0}+g_{S,1})^2}{h_S(\theta(\xi))}>0.
\]
Every summand is nonnegative, and the \(S_1\) summand is strictly positive because
\[
g_{S_1,0}+g_{S_1,1}
=-(e^{\varepsilon(\xi)}-1)(1-p(\xi))\ne0.
\]
We may therefore define, throughout \(\mathcal D\),
\[
t_3(\xi):=-\frac{L_3(\xi,0)}{a_3(\xi)},
\qquad
L_3(\xi,t_3(\xi))=0.
\]
The witnessing stationarity equation and the canonical-weight identity give
\[
0=L_3(\xi_0,t_{\mathrm{wit}})
=L_3(\xi_0,0)+t_{\mathrm{wit}}a_3(\xi_0),
\qquad
t_{\mathrm{wit}}=t_3(\xi_0).
\]

For each fixed pattern and input coordinate, the ray entry is either \(e^\varepsilon\) or \(1\), and is continuous. The masses are finite sums of products of continuous functions; the gradients are continuous as well. Their positive denominators show that \(L_3(\xi,0)\) and \(L_3(\xi,1)\) are continuous at \(\xi_0\). Since \(a_3(\xi_0)>0\), the quotient defining \(t_3\) is continuous there. It follows that
\[
\xi\longmapsto f_S(\xi,t_3(\xi))
\quad\text{is continuous at }\xi_0
\quad(S\in\mathcal S).
\]
% lean: CausalSmith.Stat.LdpAteEfficiencySurface.patternMass_pos_of_interior
% lean: CausalSmith.Stat.LdpAteEfficiencySurface.r3StationarityCoeff_pos
% lean: CausalSmith.Stat.LdpAteEfficiencySurface.r3Stationarity_affine
% lean: CausalSmith.Stat.LdpAteEfficiencySurface.continuousAt_r3Direction_region

\item \textit{A local dual solution anchored at the witness.}
Define the active dual matrix and target vector on \(\mathcal D\) by
\[
r(\xi):=e^{\varepsilon(\xi)},
\qquad
\mathsf A_3(\xi):=
\begin{pmatrix}
r(\xi)&1&1&1\\
1&r(\xi)&r(\xi)&1\\
1&1&1&r(\xi)
\end{pmatrix},
\qquad
\boldsymbol f_3(\xi):=
\begin{pmatrix}
f_{S_1}(\xi,t_3(\xi))\\
f_{S_6}(\xi,t_3(\xi))\\
f_{S_8}(\xi,t_3(\xi))
\end{pmatrix}.
\]
The active equalities at \(\xi_0\), together with
\(t_{\mathrm{wit}}=t_3(\xi_0)\), say
\[
\mathsf A_3(\xi_0)\eta^{(0)}=\boldsymbol f_3(\xi_0).
\]
Use the fixed right inverse
\[
\mathsf R_0:=
\frac{1}{(r_0-1)(r_0+2)}
\begin{pmatrix}
r_0+1&-1&-1\\
-1&r_0+1&-1\\
0&0&0\\
-1&-1&r_0+1
\end{pmatrix}.
\]
Its denominator is nonzero. Multiplication gives
\[
\mathsf A_3(\xi_0)\mathsf R_0
=
\frac{1}{(r_0-1)(r_0+2)}
\begin{pmatrix}
r_0(r_0+1)-2&0&0\\
0&r_0(r_0+1)-2&0\\
0&0&r_0(r_0+1)-2
\end{pmatrix}
=\operatorname{diag}(1,1,1).
\]
The matrix \(\mathsf A_3\) and the target \(\boldsymbol f_3\) are continuous at \(\xi_0\). Since
\[
\det(\mathsf A_3(\xi_0)\mathsf R_0)=1,
\]
continuity of the determinant gives a neighborhood of \(\xi_0\) within
\[
\mathcal U_{\mathrm{inv}}
:=
\{\xi\in\mathcal D:
\det(\mathsf A_3(\xi)\mathsf R_0)\ne0\}.
\]
On this set define
\[
\eta_{\mathrm{loc}}(\xi):=
\eta^{(0)}
+
\mathsf R_0
(\mathsf A_3(\xi)\mathsf R_0)^{-1}
\bigl(\boldsymbol f_3(\xi)-\mathsf A_3(\xi)\eta^{(0)}\bigr).
\]
The inverse in this formula is taken on the stated invertible domain. Continuity of matrix inversion and the displayed formula imply continuity of \(\eta_{\mathrm{loc}}\) at \(\xi_0\). The residual vanishes at the base point, so
\[
\eta_{\mathrm{loc}}(\xi_0)=\eta^{(0)}.
\]
Moreover, for every \(\xi\in\mathcal U_{\mathrm{inv}}\),
\[
\begin{aligned}
\mathsf A_3(\xi)\eta_{\mathrm{loc}}(\xi)
&=
\mathsf A_3(\xi)\eta^{(0)}
+
\mathsf A_3(\xi)\mathsf R_0
(\mathsf A_3(\xi)\mathsf R_0)^{-1}
\bigl(\boldsymbol f_3(\xi)-\mathsf A_3(\xi)\eta^{(0)}\bigr)\\
&=\boldsymbol f_3(\xi).
\end{aligned}
\]
Thus the three active dual equalities hold locally, and the original inactive strict inequalities hold at the base point:
\[
f_S(\xi_0,t_3(\xi_0))
<
b_S(\varepsilon(\xi_0))^\top\eta_{\mathrm{loc}}(\xi_0)
\quad(S\notin\mathcal K).
\]
% lean: CausalSmith.Stat.LdpAteEfficiencySurface.r3DualMatrix_mul_rightInverse
% lean: CausalSmith.Stat.LdpAteEfficiencySurface.exists_eventually_anchoredSolution_of_rightInverse
% lean: CausalSmith.Stat.LdpAteEfficiencySurface.exists_local_r3Dual_of_mem

\item \textit{Persistence of the certificate.}
For every \(S\in\mathcal S\), the dual value
\[
b_S(\varepsilon(\xi))^\top\eta_{\mathrm{loc}}(\xi)
=
\sum_{j=0}^3b_{S,j}(\varepsilon(\xi))
\eta_{\mathrm{loc},j}(\xi)
\]
is continuous at \(\xi_0\), being a finite sum of products of continuous functions. Consequently, each inactive strict inequality persists in a neighborhood of \(\xi_0\). Intersect these finitely many neighborhoods with the neighborhood on which the active equations hold. The strict coordinate inequalities
\[
0<p(\xi)<1,\qquad
0<\mu_0(\xi)<1,\qquad
0<\mu_1(\xi)<1,\qquad
0<\varepsilon(\xi)
\]
also persist locally. We obtain a neighborhood \(\mathcal N\) of \(\xi_0\) such that, for every \(\xi\in\mathcal N\),
\[
\begin{gathered}
\xi\in\mathcal D,\qquad
\alpha^{(3)}(\varepsilon(\xi))\in\mathcal A_{\varepsilon(\xi)},\\
\{S:\alpha^{(3)}_S(\varepsilon(\xi))\ne0\}=\mathcal K,
\qquad
\alpha^{(3)}_S(\varepsilon(\xi))>0\quad(S\in\mathcal K),\\
\sum_{S\in\mathcal S}
\alpha^{(3)}_S(\varepsilon(\xi))\sigma_S(\xi,t_3(\xi))=0,\\
b_S(\varepsilon(\xi))^\top\eta_{\mathrm{loc}}(\xi)
=f_S(\xi,t_3(\xi))
\quad(S\in\mathcal K),\\
f_S(\xi,t_3(\xi))
<
b_S(\varepsilon(\xi))^\top\eta_{\mathrm{loc}}(\xi)
\quad(S\notin\mathcal K).
\end{gathered}
\]
These quantities constitute the certificate in \cref{obj:def:support-regions}, so
\[
\mathcal N\subseteq\mathcal R_3.
\]
As \(\xi_0\) was arbitrary, \(\mathcal R_3\) is open in \(\mathbb R^4\).
% lean: CausalSmith.Stat.LdpAteEfficiencySurface.isOpen_R3_continuousCertificate

\item \textit{Primal feasibility and stationarity at the stated point.}
Define
\[
\xi_{\mathrm{low}}
:=
\left(\frac3{10},\frac1{10},\frac1{10},\log3\right),
\qquad
\theta_{\mathrm{low}}:=\theta(\xi_{\mathrm{low}})
=\left(\frac1{10},\frac1{10}\right)^\top,
\qquad
t_{\mathrm{low}}:=-\frac{339}{9985},
\]
\[
\alpha:=\alpha^{(3)}(\log3),
\qquad
\alpha_S=
\begin{cases}
1/5,&S\in\mathcal K,\\
0,&S\notin\mathcal K.
\end{cases}
\]
Since \(\log3>0\) and \(e^{\log3}=3\), the canonical-weight calculation yields
\[
\alpha\in\mathcal A_{\log3},
\qquad
\{S:\alpha_S\ne0\}=\mathcal K,
\qquad
\alpha_{S_1}=\alpha_{S_6}=\alpha_{S_8}=\frac15.
\]
The record probabilities are
\[
\pi_{\theta_{\mathrm{low}}}(3/10)
=\frac1{100}(63,7,27,3)^\top.
\]
Substitution into the pattern formulas gives
\[
\begin{array}{c|c|c}
S&h_S(\theta_{\mathrm{low}})&g_S\\ \hline
S_1&113/50&(-7/5,0)^\top\\
S_6&42/25&(7/5,-3/5)^\top\\
S_8&53/50&(0,3/5)^\top
\end{array}
\]
and hence, for every \(t\in\mathbb R\),
\[
F_{\theta_{\mathrm{low}}}(\alpha,t)
=
\frac15\left[
\frac{98}{113}t^2
+\frac{(4t-3)^2}{42}
+\frac{18}{53}(t+1)^2
\right].
\]
Here and below objectives at \(\theta_{\mathrm{low}}\) are evaluated at \(p=3/10\) and \(\varepsilon=\log3\). The weighted slope at \(t_{\mathrm{low}}\) is therefore
\[
\sum_{S\in\mathcal S}\alpha_S
\sigma_S(\xi_{\mathrm{low}},t_{\mathrm{low}})
=
\frac15\left[
\frac{196}{113}t_{\mathrm{low}}
+\frac{8(4t_{\mathrm{low}}-3)}{42}
+\frac{36}{53}(t_{\mathrm{low}}+1)
\right]
=0.
\]
This establishes the required profiling stationarity.
% lean: CausalSmith.Stat.LdpAteEfficiencySurface.lowMeanCandidate_feasible
% lean: CausalSmith.Stat.LdpAteEfficiencySurface.lowMeanCandidate_stationary

\item \textit{The exact dual certificate and region membership.}
Define the four-coordinate dual vector by
\[
\eta_{\mathrm{low}}:=
\begin{pmatrix}
-21817593/398800900\\
32820291/8773619800\\
509870781/8773619800\\
41183667/398800900
\end{pmatrix}.
\]
Its coordinate sum is
\[
\sum_{j=0}^3\eta_{\mathrm{low},j}
=
-\frac{21817593}{398800900}
+\frac{32820291}{8773619800}
+\frac{509870781}{8773619800}
+\frac{41183667}{398800900}
=\frac{441}{3994}.
\]
For every \(S\in\mathcal S\), define its slack and the specialized ray entries by
\[
\Delta_S:=
b_S(\log3)^\top\eta_{\mathrm{low}}
-f_S(\xi_{\mathrm{low}},t_{\mathrm{low}}),
\qquad
b^{\mathrm{low}}_{S,j}:=b_{S,j}(\log3)
=
\begin{cases}
3,&j\in S,\\
1,&j\notin S.
\end{cases}
\]
The slack is computed directly as
\[
\Delta_S
=
\sum_{j=0}^3 b^{\mathrm{low}}_{S,j}\eta_{\mathrm{low},j}
-
\frac{
\left[
\frac7{10}(b^{\mathrm{low}}_{S,1}-b^{\mathrm{low}}_{S,0})t_{\mathrm{low}}
+
\frac3{10}(b^{\mathrm{low}}_{S,3}-b^{\mathrm{low}}_{S,2})(t_{\mathrm{low}}+1)
\right]^2}
{
\bigl(
63b^{\mathrm{low}}_{S,0}
+7b^{\mathrm{low}}_{S,1}
+27b^{\mathrm{low}}_{S,2}
+3b^{\mathrm{low}}_{S,3}
\bigr)/100
}.
\]
For example, \(S_2=\{1\}\) has ray \((1,3,1,1)^\top\), so
\[
\Delta_{S_2}
=
\frac{441}{3994}
+2\frac{32820291}{8773619800}
-\frac{98}{57}\left(\frac{339}{9985}\right)^2
=
\frac{1932296079}{16669877620}.
\]
Applying the same displayed formula to each of the fourteen binary patterns gives
\[
\begin{array}{c|c@{\qquad}c|c}
S&\Delta_S&S&\Delta_S\\ \hline
S_1&0&S_8&0\\
S_2&1932296079/16669877620&S_9&4394565/115652261\\
S_3&7441119/877361980&S_{10}&77641557/877361980\\
S_4&7441119/877361980&S_{11}&2821826721/35971841180\\
S_5&14134743/877361980&S_{12}&380056761/877361980\\
S_6&0&S_{13}&3683761767/11405705740\\
S_7&41610/3988009&S_{14}&50813490/115652261
\end{array}
\]
Thus the active and inactive cases satisfy, respectively,
\[
b_S(\log3)^\top\eta_{\mathrm{low}}
=f_S(\xi_{\mathrm{low}},t_{\mathrm{low}})
\quad(S\in\mathcal K),
\]
\[
f_S(\xi_{\mathrm{low}},t_{\mathrm{low}})
<
b_S(\log3)^\top\eta_{\mathrm{low}}
\quad(S\notin\mathcal K).
\]
The point \(\xi_{\mathrm{low}}\) belongs to \(\mathcal D\). Together with the feasible positive active weights and stationarity established above, these equalities and strict inequalities supply its certificate in \cref{obj:def:support-regions}. Therefore
\[
\xi_{\mathrm{low}}\in\mathcal R_3.
\]
% lean: CausalSmith.Stat.LdpAteEfficiencySurface.lowMeanDual_value
% lean: CausalSmith.Stat.LdpAteEfficiencySurface.lowMeanDual_slack
% lean: CausalSmith.Stat.LdpAteEfficiencySurface.lowMeanPoint_mem_R3

\item \textit{A uniform upper bound on profiled information.}
Let \(\beta\in\mathcal A_{\log3}\). Positive pattern masses and nonnegative weights imply
\[
F_{\theta_{\mathrm{low}}}(\beta,t)
=
\sum_{S\in\mathcal S}\beta_S
\frac{(g_S^\top v(t))^2}{h_S(\theta_{\mathrm{low}})}
\ge0
\quad(t\in\mathbb R).
\]
Thus its set of objective values,
\[
\mathcal U_\beta:=
\{F_{\theta_{\mathrm{low}}}(\beta,t):t\in\mathbb R\},
\]
is nonempty and bounded below by zero. In particular,
\[
\inf_{t\in\mathbb R}F_{\theta_{\mathrm{low}}}(\beta,t)
\le F_{\theta_{\mathrm{low}}}(\beta,t_{\mathrm{low}}).
\]
The active equality and inactive strict-inequality cases both give
\[
f_S(\xi_{\mathrm{low}},t_{\mathrm{low}})
\le b_S(\log3)^\top\eta_{\mathrm{low}}
\quad(S\in\mathcal S).
\]
Multiplication by \(\beta_S\ge0\) and summation yield
\[
F_{\theta_{\mathrm{low}}}(\beta,t_{\mathrm{low}})
\le
\sum_{S\in\mathcal S}\beta_S b_S(\log3)^\top\eta_{\mathrm{low}}.
\]
Normalization identifies the latter sum:
\[
\begin{aligned}
\sum_{S\in\mathcal S}\beta_S b_S(\log3)^\top\eta_{\mathrm{low}}
&=
\sum_{j=0}^3\eta_{\mathrm{low},j}
\left(\sum_{S\in\mathcal S}\beta_S b_{S,j}(\log3)\right)\\
&=\sum_{j=0}^3\eta_{\mathrm{low},j}
=\frac{441}{3994}.
\end{aligned}
\]
Consequently, for every feasible \(\beta\),
\[
\inf_{t\in\mathbb R}F_{\theta_{\mathrm{low}}}(\beta,t)
\le
F_{\theta_{\mathrm{low}}}(\beta,t_{\mathrm{low}})
\le\frac{441}{3994}.
\]
% lean: CausalSmith.Stat.LdpAteEfficiencySurface.lowMean_information_nonneg
% lean: CausalSmith.Stat.LdpAteEfficiencySurface.feasible_dualRay_sum
% lean: CausalSmith.Stat.LdpAteEfficiencySurface.lowMean_dual_upper_bound

\item \textit{The candidate infimum and optimized value.}
Expanding the candidate objective gives
\[
F_{\theta_{\mathrm{low}}}(\alpha,t)
=
\frac{39940}{125769}t^2+\frac8{371}t+\frac{411}{3710}
=
\frac{441}{3994}
+
\frac{39940}{125769}
\left(t+\frac{339}{9985}\right)^2.
\]
The coefficient of the square is positive. Therefore
\[
F_{\theta_{\mathrm{low}}}(\alpha,t)\ge\frac{441}{3994}
\quad(t\in\mathbb R),
\qquad
F_{\theta_{\mathrm{low}}}(\alpha,t_{\mathrm{low}})
=\frac{441}{3994},
\]
and hence
\[
\inf_{t\in\mathbb R}F_{\theta_{\mathrm{low}}}(\alpha,t)
=\frac{441}{3994}.
\]
Define the set of feasible profiled values by
\[
\mathcal E_{\mathrm{prof}}:=
\left\{
\inf_{t\in\mathbb R}F_{\theta_{\mathrm{low}}}(\beta,t):
\beta\in\mathcal A_{\log3}
\right\}.
\]
The uniform bound gives
\[
z\le\frac{441}{3994}\quad(z\in\mathcal E_{\mathrm{prof}}),
\]
while the feasible candidate gives
\[
\frac{441}{3994}\in\mathcal E_{\mathrm{prof}}.
\]
Thus this nonempty, bounded-above set has supremum \(441/3994\). By the optimized-value definition in \cref{obj:def:staircase-saddle},
\[
J^*(\theta_{\mathrm{low}},3/10,\log3)
=
\sup\mathcal E_{\mathrm{prof}}
=
\frac{441}{3994}
=
\inf_{t\in\mathbb R}F_{\theta_{\mathrm{low}}}(\alpha,t).
\]
This proves attainment by the stated weight vector.
% lean: CausalSmith.Stat.LdpAteEfficiencySurface.lowMeanCandidate_information_value
% lean: CausalSmith.Stat.LdpAteEfficiencySurface.lowMeanCandidate_global_min
% lean: CausalSmith.Stat.LdpAteEfficiencySurface.lowMeanCandidate_inf_value
% lean: CausalSmith.Stat.LdpAteEfficiencySurface.lowMean_Jstar_value

\item \textit{Uniqueness of the optimizing weights.}
Suppose that \(\beta\in\mathcal A_{\log3}\) satisfies
\[
\inf_{t\in\mathbb R}F_{\theta_{\mathrm{low}}}(\beta,t)
=
J^*(\theta_{\mathrm{low}},3/10,\log3).
\]
Its objective-value set is bounded below by zero, so the infimum is bounded above by its value at \(t_{\mathrm{low}}\). The dual upper bound then gives
\[
\frac{441}{3994}
=
\inf_{t\in\mathbb R}F_{\theta_{\mathrm{low}}}(\beta,t)
\le
F_{\theta_{\mathrm{low}}}(\beta,t_{\mathrm{low}})
\le
\frac{441}{3994}.
\]
All these values are equal. In particular, normalization of the dual sum yields
\[
\sum_{S\in\mathcal S}\beta_S f_S(\xi_{\mathrm{low}},t_{\mathrm{low}})
=
\frac{441}{3994}
=
\sum_{S\in\mathcal S}\beta_S
b_S(\log3)^\top\eta_{\mathrm{low}}.
\]
For an active pattern the corresponding terms are equal. For an inactive pattern, nonnegative multiplication of its strict slack gives
\[
\beta_S f_S(\xi_{\mathrm{low}},t_{\mathrm{low}})
\le
\beta_S b_S(\log3)^\top\eta_{\mathrm{low}}.
\]
Since these finitely many termwise inequalities have equal sums, equality holds in every term:
\[
\beta_S f_S(\xi_{\mathrm{low}},t_{\mathrm{low}})
=
\beta_S b_S(\log3)^\top\eta_{\mathrm{low}}
\quad(S\in\mathcal S).
\]
If \(S\notin\mathcal K\) and \(\beta_S\ne0\), feasibility implies \(\beta_S>0\). Its strict slack would then give
\[
\beta_S f_S(\xi_{\mathrm{low}},t_{\mathrm{low}})
<
\beta_S b_S(\log3)^\top\eta_{\mathrm{low}},
\]
a contradiction. Hence
\[
\beta_S=0\quad(S\notin\mathcal K).
\]

The normalization equations at input coordinates \(0,1,3\) now reduce to
\[
\begin{aligned}
3\beta_{S_1}+\beta_{S_6}+\beta_{S_8}&=1,\\
\beta_{S_1}+3\beta_{S_6}+\beta_{S_8}&=1,\\
\beta_{S_1}+\beta_{S_6}+3\beta_{S_8}&=1.
\end{aligned}
\]
Their differences give
\[
2(\beta_{S_1}-\beta_{S_6})=0,
\qquad
2(\beta_{S_6}-\beta_{S_8})=0,
\]
and substitution gives
\[
\beta_{S_1}=\beta_{S_6}=\beta_{S_8}=\frac15.
\]
For \(S\in\mathcal K\), these weights equal \(\alpha_S\); for \(S\notin\mathcal K\), both weights are zero. Therefore \(\beta=\alpha\).
% lean: CausalSmith.Stat.LdpAteEfficiencySurface.three_output_region

\item \textit{Reciprocal variance.}
Finally, the reciprocal definition in \cref{obj:def:staircase-saddle} and the positive optimized value give
\[
V^*(\theta_{\mathrm{low}},3/10,\log3)
=
\left(J^*(\theta_{\mathrm{low}},3/10,\log3)\right)^{-1}
=
\left(\frac{441}{3994}\right)^{-1}
=
\frac{3994}{441}.
\]
% lean: CausalSmith.Stat.LdpAteEfficiencySurface.three_output_region
\end{enumerate}
\end{proof}

\begin{proof}[Proof of \cref{obj:thm:support-transition}]
Throughout, fix \(p=3/10\) and \(\varepsilon=\log 3\). Set
\[
A_5=\{S_3,S_4,S_6,S_8,S_9\},
\qquad
A_3=\{S_1,S_6,S_8\},
\qquad
\iota:\mathbb R^2\longrightarrow\mathbb R^4,
\quad
\iota(\theta)=\left(\frac3{10},\mu_0,\mu_1,\log 3\right),
\]
where \(\theta=(\mu_0,\mu_1)^\top\). We use the objective, rays, and certificates of
\cref{obj:synth_4,obj:synth_5}, with the pattern enumeration in
\cref{obj:synth_8}.

\begin{enumerate}
\item \emph{Construct the open slices.}
By \cref{obj:thm:five-output-region,obj:thm:three-output-region},
\[
\mathcal R_5,\mathcal R_3\ \text{are open},\qquad
\iota\bigl((1/2,1/2)^\top\bigr)\in\mathcal R_5,
\qquad
\iota\bigl((1/10,1/10)^\top\bigr)\in\mathcal R_3.
\]
Define
\[
\mathcal V_5
=\iota^{-1}(\mathcal R_5)\cap\{\theta:\mu_0>3/10\},
\qquad
\mathcal V_3
=\iota^{-1}(\mathcal R_3)\cap\{\theta:\mu_0<3/10\}.
\]
The map \(\iota\) is continuous, so both sets are open. Their respective centers belong to them because
\[
\frac1{10}<\frac3{10}<\frac12.
\]
They are disjoint: membership in their intersection would require both
\(\mu_0>3/10\) and \(\mu_0<3/10\). Moreover,
\[
\theta\in\mathcal V_5\Longrightarrow\iota(\theta)\in\mathcal R_5,
\qquad
\theta\in\mathcal V_3\Longrightarrow\iota(\theta)\in\mathcal R_3.
\]
% lean: CausalSmith.Stat.LdpAteEfficiencySurface.supportRegion_slices

\item \emph{Preserve the three distinct scores near the low-mean point.}
For this neighborhood construction, let
\[
\mathcal O=(0,1)^3\times(0,\infty),\qquad
\xi=(p_\xi,\mu_{0,\xi},\mu_{1,\xi},\varepsilon_\xi)\in\mathcal O,
\qquad
\theta_\xi=(\mu_{0,\xi},\mu_{1,\xi})^\top.
\]
Write the record probabilities, rays, masses, and gradients at these parameters as
\[
\begin{aligned}
\pi_\xi&=\bigl((1-p_\xi)(1-\mu_{0,\xi}),
 (1-p_\xi)\mu_{0,\xi},
 p_\xi(1-\mu_{1,\xi}),p_\xi\mu_{1,\xi}\bigr)^\top,\\
b_S(\xi)&=\mathbf1_4+(e^{\varepsilon_\xi}-1)\mathbf1_S,
\qquad h_S(\xi)=\pi_\xi^\top b_S(\xi),\\
g_S(\xi)&=(e^{\varepsilon_\xi}-1)
\begin{pmatrix}
(1-p_\xi)\{\mathbf1_S(1)-\mathbf1_S(0)\}\\
p_\xi\{\mathbf1_S(3)-\mathbf1_S(2)\}
\end{pmatrix},
\qquad
\mathbf1_2=\begin{pmatrix}1\\1\end{pmatrix}.
\end{aligned}
\]
Here \(\mathbf1_S\) is the indicator vector of the pattern \(S\). In particular,
\(h_S(\xi)>0\) for every \(S\).

Define the canonical three-pattern weights and their stationarity coefficients by
\[
\alpha^{\mathrm{can}}_S(\xi)=
\begin{cases}
(e^{\varepsilon_\xi}+2)^{-1},&S\in A_3,\\
0,&S\notin A_3,
\end{cases}
\]
\[
a_3(\xi)=\frac{2}{e^{\varepsilon_\xi}+2}
 \sum_{S\in A_3}\frac{(g_S(\xi)^\top\mathbf1_2)^2}{h_S(\xi)},
\qquad
\ell_3(\xi)=\frac{2}{e^{\varepsilon_\xi}+2}
 \sum_{S\in A_3}
 \frac{(g_S(\xi)^\top v(0))(g_S(\xi)^\top\mathbf1_2)}{h_S(\xi)}.
\]
Expansion of the quadratic objective gives
\[
\frac{\partial}{\partial t}
F_{\theta_\xi}\bigl(\alpha^{\mathrm{can}}(\xi),t\bigr)
=\ell_3(\xi)+a_3(\xi)t,
\]
where the objective is evaluated at \(p_\xi,\varepsilon_\xi\). The \(S_1\) summand yields
\[
a_3(\xi)\ge
\frac{2(e^{\varepsilon_\xi}-1)^2(1-p_\xi)^2}
 {(e^{\varepsilon_\xi}+2)h_{S_1}(\xi)}>0.
\]
Thus the unique stationary coordinate is the function
\[
t_3:\mathcal O\longrightarrow\mathbb R,
\qquad
t_3(\xi)=-\frac{\ell_3(\xi)}{a_3(\xi)}.
\]
All denominators in these formulas are positive, so \(t_3\) and the three functions
\[
r_{3,S}(\xi)=
\frac{g_S(\xi)^\top v(t_3(\xi))}{h_S(\xi)},
\qquad S\in A_3,
\]
are continuous at
\[
\xi_{\mathrm{low}}
=\left(\frac3{10},\frac1{10},\frac1{10},\log 3\right).
\]

To calculate their values, the record probabilities at this point are
\[
\pi_{\xi_{\mathrm{low}}}
=\left(\frac{63}{100},\frac7{100},\frac{27}{100},\frac3{100}\right)^\top.
\]
Using \(S_1=\{0\}\), \(S_6=\{1,2\}\), and \(S_8=\{3\}\), the definitions give
\[
\begin{array}{c|c|c}
S&h_S(\xi_{\mathrm{low}})&g_S(\xi_{\mathrm{low}})\\ \hline
S_1&113/50&(-7/5,0)^\top\\
S_6&42/25&(7/5,-3/5)^\top\\
S_8&53/50&(0,3/5)^\top
\end{array}
\]
and hence
\[
\ell_3(\xi_{\mathrm{low}})=\frac8{371},
\qquad
a_3(\xi_{\mathrm{low}})=\frac{79880}{125769},
\qquad
t_3(\xi_{\mathrm{low}})=-\frac{339}{9985}.
\]
Substitution into the projected-score formula gives
\[
r_{3,S_1}(\xi_{\mathrm{low}})=\frac{42}{1997},
\qquad
r_{3,S_6}(\xi_{\mathrm{low}})=-\frac{1491}{3994},
\qquad
r_{3,S_8}(\xi_{\mathrm{low}})=\frac{1092}{1997}.
\]
These values are pairwise distinct: the middle value is negative, the others are positive, and
\[
\frac{1092}{1997}-\frac{42}{1997}=\frac{1050}{1997}>0.
\]
For each ordered pair \(S,T\in A_3\) with \(S\ne T\), continuity of their difference therefore supplies an open neighborhood \(D_{S,T}\) satisfying
\[
\xi_{\mathrm{low}}\in D_{S,T}
\subseteq
\{\xi\in\mathcal O:r_{3,S}(\xi)-r_{3,T}(\xi)\ne0\}.
\]
Define
\[
D_3=\bigcap_{\substack{S,T\in A_3\\S\ne T}}D_{S,T}.
\]
This finite intersection is open, contains \(\xi_{\mathrm{low}}\), and satisfies
\[
\xi\in D_3,\quad S,T\in A_3,\quad S\ne T
\Longrightarrow r_{3,S}(\xi)\ne r_{3,T}(\xi).
\]
% lean: CausalSmith.Stat.LdpAteEfficiencySurface.exists_lowMean_r3_distinctScore_neighborhood

\item \emph{Choose the final neighborhoods.}
Set
\[
U_5=\mathcal V_5,
\qquad
U_3=\mathcal V_3\cap\iota^{-1}(D_3).
\]
Continuity of \(\iota\) makes \(U_3\) open. The center conditions established above give
\[
(1/2,1/2)^\top\in U_5,
\qquad
(1/10,1/10)^\top\in U_3.
\]
Since \(U_3\subseteq\mathcal V_3\), the sets \(U_5,U_3\) are disjoint. Every point in either neighborhood has interior arm means, by its corresponding region membership.
% lean: CausalSmith.Stat.LdpAteEfficiencySurface.support_transition

\item \emph{Obtain the five-pattern optimum and support rigidity.}
Fix \(\theta\in U_5\). Membership in \(\mathcal R_5\), together with
\cref{obj:def:support-regions,obj:synth_5}, supplies
\[
\alpha^{(5)}\in\mathcal A_\varepsilon,\qquad
t_{\mathrm{cert},5}\in\mathbb R,\qquad
\eta^{(5)}\in\mathbb R^4,
\]
with
\[
\alpha^{(5)}_S>0\quad(S\in A_5),\qquad
\alpha^{(5)}_S=0\quad(S\notin A_5),
\qquad
\frac{\partial}{\partial t}
F_\theta(\alpha^{(5)},t_{\mathrm{cert},5})=0,
\]
and
\[
f_S(t_{\mathrm{cert},5})=
 b_S^\top\eta^{(5)}\quad(S\in A_5),
\qquad
f_S(t_{\mathrm{cert},5})<
 b_S^\top\eta^{(5)}\quad(S\notin A_5).
\]
Here \(f_S(t)=(g_S^\top v(t))^2/h_S(\theta)\), as in
\cref{obj:synth_5}. Stationarity gives the exact expansion
\[
F_\theta(\alpha^{(5)},t)-F_\theta(\alpha^{(5)},t_{\mathrm{cert},5})
=(t-t_{\mathrm{cert},5})^2
 \sum_{S\in\mathcal S}
 \alpha^{(5)}_S\frac{(g_S^\top\mathbf1_2)^2}{h_S(\theta)}
\ge0.
\]
For every feasible \(\beta\), normalization and the dual inequalities give
\[
\begin{aligned}
\inf_{t\in\mathbb R}F_\theta(\beta,t)
&\le F_\theta(\beta,t_{\mathrm{cert},5})\\
&\le\sum_{S\in\mathcal S}\beta_S b_S^\top\eta^{(5)}
=\mathbf1_4^\top\eta^{(5)}
=F_\theta(\alpha^{(5)},t_{\mathrm{cert},5}).
\end{aligned}
\]
The last equality uses the active dual equalities and the vanishing inactive weights. Since \(\alpha^{(5)}\) attains its minimum at the certifying coordinate,
\[
J^*(\theta,p,\varepsilon)
=\inf_tF_\theta(\alpha^{(5)},t)
=F_\theta(\alpha^{(5)},t_{\mathrm{cert},5})
=\mathbf1_4^\top\eta^{(5)}.
\]

If a feasible \(\beta\) satisfies \(J^*=\inf_tF_\theta(\beta,t)\), the preceding inequalities are equalities. Consequently,
\[
0=\sum_{S\in\mathcal S}\beta_S
 \bigl(b_S^\top\eta^{(5)}-f_S(t_{\mathrm{cert},5})\bigr).
\]
Every summand is nonnegative, and every inactive slack is strictly positive. Therefore
\[
\beta_S=0\qquad(S\notin A_5).
\]
% lean: CausalSmith.Stat.LdpAteEfficiencySurface.strictCertificate_profile_optimal_and_rigid

\item \emph{Determine the five-pattern mixture uniquely.}
For a feasible \(\beta\) vanishing outside \(A_5\), the four normalization equations, using \(e^\varepsilon=3\), are
\[
\begin{aligned}
3\beta_{S_3}+\beta_{S_4}+\beta_{S_6}+\beta_{S_8}+3\beta_{S_9}&=1,\\
3\beta_{S_3}+\beta_{S_4}+3\beta_{S_6}+\beta_{S_8}+\beta_{S_9}&=1,\\
\beta_{S_3}+3\beta_{S_4}+3\beta_{S_6}+\beta_{S_8}+\beta_{S_9}&=1,\\
\beta_{S_3}+\beta_{S_4}+\beta_{S_6}+3\beta_{S_8}+3\beta_{S_9}&=1.
\end{aligned}
\]
Subtracting these equations yields
\[
\beta_{S_3}=\beta_{S_4}=\beta_{S_8},
\qquad
\beta_{S_6}=\beta_{S_9},
\qquad
5\beta_{S_3}+4\beta_{S_6}=1.
\]
Thus, with the mixture coordinate defined by
\[
\lambda(\beta)=4\beta_{S_6},
\]
the weights satisfy
\[
\beta_{S_6}=\beta_{S_9}=\frac{\lambda(\beta)}4,
\qquad
\beta_{S_3}=\beta_{S_4}=\beta_{S_8}
=\frac{1-\lambda(\beta)}5.
\]

Define the two normalized component objectives and their coefficients by
\[
B_5(t)=\frac{f_{S_6}(t)+f_{S_9}(t)}4,
\qquad
T_5(t)=\frac{f_{S_3}(t)+f_{S_4}(t)+f_{S_8}(t)}5,
\]
\[
C_{\mathrm{bin}}=\frac1{100}
 \left(\frac1{h_{S_6}(\theta)}+\frac1{h_{S_9}(\theta)}\right)>0,
\qquad
C_{\mathrm{tri}}=\frac9{125}
 \left(\frac1{h_{S_4}(\theta)}+\frac1{h_{S_8}(\theta)}\right)>0.
\]
The pattern gradients are
\[
g_{S_3}=0,\quad
g_{S_4}=(0,-3/5)^\top,\quad
g_{S_6}=(7/5,-3/5)^\top,\quad
g_{S_8}=(0,3/5)^\top,\quad
g_{S_9}=(-7/5,3/5)^\top,
\]
so direct substitution gives
\[
B_5(t)=C_{\mathrm{bin}}(4t-3)^2,
\qquad
T_5(t)=C_{\mathrm{tri}}(t+1)^2.
\]
Also,
\[
\frac{b_{S_6}+b_{S_9}}4
=\frac{b_{S_3}+b_{S_4}+b_{S_8}}5
=\mathbf1_4.
\]
The active dual equalities therefore imply
\[
B_5(t_{\mathrm{cert},5})
=T_5(t_{\mathrm{cert},5})
=\mathbf1_4^\top\eta^{(5)}.
\]
In particular,
\[
T_5'(t_{\mathrm{cert},5})
=2C_{\mathrm{tri}}(t_{\mathrm{cert},5}+1)\ne0.
\]
Indeed, a zero derivative would give \(t_{\mathrm{cert},5}=-1\), whereas the component equality would then require
\(49C_{\mathrm{bin}}=0\), contradicting \(C_{\mathrm{bin}}>0\).

For the certificate weights, the mixture representation gives
\[
F_\theta(\alpha^{(5)},t)
=\lambda(\alpha^{(5)})B_5(t)
 +(1-\lambda(\alpha^{(5)}))T_5(t).
\]
Its stationarity consequently implies
\[
\lambda(\alpha^{(5)})B_5'(t_{\mathrm{cert},5})
 +(1-\lambda(\alpha^{(5)}))T_5'(t_{\mathrm{cert},5})=0.
\]
The derivative difference is nonzero: if the two derivatives were equal, this equation would force \(T_5'(t_{\mathrm{cert},5})=0\). Solving gives
\[
\lambda(\alpha^{(5)})
=-\frac{T_5'(t_{\mathrm{cert},5})}
 {B_5'(t_{\mathrm{cert},5})-T_5'(t_{\mathrm{cert},5})}.
\]
For any feasible optimal \(\beta\), the preceding support-rigidity argument applies, and
\[
F_\theta(\beta,t_{\mathrm{cert},5})
=J^*=\inf_tF_\theta(\beta,t).
\]
Hence \(\beta\) is stationary at the same coordinate and satisfies the same mixture equation. Thus
\[
\lambda(\beta)=\lambda(\alpha^{(5)}),
\qquad
\beta=\alpha^{(5)}.
\]
This proves the required uniqueness among all feasible optimal weights.
% lean: CausalSmith.Stat.LdpAteEfficiencySurface.uniqueOptimalSupport_of_mem_R5

\item \emph{Establish exact cardinality five.}
We first record positivity of the optimized information throughout the interior fixed-parameter slice. Define the feasible singleton design by
\[
\alpha^{\mathrm{rr}}_S=
\begin{cases}
1/6,&S\in\{S_1,S_2,S_4,S_8\},\\
0,&\text{otherwise}.
\end{cases}
\]
Each input row sums to \((3+1+1+1)/6=1\). Its objective satisfies, for every \(t\in\mathbb R\),
\[
F_\theta(\alpha^{\mathrm{rr}},t)
\ge\frac{49}{150h_{S_1}(\theta)}t^2,
\qquad
F_\theta(\alpha^{\mathrm{rr}},t)
\ge\frac{3}{50h_{S_8}(\theta)}(t+1)^2.
\]
It is a coercive quadratic, so choose
\[
t_{\mathrm{rr}}\in\arg\min_{t\in\mathbb R}
F_\theta(\alpha^{\mathrm{rr}},t).
\]
A zero minimum would force both \(t_{\mathrm{rr}}=0\) and
\(t_{\mathrm{rr}}=-1\). Therefore
\[
0<F_\theta(\alpha^{\mathrm{rr}},t_{\mathrm{rr}})
\le J^*(\theta,p,\varepsilon).
\]
% lean: CausalSmith.Stat.LdpAteEfficiencySurface.Jstar_pos_interior

For the five-pattern certificate, define
\[
I_5=I_\theta(\alpha^{(5)}),
\qquad
c=\begin{pmatrix}-1\\1\end{pmatrix},
\qquad
w_5=\frac{v(t_{\mathrm{cert},5})}{J^*(\theta,p,\varepsilon)}.
\]
The matrix \(I_5\) is positive semidefinite, since it is a sum of nonnegative multiples of \(g_Sg_S^\top\). Stationarity and the certificate value give
\[
\mathbf1_2^\top I_5v(t_{\mathrm{cert},5})=0,
\qquad
v(t_{\mathrm{cert},5})^\top I_5v(t_{\mathrm{cert},5})=J^*.
\]
The first identity makes the two coordinates of \(I_5v(t_{\mathrm{cert},5})\) opposite; the second makes its second coordinate \(J^*\). Hence
\[
I_5v(t_{\mathrm{cert},5})=J^*c,
\qquad
I_5w_5=c,
\qquad
c^\top w_5=\frac1{J^*}=V^*.
\]
Every other solution \(w_{\mathrm{alt}}\in\mathbb R^2\) of \(I_5w_{\mathrm{alt}}=c\) has the same contrast value, because
\[
c^\top(w_{\mathrm{alt}}-w_5)
=w_5^\top I_5(w_{\mathrm{alt}}-w_5)=0.
\]
The contrast-variance convention in \cref{obj:thm:five-output-region} therefore gives variance \(V^*>0\), equal to its nonnegative part.

For the associated staircase channel,
\[
(P_\theta Q_{\alpha^{(5)}})(\{S\})
=\alpha^{(5)}_S h_S(\theta).
\]
The positive masses \(h_S(\theta)\) make precisely the five patterns in \(A_5\) active, so
\[
\kappa(Q_{\alpha^{(5)}})=5.
\]
Feasibility supplies unit row sums, and the ray entries \(1,3\) supply the privacy inequalities. Finally, \cref{obj:thm:five-output-region} gives
\[
\kappa(Q)\ge5
\]
for every attaining stationary \(\varepsilon\)-LDP channel on any measurable output space, since \(\iota(\theta)\in\mathcal R_5\).
% lean: CausalSmith.Stat.LdpAteEfficiencySurface.optimalCardinality_of_mem_R5

\item \emph{Obtain the unique three-pattern optimum.}
Fix \(\theta\in U_3\). Its membership in \(\mathcal R_3\) supplies certificate data
\[
\alpha^{(3)}\in\mathcal A_\varepsilon,\qquad
t_{\mathrm{cert},3}\in\mathbb R,\qquad
\eta^{(3)}\in\mathbb R^4,
\]
with support \(A_3\), positive active weights, stationarity, active dual equalities, and strict inactive dual inequalities. The quadratic expansion and normalization calculation used above give
\[
J^*
=\inf_tF_\theta(\alpha^{(3)},t)
=F_\theta(\alpha^{(3)},t_{\mathrm{cert},3})
=\mathbf1_4^\top\eta^{(3)},
\]
and, for every feasible \(\beta\),
\[
F_\theta(\beta,t_{\mathrm{cert},3})
\le\mathbf1_4^\top\eta^{(3)}=J^*.
\]
If \(J^*=\inf_tF_\theta(\beta,t)\), equality follows and
\[
0=\sum_{S\in\mathcal S}\beta_S
 \bigl(b_S^\top\eta^{(3)}-f_S(t_{\mathrm{cert},3})\bigr),
\qquad
\beta_S=0\quad(S\notin A_3).
\]

For any feasible weights supported on \(A_3\), the normalization rows are
\[
\begin{aligned}
3\beta_{S_1}+\beta_{S_6}+\beta_{S_8}&=1,\\
\beta_{S_1}+3\beta_{S_6}+\beta_{S_8}&=1,\\
\beta_{S_1}+\beta_{S_6}+3\beta_{S_8}&=1,
\end{aligned}
\]
with the second row also applying to the other record in \(S_6\). Subtraction and substitution give
\[
\beta_{S_1}=\beta_{S_6}=\beta_{S_8}=\frac15.
\]
These equations apply both to \(\alpha^{(3)}\) and to every feasible optimal \(\beta\). Consequently,
\[
\alpha^{(3)}_S=
\begin{cases}
1/5,&S\in A_3,\\
0,&S\notin A_3,
\end{cases}
\qquad
\beta=\alpha^{(3)}.
\]
% lean: CausalSmith.Stat.LdpAteEfficiencySurface.uniqueOptimalSupport_of_mem_R3

\item \emph{Construct an attaining channel with three active outputs.}
The canonical weights just obtained agree with
\(\alpha^{\mathrm{can}}(\iota(\theta))\). Their stationarity equation is
\[
0=\ell_3(\iota(\theta))+a_3(\iota(\theta))t_{\mathrm{cert},3}.
\]
Since \(a_3(\iota(\theta))>0\),
\[
t_{\mathrm{cert},3}=t_3(\iota(\theta)).
\]
The information matrix is positive semidefinite, and its certifying minimum equals \(J^*>0\). The same two-coordinate calculation therefore gives
\[
I_\theta(\alpha^{(3)})v(t_{\mathrm{cert},3})=J^*c,
\qquad
I_\theta(\alpha^{(3)})
 \frac{v(t_{\mathrm{cert},3})}{J^*}=c,
\qquad
c^\top\frac{v(t_{\mathrm{cert},3})}{J^*}=V^*>0.
\]
Thus the feasible staircase channel attains the required contrast variance. Its output probabilities satisfy
\[
(P_\theta Q_{\alpha^{(3)}})(\{S\})
=\alpha^{(3)}_S h_S(\theta),
\qquad
\kappa(Q_{\alpha^{(3)}})=|A_3|=3.
\]
% lean: CausalSmith.Stat.LdpAteEfficiencySurface.r3Certificate_staircase_attains_three

\item \emph{Prove the universal three-output lower bound.}
Let \(Q\) be any attaining stationary \(\varepsilon\)-LDP channel on a measurable output space \(Z\), and define
\[
I_Q=I_\theta(Q).
\]
Its contrast variance is the finite positive value \(V^*=1/J^*\). Hence the contrast is in the range of \(I_Q\), and we may choose
\[
w_Q\in\mathbb R^2,\qquad
I_Qw_Q=c,\qquad
c^\top w_Q=\frac1{J^*}.
\]
Positive semidefiniteness and Cauchy--Schwarz give
\[
1=(v(t_{\mathrm{cert},3})^\top I_Qw_Q)^2
\le
\bigl(v(t_{\mathrm{cert},3})^\top I_Qv(t_{\mathrm{cert},3})\bigr)
 \bigl(w_Q^\top I_Qw_Q\bigr).
\]
Since \(w_Q^\top I_Qw_Q=c^\top w_Q=1/J^*\), this implies
\[
J^*\le
v(t_{\mathrm{cert},3})^\top I_Qv(t_{\mathrm{cert},3}).
\]

Apply \cref{obj:lem:staircase-refinement}, using the interior parameters and the constant positive privacy budget. It supplies feasible weights \(\bar\alpha\) and a probability kernel \(K\) with
\[
Q=K\circ Q_{\bar\alpha},
\qquad
I_\theta(\bar\alpha)-I_Q\succeq0.
\]
The directional lower bound, refinement inequality, and certificate upper bound form the chain
\[
J^*
\le v(t_{\mathrm{cert},3})^\top I_Qv(t_{\mathrm{cert},3})
\le F_\theta(\bar\alpha,t_{\mathrm{cert},3})
\le J^*.
\]
Thus every inequality is an equality. Equality in the dual bound gives
\[
0=\sum_{S\in\mathcal S}\bar\alpha_S
 \bigl(b_S^\top\eta^{(3)}-f_S(t_{\mathrm{cert},3})\bigr),
\qquad
\bar\alpha_S=0\quad(S\notin A_3).
\]
The normalization equations above then force
\[
\bar\alpha_{S_1}=\bar\alpha_{S_6}=\bar\alpha_{S_8}=\frac15.
\]

Because \(\iota(\theta)\in D_3\) and
\(t_{\mathrm{cert},3}=t_3(\iota(\theta))\), the active projected scores satisfy
\[
\rho_S(\theta,t_{\mathrm{cert},3})
\ne\rho_T(\theta,t_{\mathrm{cert},3})
\qquad(S,T\in A_3,\ S\ne T).
\]
Equality in the refinement direction, together with
\cref{obj:lem:staircase-refinement}, consequently implies
\[
K(S,\cdot)\perp K(T,\cdot)
\qquad(S,T\in A_3,\ S\ne T).
\]

Define the output law by
\[
\nu_Q=P_\theta Q
=\sum_{S\in A_3}\frac{h_S(\theta)}5K(S,\cdot).
\]
In particular, for every measurable output event \(E\) and every \(S\in A_3\),
\[
\nu_Q(E)\ge\frac{h_S(\theta)}5K(S,E),
\qquad
K(S,\cdot)\ll\nu_Q.
\]
If \(Z\) is infinite, the cardinality convention gives
\(\kappa(Q)=+\infty\ge3\).

If \(Z\) is finite, each probability measure \(K(S,\cdot)\), for \(S\in A_3\), satisfies
\[
1=K(S,Z)\le\sum_{z\in Z}K(S,\{z\}).
\]
Thus choose, for every \(S\in A_3\), a point
\[
z_S\in Z,\qquad K(S,\{z_S\})>0.
\]
Absolute continuity transfers positivity to the output law: indeed,
\[
\nu_Q(\{z_S\})=0\Longrightarrow K(S,\{z_S\})=0,
\]
so \(\nu_Q(\{z_S\})>0\).

To prove that these points are distinct, fix distinct \(S,T\in A_3\). Mutual singularity supplies a measurable set \(E_{S,T}\subseteq Z\) with
\[
K(S,E_{S,T})=0,\qquad K(T,Z\setminus E_{S,T})=0.
\]
Suppose that \(z_S=z_T\). If \(z_S\in E_{S,T}\), monotonicity gives
\[
0<K(S,\{z_S\})\le K(S,E_{S,T})=0.
\]
If \(z_S\notin E_{S,T}\), then \(z_T\in Z\setminus E_{S,T}\), and
\[
0<K(T,\{z_T\})\le K(T,Z\setminus E_{S,T})=0.
\]
Both cases yield a contradiction. Hence the three selected points are distinct. Hence
\[
\kappa(Q)
=\#\{z\in Z:\nu_Q(\{z\})\ne0\}
\ge |A_3|=3.
\]
Together with the attaining three-output staircase channel, this proves exact optimal cardinality three throughout \(U_3\), and completes the result.
% lean: CausalSmith.Stat.LdpAteEfficiencySurface.optimalCardinality_of_mem_R3_of_distinct
\end{enumerate}
\end{proof}

\begin{proof}[Proof of \cref{obj:thm:balanced-reduction}]
Throughout, \(p=1/2\). We use the staircase notation of \cref{obj:synth_4,obj:synth_8,obj:def:staircase-saddle}.

\begin{enumerate}
\item Define the binary privacy contraction and the corresponding information value by
\[
\kappa_\varepsilon
:=\frac{e^\varepsilon-1}{e^\varepsilon+1},
\qquad
J_{\mathrm{bin}}
:=\frac{\kappa_\varepsilon^2}
{1-\kappa_\varepsilon^2\tau^2}.
\]
Since \(\varepsilon>0\), we have \(e^\varepsilon>1\), and hence
\[
0<\kappa_\varepsilon<1.
\]
The interior-mean inequalities give
\[
-1<\mu_1-\mu_0<1,
\qquad |\tau|<1.
\]
Consequently,
\[
0\le \kappa_\varepsilon^2\tau^2
\le \kappa_\varepsilon^2<1,
\qquad
1-\kappa_\varepsilon^2\tau^2>0,
\qquad
J_{\mathrm{bin}}>0.
\]
% lean: CausalSmith.Stat.LdpAteEfficiencySurface.sharp_information_pos

\item We first establish positivity of the optimized information. For every staircase pattern,
\[
1\le b_S(j)\le e^\varepsilon,
\qquad
\sum_{j=0}^3\pi_{\theta,j}=1,
\]
so
\[
1\le h_S(\theta)
=\sum_{j=0}^3\pi_{\theta,j}b_S(j)
\le e^\varepsilon.
\]
Define the singleton design by
\[
\alpha^{\mathrm{sing}}_S
:=
\begin{cases}
(e^\varepsilon+3)^{-1},& |S|=1,\\
0,& |S|\ne1,
\end{cases}
\qquad S\in\mathcal S.
\]
For each record \(j\), exactly one singleton pattern contributes
\(e^\varepsilon\) and the other three contribute \(1\). Thus
\[
\sum_{S\in\mathcal S}
\alpha^{\mathrm{sing}}_S b_S(j)
=\frac{e^\varepsilon+3}{e^\varepsilon+3}=1,
\qquad
\alpha^{\mathrm{sing}}\in\mathcal A_\varepsilon.
\]
% lean: CausalSmith.Stat.LdpAteEfficiencySurface.staircaseFeasible_singleton_design

For completeness, its objective attains a minimum. Define
\[
\begin{aligned}
a_{\mathrm{sing}}
&:=2\sum_{S\in\mathcal S}
\alpha^{\mathrm{sing}}_S
\frac{(g_S^\top v(0))(g_{S,0}+g_{S,1})}
{h_S(\theta)},\\
b_{\mathrm{sing}}
&:=\sum_{S\in\mathcal S}
\alpha^{\mathrm{sing}}_S
\frac{(g_{S,0}+g_{S,1})^2}{h_S(\theta)}
\ge0.
\end{aligned}
\]
Expanding \(v(t)=(t,t+1)^\top\) gives
\[
F_\theta(\alpha^{\mathrm{sing}},t)
=
F_\theta(\alpha^{\mathrm{sing}},0)
+a_{\mathrm{sing}}t+b_{\mathrm{sing}}t^2,
\qquad t\in\mathbb R.
\]
If \(b_{\mathrm{sing}}=0\), nonnegativity of the objective forces
\(a_{\mathrm{sing}}=0\): otherwise, at
\[
t_{\mathrm{bad}}
:=-\frac{F_\theta(\alpha^{\mathrm{sing}},0)+1}
{a_{\mathrm{sing}}},
\]
the objective would equal \(-1\). In this case \(t=0\) is a minimizer.
If \(b_{\mathrm{sing}}>0\), define
\[
t_{\mathrm{sing}}
:=-\frac{a_{\mathrm{sing}}}{2b_{\mathrm{sing}}}.
\]
Then
\[
F_\theta(\alpha^{\mathrm{sing}},t)
-F_\theta(\alpha^{\mathrm{sing}},t_{\mathrm{sing}})
=b_{\mathrm{sing}}(t-t_{\mathrm{sing}})^2\ge0.
\]
In the constant case, set \(t_{\mathrm{sing}}:=0\).
% lean: CausalSmith.Stat.LdpAteEfficiencySurface.informationObjective_exists_minimizer

The contributions from \(S=\{0\}\) and \(S=\{3\}\) satisfy
\[
\begin{aligned}
0\le
\frac{(e^\varepsilon-1)^2t_{\mathrm{sing}}^2}
{4(e^\varepsilon+3)h_{\{0\}}(\theta)}
&\le F_\theta(\alpha^{\mathrm{sing}},t_{\mathrm{sing}}),\\
0\le
\frac{(e^\varepsilon-1)^2(t_{\mathrm{sing}}+1)^2}
{4(e^\varepsilon+3)h_{\{3\}}(\theta)}
&\le F_\theta(\alpha^{\mathrm{sing}},t_{\mathrm{sing}}).
\end{aligned}
\]
Both coefficients are strictly positive. A zero minimum would therefore
force \(t_{\mathrm{sing}}=0\) and \(t_{\mathrm{sing}}+1=0\), a contradiction.
By the definition of \(J^*\),
\[
0<
\min_{t\in\mathbb R}F_\theta(\alpha^{\mathrm{sing}},t)
\le J^*(\theta,1/2,\varepsilon).
\]
% lean: CausalSmith.Stat.LdpAteEfficiencySurface.Jstar_pos_interior

\item We next bound the balanced objective. Fix
\(\alpha\in\mathcal A_\varepsilon\), and write its staircase channel as
\[
Q_\alpha(x_j,\{S\})
:=\alpha_S b_S(j),
\qquad j\in\{0,1,2,3\},\quad S\in\mathcal S.
\]
The feasibility equations make its rows probability measures. Since
\(b_S(j)\le e^\varepsilon b_S(j')\), summing the atomic inequalities gives
\[
Q_\alpha(x_j,A)
\le e^\varepsilon Q_\alpha(x_{j'},A),
\qquad A\subseteq\mathcal S.
\]
Define the grouped binary channel by
\[
\begin{aligned}
\overline Q_\alpha(+1,A)
&:=\frac{Q_\alpha(x_0,A)+Q_\alpha(x_3,A)}2,\\
\overline Q_\alpha(-1,A)
&:=\frac{Q_\alpha(x_1,A)+Q_\alpha(x_2,A)}2.
\end{aligned}
\]
Its rows are probability measures. Pairing the original row inequalities
gives
\[
\begin{aligned}
\overline Q_\alpha(+1,A)
&\le \frac{e^\varepsilon}{2}
\{Q_\alpha(x_1,A)+Q_\alpha(x_2,A)\}
=e^\varepsilon\overline Q_\alpha(-1,A),\\
\overline Q_\alpha(-1,A)
&\le e^\varepsilon\overline Q_\alpha(+1,A).
\end{aligned}
\]
The inequalities with identical binary inputs follow from
\(1\le e^\varepsilon\).
% lean: CausalSmith.Stat.LdpAteEfficiencySurface.BalancedBinary.groupedChannel_private

Introduce the reference measure and its nonnegative row densities:
\[
\nu_\alpha
:=\overline Q_\alpha(+1,\cdot)
+\overline Q_\alpha(-1,\cdot),
\qquad
f_{\alpha,\pm}
:=\frac{d\overline Q_\alpha(\pm1,\cdot)}{d\nu_\alpha}.
\]
Define, on the entire output space,
\[
a_\alpha(z)
:=\frac{f_{\alpha,+}(z)+f_{\alpha,-}(z)}2,
\qquad
\delta_\alpha(z)
:=
\begin{cases}
\displaystyle
\frac{f_{\alpha,+}(z)-f_{\alpha,-}(z)}
{f_{\alpha,+}(z)+f_{\alpha,-}(z)},
& f_{\alpha,+}(z)+f_{\alpha,-}(z)>0,\\[6pt]
0,& f_{\alpha,+}(z)+f_{\alpha,-}(z)=0.
\end{cases}
\]
Row normalization yields
\[
\int a_\alpha\,d\nu_\alpha=1,
\qquad
a_\alpha\delta_\alpha
=\frac{f_{\alpha,+}-f_{\alpha,-}}2,
\qquad
\int a_\alpha\delta_\alpha\,d\nu_\alpha=0.
\]
Privacy implies, almost everywhere,
\[
f_{\alpha,+}\le e^\varepsilon f_{\alpha,-},
\qquad
f_{\alpha,-}\le e^\varepsilon f_{\alpha,+}.
\]
At a point where their sum is positive, these inequalities imply
\[
(e^\varepsilon+1)|f_{\alpha,+}-f_{\alpha,-}|
\le
(e^\varepsilon-1)(f_{\alpha,+}+f_{\alpha,-}),
\]
and therefore \(|\delta_\alpha|\le\kappa_\varepsilon\).
Where the sum vanishes, both densities vanish and the same bound follows
from \(\delta_\alpha=0\).
% lean: Causalean.Stat.Privacy.Binary.ae_abs_contrast_le

The grouped mixture density and its derivative with respect to the
binary sign parameter are
\[
m_{\alpha,\tau}
:=\frac{1+\tau}{2}f_{\alpha,+}
+\frac{1-\tau}{2}f_{\alpha,-}
=a_\alpha(1+\tau\delta_\alpha),
\qquad
\partial_\tau m_{\alpha,\tau}=a_\alpha\delta_\alpha.
\]
Thus its Fisher information is
\[
\mathcal I_\alpha
:=\int
\frac{a_\alpha\delta_\alpha^2}
{1+\tau\delta_\alpha}\,d\nu_\alpha.
\]
Zero-density outputs contribute zero. Since
\(|\tau|<1\) and \(|\delta_\alpha|\le\kappa_\varepsilon<1\), the remaining
denominators are positive. The pointwise inequality needed below follows
from the identity
\[
J_{\mathrm{bin}}(1-\tau\delta_\alpha)
-\frac{\delta_\alpha^2}{1+\tau\delta_\alpha}
=
\frac{\kappa_\varepsilon^2-\delta_\alpha^2}
{(1-\kappa_\varepsilon^2\tau^2)(1+\tau\delta_\alpha)}
\ge0.
\]
This identity also applies when \(\tau=0\).
Multiplying by \(a_\alpha\ge0\) and integrating gives
\[
\begin{aligned}
\mathcal I_\alpha
&\le
J_{\mathrm{bin}}
\int(a_\alpha-\tau a_\alpha\delta_\alpha)\,d\nu_\alpha\\
&=J_{\mathrm{bin}}.
\end{aligned}
\]
The bound \(|a_\alpha\delta_\alpha|\le\kappa_\varepsilon a_\alpha\)
ensures the integrability of the displayed upper bound.
% lean: Causalean.Stat.Privacy.Binary.information_le_sharp

We now identify this information with the staircase objective. Balance
gives
\[
\tau=1-2\mu_0,
\qquad
\pi_\theta
=\frac12(1-\mu_0,\mu_0,\mu_0,1-\mu_0)^\top.
\]
Hence, for every \(S\in\mathcal S\),
\[
\begin{aligned}
h_S(\theta)
&=\frac{(1-\mu_0)\{b_S(0)+b_S(3)\}
+\mu_0\{b_S(1)+b_S(2)\}}2,\\
g_S^\top v(-1/2)
&=\frac{b_S(0)+b_S(3)-b_S(1)-b_S(2)}4.
\end{aligned}
\]
Define the grouped atomic probabilities by
\[
q_{\alpha,+}(S)
:=\frac{\alpha_S\{b_S(0)+b_S(3)\}}2,
\qquad
q_{\alpha,-}(S)
:=\frac{\alpha_S\{b_S(1)+b_S(2)\}}2.
\]
The discrete Fisher-information formula therefore gives
\[
\begin{aligned}
\mathcal I_\alpha
&=
\sum_{S\in\mathcal S}
\frac{\bigl\{(q_{\alpha,+}(S)-q_{\alpha,-}(S))/2\bigr\}^2}
{\{(1+\tau)q_{\alpha,+}(S)
+(1-\tau)q_{\alpha,-}(S)\}/2}\\
&=
\sum_{S\in\mathcal S}
\alpha_S\frac{(g_S^\top v(-1/2))^2}{h_S(\theta)}
=F_\theta(\alpha,-1/2).
\end{aligned}
\]
Here a summand with \(\alpha_S=0\) is defined as zero; both grouped
probabilities then vanish. When \(\alpha_S>0\), its mixture probability
is \(\alpha_Sh_S(\theta)>0\), and the displayed equality follows by
cancellation. We have proved
\[
F_\theta(\alpha,-1/2)\le J_{\mathrm{bin}}
\qquad\text{for every }\alpha\in\mathcal A_\varepsilon.
\]
% lean: CausalSmith.Stat.LdpAteEfficiencySurface.BalancedBinary.informationObjective_balanced_le_sharp

\item The feasible set is nonempty by the singleton construction above.
It is closed, and its row constraints imply
\[
0\le\alpha_S\le\sum_{S'\in\mathcal S}\alpha_{S'}\le1.
\]
Thus it is compact. Since the objective at a fixed profiling coordinate
is continuous and linear in the weights, choose
\[
\alpha^{\mathrm{env}}
\in\operatorname*{arg\,max}_{\alpha\in\mathcal A_\varepsilon}
F_\theta(\alpha,-1/2).
\]
For each feasible \(\alpha\), its profiled minimum is at most its value
at \(-1/2\). Consequently,
\[
\begin{aligned}
J^*(\theta,1/2,\varepsilon)
&\le
\max_{\alpha\in\mathcal A_\varepsilon}
F_\theta(\alpha,-1/2)\\
&=F_\theta(\alpha^{\mathrm{env}},-1/2)
\le J_{\mathrm{bin}}.
\end{aligned}
\]
Both information values are positive, so taking reciprocals yields
\[
J_{\mathrm{bin}}^{-1}
\le V^*(\theta,1/2,\varepsilon).
\]
% lean: CausalSmith.Stat.LdpAteEfficiencySurface.BalancedBinary.Jstar_balanced_le_sharp

\item We establish the channel properties needed for the matching
bound. For the ordered record alphabet in \cref{obj:synth_8}, define
\[
(\sigma_0,\sigma_1,\sigma_2,\sigma_3):=(1,-1,-1,1).
\]
These are the values of \((2W_i-1)(2Y_i-1)\).
The randomized-response channel has atomic probabilities
\[
Q(x_j,\{z\})
=
\begin{cases}
\displaystyle\frac{e^\varepsilon}{e^\varepsilon+1},
&z=\sigma_j,\\[6pt]
\displaystyle\frac1{e^\varepsilon+1},
&z=-\sigma_j,
\end{cases}
\qquad z\in\{-1,1\}.
\]
They are positive and satisfy
\[
\sum_{z\in\{-1,1\}}Q(x_j,\{z\})
=\frac{e^\varepsilon+1}{e^\varepsilon+1}=1.
\]
Thus \(Q\) is a Markov channel, and using this same kernel at each
release makes it stationary. For every pair of rows and every output,
\[
\frac{Q(x_j,\{z\})}{Q(x_{j'},\{z\})}
\in\{e^{-\varepsilon},1,e^\varepsilon\},
\]
so
\[
Q(x_j,\{z\})\le e^\varepsilon Q(x_{j'},\{z\}).
\]
Summing over the outputs in any measurable event gives
\[
Q(x_j,A)\le e^\varepsilon Q(x_{j'},A).
\]
% lean: CausalSmith.Stat.LdpAteEfficiencySurface.balancedRR_stationary

\item Group this channel in the same way:
\[
\overline Q_{\mathrm{RR}}(+1,\cdot)
:=\frac{Q(x_0,\cdot)+Q(x_3,\cdot)}2=Q(x_0,\cdot),
\qquad
\overline Q_{\mathrm{RR}}(-1,\cdot)
:=\frac{Q(x_1,\cdot)+Q(x_2,\cdot)}2=Q(x_1,\cdot).
\]
The equalities hold because the paired records have the same input
sign. Its binary-input rows are therefore ordinary randomized response.
Mixing them with probabilities \((1+\tau)/2\) and \((1-\tau)/2\) gives
\[
P_\theta(Z=z)=\frac{1+z\kappa_\varepsilon\tau}{2},
\qquad
\partial_\tau P_\theta(Z=z)=\frac{z\kappa_\varepsilon}{2},
\qquad z\in\{-1,1\}.
\]
The two mixture probabilities are positive. Its scalar Fisher
information is consequently
\[
\begin{aligned}
\mathcal I_{\mathrm{RR}}
&:=
\sum_{z\in\{-1,1\}}
\frac{\{\partial_\tau P_\theta(Z=z)\}^2}{P_\theta(Z=z)}\\
&=
\frac{\kappa_\varepsilon^2}{2(1+\kappa_\varepsilon\tau)}
+\frac{\kappa_\varepsilon^2}{2(1-\kappa_\varepsilon\tau)}
=J_{\mathrm{bin}}.
\end{aligned}
\]
% lean: CausalSmith.Stat.LdpAteEfficiencySurface.grouped_balancedRR_information_eq

To relate this scalar information to the full information matrix, define
\[
D_{\mathrm{RR}}
:=\sum_{j=0}^3Q(x_j,\cdot),
\qquad
f_{\mathrm{RR},j}
:=\frac{dQ(x_j,\cdot)}{dD_{\mathrm{RR}}},
\qquad
m_{\mathrm{RR},\theta}
:=\sum_{j=0}^3\pi_{\theta,j}f_{\mathrm{RR},j}.
\]
The two-coordinate density derivative at fixed assignment probability is
\[
\dot m_{\mathrm{RR},\theta}
:=
\begin{pmatrix}
(f_{\mathrm{RR},1}-f_{\mathrm{RR},0})/2\\
(f_{\mathrm{RR},3}-f_{\mathrm{RR},2})/2
\end{pmatrix}.
\]
Under balance,
\[
m_{\mathrm{RR},\theta}
=
\frac{(1-\mu_0)(f_{\mathrm{RR},0}+f_{\mathrm{RR},3})
+\mu_0(f_{\mathrm{RR},1}+f_{\mathrm{RR},2})}{2},
\]
and
\[
v(-1/2)^\top\dot m_{\mathrm{RR},\theta}
=
\frac{f_{\mathrm{RR},0}+f_{\mathrm{RR},3}
-f_{\mathrm{RR},1}-f_{\mathrm{RR},2}}4.
\]
The grouped reference measure and its row densities are
\[
\nu_{\mathrm{RR}}:=\frac12D_{\mathrm{RR}},
\qquad
\overline f_{\mathrm{RR},+}
:=f_{\mathrm{RR},0}+f_{\mathrm{RR},3},
\qquad
\overline f_{\mathrm{RR},-}
:=f_{\mathrm{RR},1}+f_{\mathrm{RR},2}.
\]
Its mixture density satisfies
\[
\overline m_{\mathrm{RR},\tau}
:=\frac{1+\tau}{2}\overline f_{\mathrm{RR},+}
+\frac{1-\tau}{2}\overline f_{\mathrm{RR},-}
=2m_{\mathrm{RR},\theta}.
\]
Therefore, writing
\[
I_{\mathrm{RR}}:=I_\theta(Q),
\]
we obtain
\[
\begin{aligned}
\mathcal I_{\mathrm{RR}}
&=
\int
\frac{\{(\overline f_{\mathrm{RR},+}
-\overline f_{\mathrm{RR},-})/2\}^2}
{2m_{\mathrm{RR},\theta}}\,d\nu_{\mathrm{RR}}\\
&=
\int
\frac{\{v(-1/2)^\top\dot m_{\mathrm{RR},\theta}\}^2}
{m_{\mathrm{RR},\theta}}\,dD_{\mathrm{RR}}\\
&=v(-1/2)^\top I_{\mathrm{RR}}v(-1/2)
=J_{\mathrm{bin}}.
\end{aligned}
\]
The integrands are defined as zero at zero-density outputs. Indeed, if
the sum of the four nonnegative row densities vanishes, all four vanish.
If that sum is positive and
\(f_{\mathrm{RR},0}+f_{\mathrm{RR},3}=0\), then
\[
m_{\mathrm{RR},\theta}
=\frac{\mu_0}{2}
(f_{\mathrm{RR},1}+f_{\mathrm{RR},2})>0.
\]
If \(f_{\mathrm{RR},0}+f_{\mathrm{RR},3}>0\), then
\[
m_{\mathrm{RR},\theta}
\ge\frac{1-\mu_0}{2}
(f_{\mathrm{RR},0}+f_{\mathrm{RR},3})>0.
\]
Thus the change of reference measure and the divisions above cover all
outputs.
% lean: CausalSmith.Stat.LdpAteEfficiencySurface.BalancedBinary.grouped_information_eq_channelQuadratic

The paired-row identities also give, almost everywhere,
\[
f_{\mathrm{RR},0}=f_{\mathrm{RR},3},
\qquad
f_{\mathrm{RR},1}=f_{\mathrm{RR},2}.
\]
It follows that, for every \(t\in\mathbb R\),
\[
v(t)^\top\dot m_{\mathrm{RR},\theta}
=
\frac{t(f_{\mathrm{RR},1}-f_{\mathrm{RR},0})
+(t+1)(f_{\mathrm{RR},3}-f_{\mathrm{RR},2})}{2}
=
v(-1/2)^\top\dot m_{\mathrm{RR},\theta}.
\]
Integrating the squared directional derivative therefore gives
\[
v(t)^\top I_{\mathrm{RR}}v(t)=J_{\mathrm{bin}},
\qquad t\in\mathbb R.
\]
% lean: CausalSmith.Stat.LdpAteEfficiencySurface.balancedRR_quadratic_eq

\item We now evaluate the contrast variance using this positive
minimizing quadratic value. The matrix \(I_{\mathrm{RR}}\) is symmetric,
and for every \(\xi_{\mathrm{dir}}\in\mathbb R^2\),
\[
\xi_{\mathrm{dir}}^\top I_{\mathrm{RR}}\xi_{\mathrm{dir}}
=
\int
\frac{(\xi_{\mathrm{dir}}^\top\dot m_{\mathrm{RR},\theta})^2}
{m_{\mathrm{RR},\theta}}\,dD_{\mathrm{RR}}
\ge0.
\]
Thus it is positive semidefinite.
% lean: CausalSmith.Stat.LdpAteEfficiencySurface.channelFisherInfo_posSemidef

Define the perturbation coefficients by
\[
a_{\mathrm{geom}}
:=\sum_{k,\ell=0}^1(I_{\mathrm{RR}})_{k\ell},
\qquad
b_{\mathrm{geom}}
:=2\{(I_{\mathrm{RR}}v(-1/2))_0
+(I_{\mathrm{RR}}v(-1/2))_1\}.
\]
Symmetry and expansion give, for every \(\eta\in\mathbb R\),
\[
\begin{aligned}
0
&\le
v(-1/2+\eta)^\top I_{\mathrm{RR}}v(-1/2+\eta)
-v(-1/2)^\top I_{\mathrm{RR}}v(-1/2)\\
&=\eta b_{\mathrm{geom}}+\eta^2a_{\mathrm{geom}}.
\end{aligned}
\]
If \(b_{\mathrm{geom}}\ne0\), choose
\[
\eta_{\mathrm{test}}
:=-\frac{b_{\mathrm{geom}}}{2(|a_{\mathrm{geom}}|+1)}.
\]
Using \(a_{\mathrm{geom}}\le|a_{\mathrm{geom}}|\), the last expression
would satisfy
\[
\begin{aligned}
\eta_{\mathrm{test}}b_{\mathrm{geom}}
+\eta_{\mathrm{test}}^2a_{\mathrm{geom}}
&=
\frac{b_{\mathrm{geom}}^2
\{a_{\mathrm{geom}}-2(|a_{\mathrm{geom}}|+1)\}}
{4(|a_{\mathrm{geom}}|+1)^2}\\
&\le-\frac{b_{\mathrm{geom}}^2}
{4(|a_{\mathrm{geom}}|+1)}<0.
\end{aligned}
\]
Hence \(b_{\mathrm{geom}}=0\), so the coordinates of
\(I_{\mathrm{RR}}v(-1/2)\) sum to zero. Since
\[
J_{\mathrm{bin}}
=v(-1/2)^\top I_{\mathrm{RR}}v(-1/2)
=(I_{\mathrm{RR}}v(-1/2))_1,
\]
we obtain
\[
I_{\mathrm{RR}}v(-1/2)
=J_{\mathrm{bin}}(-1,1)^\top
=J_{\mathrm{bin}}c.
\]
Define
\[
w_{\mathrm{sol}}
:=\frac{v(-1/2)}{J_{\mathrm{bin}}}\in\mathbb R^2.
\]
Then
\[
I_{\mathrm{RR}}w_{\mathrm{sol}}=c,
\qquad
c^\top w_{\mathrm{sol}}=J_{\mathrm{bin}}^{-1},
\qquad
c\in\operatorname{range}(I_{\mathrm{RR}}).
\]
For any other solution
\(\widetilde w_{\mathrm{sol}}\in\mathbb R^2\) of
\(I_{\mathrm{RR}}\widetilde w_{\mathrm{sol}}=c\), symmetry gives
\[
c^\top\widetilde w_{\mathrm{sol}}
=(I_{\mathrm{RR}}w_{\mathrm{sol}})^\top
\widetilde w_{\mathrm{sol}}
=w_{\mathrm{sol}}^\top
I_{\mathrm{RR}}\widetilde w_{\mathrm{sol}}
=w_{\mathrm{sol}}^\top c
=J_{\mathrm{bin}}^{-1}.
\]
Thus the solution-value infimum in the contrast-variance convention is
\[
c^\top I_{\mathrm{RR}}^+c
=
\inf_{\{\widetilde w_{\mathrm{sol}}:\,
I_{\mathrm{RR}}\widetilde w_{\mathrm{sol}}=c\}}
c^\top\widetilde w_{\mathrm{sol}}
=J_{\mathrm{bin}}^{-1}.
\]
% lean: CausalSmith.Stat.LdpAteEfficiencySurface.contrastVariance_eq_reciprocal_of_minimizer
% lean: CausalSmith.Stat.LdpAteEfficiencySurface.balancedRR_variance_eq_sharp

\item Apply \cref{obj:thm:finite-oracle} to \(Q\), with the constant privacy
sequence
\[
\varepsilon_n:=\varepsilon,
\qquad n\in\mathbb N.
\]
Its hypotheses hold because \(0<1/2<1\), the means are interior, the
budget is positive, and \(Q\) satisfies the stationary privacy
inequalities established above. The oracle bound gives
\[
V^*(\theta,1/2,\varepsilon)
\le c^\top I_{\mathrm{RR}}^+c
=J_{\mathrm{bin}}^{-1}.
\]
Together with the reciprocal lower bound, this proves
\[
V^*(\theta,1/2,\varepsilon)
=J_{\mathrm{bin}}^{-1}.
\]
% lean: CausalSmith.Stat.LdpAteEfficiencySurface.Vstar_balanced_eq_sharp_reciprocal

\item Since \(\kappa_\varepsilon>0\) and
\(1-\kappa_\varepsilon^2\tau^2>0\), ordinary fraction algebra gives
\[
J_{\mathrm{bin}}^{-1}
=
\frac{1-\kappa_\varepsilon^2\tau^2}{\kappa_\varepsilon^2}
=\kappa_\varepsilon^{-2}-\tau^2.
\]
Moreover,
\[
\sinh(\varepsilon/2)>0,
\qquad
e^\varepsilon=(e^{\varepsilon/2})^2.
\]
Expanding the hyperbolic functions and multiplying numerator and
denominator by \(e^{\varepsilon/2}\) yields
\[
\frac{\cosh(\varepsilon/2)}{\sinh(\varepsilon/2)}
=
\frac{e^{\varepsilon/2}+e^{-\varepsilon/2}}
{e^{\varepsilon/2}-e^{-\varepsilon/2}}
=
\frac{e^\varepsilon+1}{e^\varepsilon-1}
=\kappa_\varepsilon^{-1}.
\]
Consequently,
\[
V^*(\theta,1/2,\varepsilon)
=
\left(\frac{\cosh(\varepsilon/2)}{\sinh(\varepsilon/2)}\right)^2
-\tau^2.
\]
% lean: CausalSmith.Stat.LdpAteEfficiencySurface.contraction_inv_eq_coth

\item The row normalization and event inequalities above establish the
claimed stationary Markov privacy property. The constructed information
solution establishes estimability, and the two variance identities give
\[
c\in\operatorname{range}(I_\theta(Q)),
\qquad
c^\top I_\theta(Q)^+c
=J_{\mathrm{bin}}^{-1}
=V^*(\theta,1/2,\varepsilon).
\]
% lean: CausalSmith.Stat.LdpAteEfficiencySurface.balanced_reduction
\end{enumerate}
\end{proof}

\begin{proof}[Proof of \cref{obj:thm:laplace-comparison}]
Throughout, reciprocals at zero are interpreted as zero, empty sums as zero, and the real square root as zero on negative arguments.

\begin{enumerate}
\item At the balanced parameter point, the arm means are interior, their sum is one, and the contrast is zero. Thus \cref{obj:thm:balanced-reduction}, with privacy parameter one, gives
\[
V^*((1/2,1/2)^\top,1/2,1)
=
\left(\frac{\cosh(1/2)}{\sinh(1/2)}\right)^2.
\]
We retain this identity for the final comparison.
% lean: CausalSmith.Stat.LdpAteEfficiencySurface.balanced_reduction

\item The specified public experiment and its private-record representation are retained directly from the hypotheses:
\[
\mathcal L(X_1,\ldots,X_n)=P_\theta^{\otimes n},
\qquad n\ge1,
\]
where
\[
\bigl(P_\theta(\{x_0\}),P_\theta(\{x_1\}),
P_\theta(\{x_2\}),P_\theta(\{x_3\})\bigr)
=
\bigl(q(1-\mu_0),q\mu_0,p(1-\mu_1),p\mu_1\bigr).
\]
These are the public product experiments in \cref{obj:def:binary-trial-model}.
% lean: CausalSmith.Stat.LdpAteEfficiencySurface.laplace_comparison

\item Define the trial contribution and the Laplace scale by
\[
A_{\mathrm{trial},i}
=
\frac{W_iY_i}{p}-\frac{(1-W_i)Y_i}{q},
\qquad
b_{\mathrm{Lap}}=\frac{\Delta_A}{\varepsilon}>0.
\]
Binary assignment and outcome consistency imply, almost surely,
\[
A_{\mathrm{trial},i}
=
\begin{cases}
0,&X_i=x_0,\\
-q^{-1},&X_i=x_1,\\
0,&X_i=x_2,\\
p^{-1},&X_i=x_3.
\end{cases}
\]
Summing against the four record probabilities therefore gives
\[
\mathbb E[A_{\mathrm{trial},i}]
=
-q^{-1}(q\mu_0)+p^{-1}(p\mu_1)=\tau,
\]
\[
\mathbb E[A_{\mathrm{trial},i}^2]
=
q^{-2}(q\mu_0)+p^{-2}(p\mu_1)
=
\frac{\mu_0}{q}+\frac{\mu_1}{p}.
\]
The trial contribution is bounded, so both moments are finite.

The Laplace density is
\[
f_{\mathrm{Lap}}(\ell)
=
\frac{1}{2b_{\mathrm{Lap}}}
\exp\!\left(-\frac{|\ell|}{b_{\mathrm{Lap}}}\right),
\qquad \ell\in\mathbb R.
\]
Its symmetry gives \(\mathbb E[L_i]=0\). The gamma integral at shape three gives
\[
\int_0^\infty
\ell^2e^{-\ell/b_{\mathrm{Lap}}}\,d\ell
=
2b_{\mathrm{Lap}}^3,
\]
and hence
\[
\mathbb E[L_i^2]
=
\frac{1}{b_{\mathrm{Lap}}}
\int_0^\infty
\ell^2e^{-\ell/b_{\mathrm{Lap}}}\,d\ell
=
2b_{\mathrm{Lap}}^2
=
\frac{2\Delta_A^2}{\varepsilon^2}.
\]
Since \(A_{\mathrm{trial},i}\) is a function of \(X_i\), its independence from \(L_i\) yields
\[
\mathbb E[L_iA_{\mathrm{trial},i}]
=
\mathbb E[L_i]\mathbb E[A_{\mathrm{trial},i}]
=0.
\]
For the releases of \cref{obj:def:oa-release}, it follows that
\[
\mathbb E[\widetilde A_i]=\tau,
\qquad
\mathbb E[\widetilde A_i^2]
=
\frac{\mu_1}{p}+\frac{\mu_0}{q}
+\frac{2\Delta_A^2}{\varepsilon^2}.
\]
Define their centered versions by
\[
\xi_{\mathrm{OA},i}=\widetilde A_i-\tau.
\]
Expanding the centered square gives
\[
\mathbb E[\xi_{\mathrm{OA},i}]=0,
\qquad
\mathbb E[\xi_{\mathrm{OA},i}^2]
=
\mathbb E[\widetilde A_i^2]
-2\tau\mathbb E[\widetilde A_i]+\tau^2
=
V_{\mathrm{OA}}(\theta,p,\varepsilon)<\infty.
\]

The pairs \((X_i,L_i)\) are independent across subjects. Each pair has the same product law, because the record law and the noise law are common across subjects and its two coordinates are independent. Applying the release function to each pair shows that the centered releases are independent and identically distributed. The iid central limit theorem therefore gives
\[
\frac{1}{\sqrt n}\sum_{i=1}^n\xi_{\mathrm{OA},i}
\rightsquigarrow
N\!\left(0,V_{\mathrm{OA}}(\theta,p,\varepsilon)\right).
\]
For every \(n\ge1\),
\[
\frac{1}{\sqrt n}\sum_{i=1}^n\xi_{\mathrm{OA},i}
=
\sqrt n(\widehat\tau_{\mathrm{OA}}-\tau).
\]
At \(n=0\), both normalized expressions are zero under the stated conventions. This proves the asserted convergence in law.
% lean: CausalSmith.Stat.LdpAteEfficiencySurface.oaRelease_clt

\item Use the staircase notation of \cref{obj:synth_4}. Interior assignment and means give nonnegative record probabilities, and every component of \(b_S\) is nonnegative. Thus
\[
h_S(\theta)
=
\sum_{j=0}^3\pi_{\theta,j}(b_S)_j\ge0,
\]
and, for every feasible \(\alpha\) and every \(t\in\mathbb R\),
\[
F_\theta(\alpha,t)
=
\sum_{S\in\mathcal S}
\alpha_S\frac{(g_S^\top v(t))^2}{h_S(\theta)}
\ge0.
\]
Taking the profiled minimum and optimized maximum in
\cref{obj:def:staircase-saddle} gives
\[
J^*(\theta,p,\varepsilon)\ge0,
\qquad
V^*(\theta,p,\varepsilon)=J^*(\theta,p,\varepsilon)^{-1}\ge0.
\]

For every positive integer \(n\), factoring out \(\sqrt n\) from both square roots and canceling the positive common factor \(z_{1-\gamma/2}/\sqrt n\) gives
\[
\begin{aligned}
\frac{z_{1-\gamma/2}\sqrt{V_{\mathrm{OA}}(\theta,p,\varepsilon)/n}}
     {z_{1-\gamma/2}\sqrt{V^*(\theta,p,\varepsilon)/n}}
&=
\frac{\sqrt{V_{\mathrm{OA}}(\theta,p,\varepsilon)}}
     {\sqrt{V^*(\theta,p,\varepsilon)}}\\
&=
\sqrt{\frac{V_{\mathrm{OA}}(\theta,p,\varepsilon)}
           {V^*(\theta,p,\varepsilon)}}.
\end{aligned}
\]
The identities also respect the zero-reciprocal convention. The sequence is consequently equal to its claimed limit for every \(n\ge1\).
% lean: CausalSmith.Stat.LdpAteEfficiencySurface.Vstar_nonneg
% lean: CausalSmith.Stat.LdpAteEfficiencySurface.laplace_comparison

\item Apply the weak law to the centered releases and their squares. The moment identities above yield
\[
\frac1n\sum_{i=1}^n\xi_{\mathrm{OA},i}
\xrightarrow{\Pr}0,
\qquad
\frac1n\sum_{i=1}^n\xi_{\mathrm{OA},i}^2
\xrightarrow{\Pr}V_{\mathrm{OA}}(\theta,p,\varepsilon).
\]
For \(n\ge1\), centering the empirical variance at \(\tau\) gives the exact identity
\[
\begin{aligned}
\widehat V_{\mathrm{OA}}
&=
\frac1n\sum_{i=1}^n
(\widetilde A_i-\widehat\tau_{\mathrm{OA}})^2\\
&=
\frac1n\sum_{i=1}^n\xi_{\mathrm{OA},i}^2
-
\left(\frac1n\sum_{i=1}^n\xi_{\mathrm{OA},i}\right)^2\\
&=
\frac1n\sum_{i=1}^n\widetilde A_i^2
-\widehat\tau_{\mathrm{OA}}^2.
\end{aligned}
\]
At \(n=0\), each empirical expression is zero. Continuity of squaring and subtraction proves variance consistency. In particular, for every \(\delta>0\),
\[
0\le
\Pr\!\left(
|\widehat V_{\mathrm{OA}}-V_{\mathrm{OA}}|>\delta
\right)
\le
\Pr\!\left(
|\widehat V_{\mathrm{OA}}-V_{\mathrm{OA}}|\ge\delta
\right)
\longrightarrow0.
\]
% lean: CausalSmith.Stat.LdpAteEfficiencySurface.oaVariance_consistency

\item Choose, for every \(n\in\mathbb N\),
\[
m_n=\lfloor\sqrt n\rfloor.
\]
For every integer \(r_{\mathrm{pilot}}\ge0\),
\[
n\ge r_{\mathrm{pilot}}^2
\quad\Longrightarrow\quad
m_n\ge r_{\mathrm{pilot}},
\]
so \(m_n\to\infty\). Moreover, for \(n\ge1\),
\[
0\le\frac{m_n}{n}
\le\frac{\sqrt n}{n}
=\frac1{\sqrt n}\longrightarrow0.
\]
For \(n\ge2\),
\[
1\le m_n\le\sqrt n<n.
\]
The remaining values are \(m_0=0\) and \(m_1=1\). Thus the schedule satisfies
\cref{obj:ass:pilot-diverges,obj:ass:pilot-sublinear} over their full stated ranges.

Take the privacy sequence
\[
\varepsilon_n=\varepsilon,\qquad n\in\mathbb N.
\]
It satisfies \cref{obj:ass:fixed-privacy}. The fixed measurable strong-saddle selector, the interior parameters, the chosen schedule, and \(0<\gamma<1\) meet the hypotheses of \cref{obj:thm:pilot-attainment}. Its variance-consistency conclusion supplies, under the transcript law of the selected procedure,
\[
\Pr_\theta\!\left(
|\widehat V^*-V^*(\theta,p,\varepsilon)|>\delta
\right)\longrightarrow0
\qquad(\delta>0).
\]
% lean: CausalSmith.Stat.LdpAteEfficiencySurface.natSqrt_pilotDiverges
% lean: CausalSmith.Stat.LdpAteEfficiencySurface.natSqrt_pilotSublinear
% lean: CausalSmith.Stat.LdpAteEfficiencySurface.pilot_attainment

\item Fix any sequence of couplings \((\nu_n)_{n\ge0}\) with the specified trial and transcript marginals. Each \(\nu_n\) is a probability measure because its first marginal is the trial probability law. For measurable events \(E\) in the corresponding product space, write
\[
\Pr_{\nu_n}(E)=\nu_n(E).
\]
Regard \(\widehat V_{\mathrm{OA}}\) and \(\widehat V^*\) as functions of the first and second components, respectively.

We first establish strict positivity of the limiting denominator. In the notation of \cref{obj:synth_4},
\[
d=e^\varepsilon-1>0,
\qquad
h_S(\theta)
=
\sum_{j=0}^3\pi_{\theta,j}(b_S)_j>0,
\]
because all record probabilities are positive and all components of \(b_S\) are at least one. Define the singleton weights, for every \(S\in\mathcal S\), by
\[
\alpha_{\mathrm{RR},S}
=
\begin{cases}
(e^\varepsilon+3)^{-1},&|S|=1,\\
0,&|S|\ne1.
\end{cases}
\]
They are nonnegative, and for every record coordinate \(j\),
\[
\sum_{S\in\mathcal S}
\alpha_{\mathrm{RR},S}(b_S)_j
=
\frac{e^\varepsilon+3}{e^\varepsilon+3}=1.
\]
Hence \(\alpha_{\mathrm{RR}}\in\mathcal A_\varepsilon\).

To justify attainment of the profiling minimum, define, for any feasible \(\alpha\),
\[
\lambda_{\mathrm{lin}}(\alpha)
=
\sum_{S\in\mathcal S}
\frac{2\alpha_S(g_S^\top v(0))
\bigl((g_S)_0+(g_S)_1\bigr)}{h_S(\theta)},
\]
\[
\lambda_{\mathrm{quad}}(\alpha)
=
\sum_{S\in\mathcal S}
\frac{\alpha_S\bigl((g_S)_0+(g_S)_1\bigr)^2}
     {h_S(\theta)}
\ge0.
\]
Expanding \(v(t)=(t,t+1)^\top\) gives, for every real \(t\),
\[
F_\theta(\alpha,t)
=
F_\theta(\alpha,0)
+t\lambda_{\mathrm{lin}}(\alpha)
+t^2\lambda_{\mathrm{quad}}(\alpha).
\]
If \(\lambda_{\mathrm{quad}}(\alpha)=0\) and
\(\lambda_{\mathrm{lin}}(\alpha)\ne0\), substitution of
\[
t=-\frac{F_\theta(\alpha,0)+1}
        {\lambda_{\mathrm{lin}}(\alpha)}
\]
would give \(F_\theta(\alpha,t)=-1\), contradicting nonnegativity. Thus in this case the linear coefficient is also zero and \(t=0\) is a minimizer. If \(\lambda_{\mathrm{quad}}(\alpha)>0\), define
\[
t_{\min}(\alpha)
=
-\frac{\lambda_{\mathrm{lin}}(\alpha)}
       {2\lambda_{\mathrm{quad}}(\alpha)}.
\]
Then
\[
F_\theta(\alpha,t)-F_\theta(\alpha,t_{\min}(\alpha))
=
\lambda_{\mathrm{quad}}(\alpha)
\bigl(t-t_{\min}(\alpha)\bigr)^2\ge0.
\]
In both cases a minimizing real coordinate exists. Choose such a coordinate
\(t_{\mathrm{RR}}\) for \(\alpha_{\mathrm{RR}}\).

All summands in \(F_\theta(\alpha_{\mathrm{RR}},t_{\mathrm{RR}})\) are nonnegative. The singleton patterns \(\{0\}\) and \(\{3\}\) therefore give
\[
0\le
\frac{\alpha_{\mathrm{RR},\{0\}}(dq)^2}
     {h_{\{0\}}(\theta)}\,t_{\mathrm{RR}}^2
\le F_\theta(\alpha_{\mathrm{RR}},t_{\mathrm{RR}}),
\]
\[
0\le
\frac{\alpha_{\mathrm{RR},\{3\}}(dp)^2}
     {h_{\{3\}}(\theta)}\,(t_{\mathrm{RR}}+1)^2
\le F_\theta(\alpha_{\mathrm{RR}},t_{\mathrm{RR}}).
\]
Both coefficients are strictly positive. If the objective value were zero, the first inequality would force \(t_{\mathrm{RR}}=0\), while the second would force \(t_{\mathrm{RR}}+1=0\). Consequently,
\[
0<
F_\theta(\alpha_{\mathrm{RR}},t_{\mathrm{RR}})
=
\min_{t\in\mathbb R}F_\theta(\alpha_{\mathrm{RR}},t)
\le J^*(\theta,p,\varepsilon).
\]
The final inequality follows from the optimized maximum in
\cref{obj:def:staircase-saddle}. Taking the reciprocal proves
\[
V^*(\theta,p,\varepsilon)>0.
\]
% lean: CausalSmith.Stat.LdpAteEfficiencySurface.Jstar_pos_interior

\item The formula in \cref{obj:synth_1} expresses \(\widehat V^*\) as a nonnegative reciprocal times a sum of squares. Thus, for every sample size and every transcript,
\[
\widehat V^*\ge0.
\]
When \(m_n\ge n\), the main-sample sum is empty and its reciprocal normalization is zero, so \(\widehat V^*=0\).

For real numerator and denominator inputs, define the total ratio map
\[
f_{\mathrm{ratio}}(r_{\mathrm{num}},r_{\mathrm{den}})
=
\begin{cases}
\sqrt{r_{\mathrm{num}}/r_{\mathrm{den}}},
&r_{\mathrm{den}}>0,\\
0,&r_{\mathrm{den}}\le0,
\end{cases}
\qquad
(r_{\mathrm{num}},r_{\mathrm{den}})\in\mathbb R^2.
\]
On the half-plane \(r_{\mathrm{den}}\ge0\), this agrees with the displayed square-root ratio: at a zero denominator both expressions are zero. Since \(V^*>0\), the map agrees with the continuous square-root-of-quotient map in a neighborhood of \((V_{\mathrm{OA}},V^*)\), and is therefore continuous at that pair.

The marginal identities and the union bound give, for every \(\eta>0\),
\[
\begin{aligned}
&\Pr_{\nu_n}\!\left(
\max\{|\widehat V_{\mathrm{OA}}-V_{\mathrm{OA}}|,
       |\widehat V^*-V^*|\}\ge\eta
\right)\\
&\qquad\le
\Pr\!\left(|\widehat V_{\mathrm{OA}}-V_{\mathrm{OA}}|\ge\eta\right)
+
\Pr_\theta\!\left(|\widehat V^*-V^*|\ge\eta\right)\\
&\qquad\le
\Pr\!\left(|\widehat V_{\mathrm{OA}}-V_{\mathrm{OA}}|>\eta/2\right)
+
\Pr_\theta\!\left(|\widehat V^*-V^*|>\eta/2\right)
\longrightarrow0.
\end{aligned}
\]
Hence the pair of variance estimates converges in probability under every prescribed coupling.

For any \(\delta>0\), continuity supplies an \(\eta>0\) such that
\[
\max\{|r_{\mathrm{num}}-V_{\mathrm{OA}}|,
       |r_{\mathrm{den}}-V^*|\}<\eta
\quad\Longrightarrow\quad
\left|
f_{\mathrm{ratio}}(r_{\mathrm{num}},r_{\mathrm{den}})
-\sqrt{V_{\mathrm{OA}}/V^*}
\right|<\delta.
\]
Using the pointwise nonnegativity of \(\widehat V^*\), we obtain
\[
\begin{aligned}
&\Pr_{\nu_n}\!\left(
\left|
\sqrt{\widehat V_{\mathrm{OA}}/\widehat V^*}
-\sqrt{V_{\mathrm{OA}}/V^*}
\right|>\delta
\right)\\
&\qquad\le
\Pr_{\nu_n}\!\left(
\max\{|\widehat V_{\mathrm{OA}}-V_{\mathrm{OA}}|,
       |\widehat V^*-V^*|\}\ge\eta
\right)
\longrightarrow0.
\end{aligned}
\]
This proves the empirical ratio assertion for the exhibited schedule and every sequence of the stated couplings.
% lean: CausalSmith.Stat.LdpAteEfficiencySurface.VhatStar_nonneg
% lean: Causalean.Stat.Modes.coupled_sqrtRatio_tendsto_strict_abs

\item Finally, at \(p=q=1/2\), \(\mu_0=\mu_1=1/2\), and \(\varepsilon=1\),
\[
\tau=0,\qquad \Delta_A=4,
\]
so direct substitution gives
\[
V_{\mathrm{OA}}
=
\frac{1/2}{1/2}+\frac{1/2}{1/2}
+2\cdot4^2
=34.
\]
To compare this with the balanced value, use
\[
\sinh(1/2)>1/2>0,
\qquad
\cosh^2(1/2)-\sinh^2(1/2)=1.
\]
In particular,
\[
33\sinh^2(1/2)>\frac{33}{4}>1,
\]
and therefore
\[
\cosh^2(1/2)
=
1+\sinh^2(1/2)
<
34\sinh^2(1/2).
\]
Division by the positive quantity \(\sinh^2(1/2)\) yields
\[
V^*
=
\left(\frac{\cosh(1/2)}{\sinh(1/2)}\right)^2
<34=V_{\mathrm{OA}},
\]
as required.
% lean: CausalSmith.Stat.LdpAteEfficiencySurface.VOA_balanced_value
% lean: CausalSmith.Stat.LdpAteEfficiencySurface.balanced_coth_sq_lt_34
\end{enumerate}
\end{proof}
